% Lesson 44 — Grammar: Sequence of tenses, word order — Anglais Tle SH — Cours signé OnBuch+
% Programme officiel MINESEC. Compiler avec : tectonic EN44-grammar-sequence-of-tenses-word-order.tex
\newcommand{\DOCMATIERE}{Anglais}
\newcommand{\DOCNIVEAU}{Tle SH}
\newcommand{\DOCMODULE}{Chapter 4 : Citizenship and Human Rights}
\newcommand{\DOCLECON}{EN44}
\newcommand{\DOCTITRE}{Lesson 44 — Grammar: Sequence of tenses, word order}
% OnBuch+ — préambule « cours signés » ENRICHI (compilé avec tectonic / XeLaTeX)
% Système visuel « Pop » d'OnBuch+ : fond crème (jamais blanc), bordures encre
% épaisses, ombres « dures » décalées, titres Archivo Black en capitales.
% Version MATHÉMATIQUES : graphes (pgfplots/tikz), boîtes pédagogiques
% (définition, propriété, méthode, attention, le savais-tu, exercice résolu,
% exercice, TP, bilan), page de garde et signature OnBuch+ ancrée en bas.
%
% Macros attendues dans main.tex AVANT \input{preamble} :
%   \DOCMATIERE  \DOCNIVEAU  \DOCMODULE  \DOCLECON (numéro)  \DOCTITRE
\documentclass[11pt]{article}

\usepackage{fontspec}
\usepackage[english,french]{babel} % français = langue principale ; l'anglais via \eng{..} / textebox
\usepackage[a4paper,margin=2cm,top=3.4cm,bottom=2.8cm,headheight=1.4cm,footskip=1.3cm]{geometry}
\usepackage{amsmath,amssymb,amsfonts}
\usepackage{mathtools} % \xmapsto, \coloneqq... souvent utilisés par les modèles
\usepackage{cancel} % simplifications barrées (bilans de forces)
% \abs{z} (module, valeur absolue) et \norm{u} (norme), utilisés par les modèles
\providecommand{\abs}{}\let\abs\relax\DeclarePairedDelimiter{\abs}{\lvert}{\rvert}
\providecommand{\norm}{}\let\norm\relax\DeclarePairedDelimiter{\norm}{\lVert}{\rVert}
\usepackage[table]{xcolor} % [table] : fournit aussi \rowcolors (alternance de lignes)
\usepackage{pagecolor}
\usepackage{tikz}
\usetikzlibrary{arrows.meta,positioning,calc,shapes.geometric,shapes.misc,shapes.arrows,decorations.pathmorphing,decorations.pathreplacing,decorations.markings,patterns,shadows,mindmap,trees,fit,backgrounds,matrix,intersections,angles,quotes}
\usepackage{pgfplots}
\usepackage[version=4]{mhchem} % équations nucléaires \ce{^{235}_{92}U}
\usepackage[european,siunitx]{circuitikz} % schémas électriques
\pgfplotsset{compat=1.18}
\usepgfplotslibrary{fillbetween,groupplots} % aires entre courbes (calcul intégral)
\usepackage{enumitem}
\usepackage{fancyhdr}
% [most] charge toutes les bibliothèques tcolorbox (listings, minted, raster,
% documentation...) et gonfle inutilement la mémoire TeX sur un document long
% (~300 boîtes) : on ne charge que ce qui est réellement utilisé.
\usepackage[skins,breakable]{tcolorbox}
\usepackage{graphicx}
\usepackage{colortbl}
\usepackage{booktabs}
\usepackage{tabularx}
\usepackage{diagbox} % case d'en-tête diagonale des tableaux à double entrée
% Colonnes extensibles centrée / gauche / droite (C, L, R), souvent utilisées par les modèles
\newcolumntype{C}{>{\centering\arraybackslash}X}
\newcolumntype{L}{>{\raggedright\arraybackslash}X}
\newcolumntype{R}{>{\raggedleft\arraybackslash}X}
\usepackage{array}
\usepackage{multicol}
\usepackage{multirow}
\usepackage{siunitx}
% parse-numbers=false : accepte aussi \SI{2,5 \times 10^{-3}}{...}
\sisetup{locale=FR,output-decimal-marker={,},group-minimum-digits=4,parse-numbers=false}
\DeclareSIUnit\ohm{\ensuremath{\symup{\Omega}}} % U+2126 absent de la police
\DeclareSIUnit\division{div} % réglages d'oscilloscope (ms/div, V/div)
\DeclareSIUnit\tr{tr} % tours (vitesse de rotation, ex. \SI{1500}{\tr\per\minute})
\DeclareSIUnit\molar{M} % concentration molaire (ex. \SI{3}{\molar}, \SI{1}{\micro\molar})
\DeclareSIUnit\an{an} % année (le modèle confond parfois avec \year, un compteur TeX) — ex. \SI{5}{\kilo\meter\per\an}
\DeclareSIUnit\FCFA{FCFA} % franc CFA (ex. \SI{500}{\FCFA\per\kilo\gram})
\DeclareSIUnit\degreeBrix{\textdegree Brix} % teneur en sucre d'un sirop/jus (réfractométrie)
\DeclareSIUnit\month{mois} % le modèle utilise parfois \month, un compteur TeX, comme unité de durée
\usepackage{qrcode}
% \degree : commande fréquemment utilisée par les modèles pour un angle ou une
% température en dehors de \SI{}{} (non fournie par les paquets chargés).
\providecommand{\degree}{^{\circ}}
% \qed (fin de démonstration), utilisé par les modèles sans amsthm
\providecommand{\qed}{\hfill$\square$}
% \barre : double barre des valeurs interdites dans un tableau de variations (tkz-tab)
\providecommand{\barre}{\ensuremath{\|}}
% environnement proof (amsthm non chargé) utilisé par les modèles
\ifdefined\proof\else\newenvironment{proof}[1][Démonstration]{\par\noindent\textit{#1.}\ }{\qed\par}\fi
\usepackage{lastpage}
\usepackage{hyperref}
\hypersetup{hidelinks}

% ============================================================
% Palette « Pop » OnBuch+ — mêmes hex que la classe P (plus_ui.dart)
% ============================================================
\definecolor{poporange}{HTML}{F59321}
\definecolor{popdark}{HTML}{E07A0C}
\definecolor{popcream}{HTML}{FFF6E8}
\definecolor{popink}{HTML}{1C1714}
\definecolor{poppurple}{HTML}{7A5AE0}
\definecolor{popgreen}{HTML}{2E9E6A}
\definecolor{poppink}{HTML}{E0457B}
\definecolor{popblue}{HTML}{2F7DE1}
\definecolor{popgold}{HTML}{C98A12}
\definecolor{popmuted}{HTML}{8A7F76}
\definecolor{popline}{HTML}{EADFD2}
\definecolor{popgray}{HTML}{8A7F76}
\colorlet{popgrey}{popgray}
% Alias tolérants (le modèle nomme parfois les couleurs différemment)
\colorlet{popwhite}{white}
\colorlet{popblack}{popink}
\colorlet{popred}{poppink}
\colorlet{popyellow}{popgold}
% Teintes claires (fonds de tableaux/encadrés) — le modèle nomme parfois
% les couleurs avec un suffixe L (ex. popblueL) sans les avoir définies.
\colorlet{poporangeL}{poporange!20}
\colorlet{popdarkL}{popdark!20}
\colorlet{poppurpleL}{poppurple!20}
\colorlet{popgreenL}{popgreen!20}
\colorlet{poppinkL}{poppink!20}
\colorlet{popblueL}{popblue!20}
\colorlet{popgoldL}{popgold!20}
\colorlet{popredL}{popred!20}
\colorlet{popgrayL}{popgray!20}
% Teintes claires pour fonds de tableaux / zones
\colorlet{popblueL}{popblue!10!white}
\colorlet{poporangeL}{poporange!14!white}
\colorlet{popgreenL}{popgreen!12!white}
\colorlet{poppurpleL}{poppurple!12!white}
\colorlet{poppinkL}{poppink!12!white}
\colorlet{popgoldL}{popgold!14!white}

\pagecolor{popcream}

% ============================================================
% Typographie
% ============================================================
\defaultfontfeatures{Ligatures=TeX}
\setmainfont{PlusJakartaSans-Medium.ttf}[Path=../../../fonts/,BoldFont=PlusJakartaSans-Bold.ttf]
% Maths en sans-serif (Fira Math) avec chiffres et lettres droites de Plus
% Jakarta, pour que les formules se fondent dans le texte.
\usepackage[math-style=ISO,bold-style=ISO]{unicode-math}
\setmathfont{FiraMath-Regular.otf}[Path=../../../fonts/]
\setmathfont{PlusJakartaSans-Medium.ttf}[Path=../../../fonts/,range={up/num,up/{Latin,latin},\mathup}]
\usepackage{icomma} % virgule décimale sans espace en mode math
% Caractères mathématiques Unicode absents des polices de texte (Plus Jakarta,
% Archivo Black) : rendus par leur équivalent mathématique, en texte comme en
% formule — sinon ils s'affichent en carré vide (ex. « Arithmétique dans ℤ »).
\usepackage{newunicodechar}
\newunicodechar{ℝ}{\ensuremath{\mathbb{R}}}
\newunicodechar{ℤ}{\ensuremath{\mathbb{Z}}}
\newunicodechar{ℂ}{\ensuremath{\mathbb{C}}}
\newunicodechar{ℕ}{\ensuremath{\mathbb{N}}}
\newunicodechar{ℚ}{\ensuremath{\mathbb{Q}}}
\newunicodechar{𝔻}{\ensuremath{\mathbb{D}}}
\newunicodechar{α}{\ensuremath{\alpha}}
\newunicodechar{β}{\ensuremath{\beta}}
\newunicodechar{γ}{\ensuremath{\gamma}}
\newunicodechar{δ}{\ensuremath{\delta}}
\newunicodechar{ε}{\ensuremath{\varepsilon}}
\newunicodechar{θ}{\ensuremath{\theta}}
\newunicodechar{λ}{\ensuremath{\lambda}}
\newunicodechar{μ}{\ensuremath{\mu}}
\newunicodechar{σ}{\ensuremath{\sigma}}
\newunicodechar{τ}{\ensuremath{\tau}}
\newunicodechar{φ}{\ensuremath{\varphi}}
\newunicodechar{ψ}{\ensuremath{\psi}}
\newunicodechar{ω}{\ensuremath{\omega}}
\newunicodechar{Σ}{\ensuremath{\Sigma}}
% majuscules produites par \MakeUppercase dans les titres (α-aminés -> Α)
\newunicodechar{Α}{\ensuremath{\alpha}}
\newunicodechar{Β}{\ensuremath{\beta}}
\newunicodechar{Γ}{\ensuremath{\Gamma}}
\newunicodechar{Δ}{\ensuremath{\Delta}}
\newunicodechar{Ε}{\ensuremath{\varepsilon}}
\newunicodechar{Θ}{\ensuremath{\Theta}}
\newunicodechar{Λ}{\ensuremath{\Lambda}}
\newunicodechar{Μ}{\ensuremath{\mu}}
\newunicodechar{Τ}{\ensuremath{\tau}}
\newunicodechar{Φ}{\ensuremath{\Phi}}
\newunicodechar{Ψ}{\ensuremath{\Psi}}
\newunicodechar{Ω}{\ensuremath{\Omega}}
\newunicodechar{↦}{\ensuremath{\mapsto}}
\newunicodechar{∈}{\ensuremath{\in}}
\newunicodechar{∉}{\ensuremath{\notin}}
\newunicodechar{∧}{\ensuremath{\wedge}}
\newunicodechar{⇔}{\ensuremath{\Leftrightarrow}}
\newunicodechar{⟺}{\ensuremath{\Longleftrightarrow}}
\newunicodechar{⇒}{\ensuremath{\Rightarrow}}
\newunicodechar{⟹}{\ensuremath{\Longrightarrow}}
\newunicodechar{∘}{\ensuremath{\circ}}
\newunicodechar{∪}{\ensuremath{\cup}}
\newunicodechar{∩}{\ensuremath{\cap}}
\newunicodechar{≡}{\ensuremath{\equiv}}
\newunicodechar{⊂}{\ensuremath{\subset}}
\newunicodechar{⊥}{\ensuremath{\perp}}
\newunicodechar{∃}{\ensuremath{\exists}}
\newunicodechar{∀}{\ensuremath{\forall}}
\newunicodechar{∛}{\ensuremath{\sqrt[3]{\,}}}
\newunicodechar{∜}{\ensuremath{\sqrt[4]{\,}}}
\newunicodechar{⃗}{\ensuremath{{}^{\to}}}
\newunicodechar{ⁿ}{\ifmmode^{n}\else\textsuperscript{\ensuremath{n}}\fi}
\newunicodechar{ˣ}{\ifmmode^{x}\else\textsuperscript{\ensuremath{x}}\fi}
\newunicodechar{ᵇ}{\ifmmode^{b}\else\textsuperscript{\ensuremath{b}}\fi}
\newunicodechar{ʳ}{\ifmmode^{r}\else\textsuperscript{\ensuremath{r}}\fi}
\newunicodechar{ᵘ}{\ifmmode^{u}\else\textsuperscript{\ensuremath{u}}\fi}
\newunicodechar{ʸ}{\ifmmode^{y}\else\textsuperscript{\ensuremath{y}}\fi}
\newunicodechar{ᵃ}{\ifmmode^{a}\else\textsuperscript{\ensuremath{a}}\fi}
\newunicodechar{ᵉ}{\ifmmode^{e}\else\textsuperscript{\ensuremath{e}}\fi}
\newunicodechar{ᵏ}{\ifmmode^{k}\else\textsuperscript{\ensuremath{k}}\fi}
\newunicodechar{ᵖ}{\ifmmode^{p}\else\textsuperscript{\ensuremath{p}}\fi}
\newunicodechar{ᴺ}{\ifmmode^{N}\else\textsuperscript{\ensuremath{N}}\fi}
\newunicodechar{ᶜ}{\ifmmode^{c}\else\textsuperscript{\ensuremath{c}}\fi}
\newunicodechar{ʰ}{\ifmmode^{h}\else\textsuperscript{\ensuremath{h}}\fi}
\newunicodechar{ᵠ}{\ifmmode^{q}\else\textsuperscript{\ensuremath{q}}\fi}
\newunicodechar{ₙ}{\ifmmode_{n}\else\textsubscript{\ensuremath{n}}\fi}
\newunicodechar{ₐ}{\ifmmode_{a}\else\textsubscript{\ensuremath{a}}\fi}
\newunicodechar{ᵢ}{\ifmmode_{i}\else\textsubscript{\ensuremath{i}}\fi}
\newunicodechar{ᵣ}{\ifmmode_{r}\else\textsubscript{\ensuremath{r}}\fi}
\newunicodechar{ₖ}{\ifmmode_{k}\else\textsubscript{\ensuremath{k}}\fi}
\newunicodechar{ᵦ}{\ifmmode_{\beta}\else\textsubscript{\ensuremath{\beta}}\fi}
\newunicodechar{ₓ}{\ifmmode_{x}\else\textsubscript{\ensuremath{x}}\fi}
\newfontfamily\popdisplay{ArchivoBlack-Regular.ttf}[Path=../../../fonts/]
\newfontfamily\poplogo{SpaceGrotesk-Bold.ttf}[Path=../../../fonts/]
\newfontfamily\popsemi{PlusJakartaSans-SemiBold.ttf}[Path=../../../fonts/]
\newfontfamily\popxbold{PlusJakartaSans-ExtraBold.ttf}[Path=../../../fonts/]
\newfontfamily\popxboldls{PlusJakartaSans-ExtraBold.ttf}[Path=../../../fonts/,LetterSpace=3.0]

\setlist[enumerate]{leftmargin=1.7em,itemsep=5pt,topsep=4pt}
\setlist[itemize]{leftmargin=1.7em,itemsep=4pt,topsep=4pt,label={\color{poporange}\textbullet}}
\setlength{\parindent}{0pt}
\setlength{\parskip}{5pt}
\renewcommand{\arraystretch}{1.25}

% pgfplots : style Pop par défaut
\pgfplotsset{
  every axis/.append style={axis line style={popink,line width=1pt},tick style={popink},
    grid style={popline,line width=0.5pt},label style={font=\small},tick label style={font=\footnotesize}},
  popcurve/.style={poporange,line width=1.8pt,no markers,smooth},
}
% Même style disponible dans un \draw TikZ simple (hors pgfplots)
\tikzset{popcurve/.style={poporange,line width=1.8pt}}
\tikzset{popgrid/.style={popline,very thin}}
\colorlet{popcurve}{poporange} % le nom du style est parfois utilisé comme couleur

% ============================================================
% Pill / squiggle
% ============================================================
\newcommand{\poppill}[3]{%
  \tikz[baseline=-0.55ex]\node[draw=popink,line width=1.2pt,rounded corners=2.6pt,%
    fill=#2,inner xsep=6.5pt,inner ysep=3pt]{%
    \popxboldls\fontsize{7}{7}\selectfont\color{#3}\MakeUppercase{#1}};%
}
\newcommand{\popsquiggle}[1]{%
  \tikz[baseline=-0.4ex]{
    \draw[#1,line width=2.6pt,line cap=round]
      plot [smooth] coordinates {(0,0.09) (0.34,0.01) (0.68,0.21) (1.02,0.01) (1.36,0.10)};
  }%
}
% Mot-clé surligné (marqueur orange sous le texte)
\newcommand{\cle}[1]{{\popxbold\color{popdark}#1}}

% ============================================================
% En-tête / pied
% ============================================================
\pagestyle{fancy}
\fancyhf{}
\renewcommand{\headrulewidth}{0pt}
\renewcommand{\footrulewidth}{0pt}
\lhead{\raisebox{-0.15ex}{\poplogo\fontsize{13}{13}\selectfont\color{popink}OnBuch\color{poporange}+}}
\rhead{\poppill{\DOCMATIERE\ \textbullet\ \DOCNIVEAU\ \textbullet\ Leçon \DOCLECON}{white}{popink}}
\lfoot{\popxbold\fontsize{7.6}{7.6}\selectfont\color{popmuted}\MakeUppercase{Cours signé OnBuch+}\ \textbullet\ {\popsemi\DOCTITRE}}
\rfoot{\tikz[baseline=-0.6ex]\node[rounded corners=8.5pt,fill=popink,minimum height=17pt,minimum width=17pt,inner xsep=5pt,inner sep=0]{\popxbold\fontsize{8}{8}\selectfont\color{popcream}\thepage\,/\,\pageref*{LastPage}};}

% ============================================================
% Titres
% ============================================================
\newcommand{\chaptitre}[2]{%
  \begin{tcolorbox}[enhanced,colback=poporange,colframe=popink,boxrule=2.6pt,arc=11pt,
      drop shadow=popink,left=15pt,right=15pt,top=12pt,bottom=11pt,before skip=3pt,after skip=18pt]
    \poppill{#2}{popink}{popcream}\par\vspace{9pt}
    {\raggedright\hyphenpenalty=10000\exhyphenpenalty=10000\popdisplay\fontsize{22}{24}\selectfont\color{popcream}\MakeUppercase{#1}\par}
  \end{tcolorbox}
}
% Section numérotée : pastille orange avec le numéro + titre Archivo
\newcounter{popsec}
\newcommand{\coursec}[1]{%
  \stepcounter{popsec}\setcounter{popsub}{0}%
  \par\vspace{16pt}\noindent
  \begin{minipage}{\linewidth}
  \tikz[baseline=-0.7ex]\node[draw=popink,line width=1.6pt,fill=poporange,rounded corners=4pt,minimum size=22pt,inner sep=0]{\popdisplay\fontsize{12}{12}\selectfont\color{popcream}\thepopsec};%
  \hspace{8pt}{\popdisplay\fontsize{15}{17}\selectfont\color{popink}\MakeUppercase{#1}}\par
  \vspace{4pt}\noindent{\color{poporange}\rule{36pt}{3.6pt}}
  \end{minipage}\par\nopagebreak\vspace{8pt}
}
\newcounter{popsub}
\newcommand{\courssub}[1]{%
  \stepcounter{popsub}%
  \par\vspace{9pt}\noindent{\popxbold\fontsize{11.5}{13}\selectfont\color{popink}{\color{poporange}\thepopsec.\thepopsub}\hspace{6pt}#1}\par\nopagebreak\vspace{4pt}
}

% ============================================================
% Boîtes pédagogiques — même recette : fond blanc, bordure épaisse colorée,
% ombre dure encre, badge plein. Titre optionnel [#1].
% ============================================================
\tcbset{popbase/.style={breakable,enhanced,colback=white,boxrule=2.3pt,arc=9pt,
  shadow={3pt}{-3pt}{0pt}{popink},left=14pt,right=14pt,top=11pt,bottom=11pt,before skip=11pt,after skip=15pt,
  fontupper=\popsemi\fontsize{10.6}{14.5}\selectfont\color{popink},
  fontlower=\popsemi\fontsize{10.6}{14.5}\selectfont\color{popink}}}
% Coche dessinée (le glyphe ✓ n'existe pas dans les polices du document)
\newcommand{\popcheck}[1]{\tikz[baseline=-0.5ex]\draw[#1,line width=1.6pt,line cap=round,line join=round](-0.13,0.0)--(-0.04,-0.09)--(0.13,0.1);}
\providecommand{\coche}{\popcheck{popgreen}}
\providecommand{\resultat}[1]{\par\smallskip\noindent\fcolorbox{popgreen}{popgreenL}{\parbox{\dimexpr\linewidth-2\fboxsep-2\fboxrule\relax}{\textbf{Résultat.} #1}}\par\smallskip}
\newcommand{\popboxhead}[3]{\poppill{#1}{#2}{white}\ifx\relax#3\relax\else\hspace{6pt}{\popxbold\fontsize{10.6}{12}\selectfont\color{popink}#3}\fi\par\vspace{7pt}}

\newtcolorbox{aretenir}[1][]{popbase,colframe=popgreen,before upper={\popboxhead{À retenir}{popgreen}{#1}}}
% Texte en anglais (document de lecture, dialogue, extrait) : cadre
% d'étude. Un titre optionnel (source, auteur) s'affiche dans l'en-tête.
\newtcolorbox{textebox}[1][]{popbase,colframe=popblue,before upper={\popboxhead{Text}{popblue}{#1}\begin{otherlanguage}{english}},after upper={\end{otherlanguage}}}
% Marques ✓ / ✗ (forme correcte / fautive) souvent utilisées par les modèles
\providecommand{\cross}{\textcolor{poppink}{\ensuremath{\times}}}
\providecommand{\xmark}{\cross}\providecommand{\crossmark}{\cross}\providecommand{\cmark}{\ensuremath{\checkmark}}
% Ligne de réponse / justification à compléter (inventée par les modèles)
\providecommand{\justification}[1]{\par\noindent\textit{Justification~:} \dotfill\par}
% \textipa (package tipa non chargé) : la police ne couvre pas l'API -> simple passage
\providecommand{\textipa}[1]{#1}
% Soulignement (ulem non chargé) : \uline, \ul
\providecommand{\uline}[1]{\underline{#1}}\providecommand{\ul}[1]{\underline{#1}}
% \glossaire{\item ...} inventée par les modèles : liste de glossaire
\providecommand{\glossaire}[1]{\begin{itemize}#1\end{itemize}}
% \alert (beamer) inventée par les modèles : texte mis en évidence
\providecommand{\alert}[1]{\textbf{\textcolor{poppink}{#1}}}
% variantes ASCII d'ordinaux écrites en macro
\providecommand{\iere}{\textsuperscript{re}}\providecommand{\ieme}{\textsuperscript{e}}\providecommand{\ere}{\textsuperscript{re}}\providecommand{\eme}{\textsuperscript{e}}\providecommand{\er}{\textsuperscript{er}}\providecommand{\ie}{\textsuperscript{e}}
% Ordinaux français écrits en macro par les modèles : 1\ière, 3\ième
\newcommand{\ière}{\textsuperscript{re}}\newcommand{\ième}{\textsuperscript{e}}
% \exemple (inventée par les modèles) : étiquette « Exemple : »
\providecommand{\exemple}{\textbf{Exemple~:}\ }
% \square utilisable aussi hors mode math (cases à cocher)
\AtBeginDocument{\let\squareMath\square\renewcommand{\square}{\ensuremath{\squareMath}}}
% Mot ou phrase en anglais à l'intérieur d'un paragraphe français
\newcommand{\eng}[1]{\ifmmode\text{\textit{#1}}\else\begingroup\selectlanguage{english}\itshape #1\endgroup\fi}
\let\esp\eng % alias toléré
% flèches d'intonation inventées par les modèles
\providecommand{\textsearrow}{}\renewcommand{\textsearrow}{\ensuremath{\searrow}}\providecommand{\textnearrow}{}\renewcommand{\textnearrow}{\ensuremath{\nearrow}}
% \fr{..} : mot français (inventé par les modèles)
\providecommand{\fr}[1]{\textit{#1}}
\newtcolorbox{exemplebox}[1][]{popbase,colframe=poporange,before upper={\popboxhead{Exemple}{poporange}{#1}}}
\newtcolorbox{definition}[1][]{popbase,colframe=popblue,before upper={\popboxhead{Définition}{popblue}{#1}}}
\newtcolorbox{propriete}[1][]{popbase,colframe=poppurple,before upper={\popboxhead{Propriété}{poppurple}{#1}}}
\newtcolorbox{methode}[1][]{popbase,colframe=popgold,before upper={\popboxhead{Méthode}{popgold}{#1}}}
\newtcolorbox{attention}[1][]{popbase,colframe=poppink,before upper={\popboxhead{Attention}{poppink}{#1}}}
\newtcolorbox{experience}[1][]{popbase,colframe=popblue,colback=popblueL,before upper={\popboxhead{Expérience}{popblue}{#1}}}
% « Le savais-tu ? » : carte violette pleine (contexte, culture, Cameroun)
\newtcolorbox{savaistu}[1][]{popbase,colback=poppurple,colframe=popink,
  fontupper=\popsemi\fontsize{10.4}{14.2}\selectfont\color{white},
  before upper={\poppill{Le savais-tu\,?}{white}{poppurple}\ifx\relax#1\relax\else\hspace{6pt}{\popxbold\color{white}#1}\fi\par\vspace{7pt}}}
% Exercice résolu : énoncé en haut, \tcblower puis la solution sur fond crème
\newcounter{popexo}\newcounter{popexr}
\newtcolorbox{exoresolu}[1][]{popbase,colframe=poporange,colbacklower=popcream,
  segmentation style={popink,line width=1.2pt,dashed},
  before upper={\stepcounter{popexr}\popboxhead{Exercice résolu \thepopexr}{poporange}{#1}},
  before lower={\poppill{Solution}{popink}{popcream}\par\vspace{6pt}}}
% Exercice (non résolu en place) — la correction est dans la partie Corrigés
\newtcolorbox{exercice}[1][]{popbase,colframe=popink,
  before upper={\stepcounter{popexo}\popboxhead{Exercice \thepopexo}{popink}{#1}}}
% Corrigé
\newtcolorbox{corrige}[1][]{popbase,colframe=popgreen,colback=popgreenL,
  before upper={\popboxhead{Corrigé}{popgreen}{#1}}}
% Objectifs / prérequis / bilan
\newtcolorbox{objectifs}{popbase,colframe=popink,colback=poporangeL,
  before upper={\popboxhead{Objectifs de la leçon}{popink}{}}}
\newtcolorbox{prerequis}{popbase,colframe=popink,colback=popblueL,
  before upper={\popboxhead{Prérequis}{popblue}{}}}
\newtcolorbox{bilan}{popbase,colframe=popink,colback=white,boxrule=2.8pt,
  before upper={\popboxhead{Fiche bilan}{poporange}{L'essentiel à connaître pour l'examen}}}
% Figure encadrée avec légende
\newtcolorbox{popfigure}[1][]{enhanced,breakable=false,colback=white,colframe=popline,boxrule=1.4pt,arc=8pt,
  left=10pt,right=10pt,top=10pt,bottom=8pt,before skip=10pt,after skip=12pt,halign=center,
  lower separated=false,halign lower=center,
  fontlower=\popsemi\fontsize{9}{11.5}\selectfont\color{popmuted},
  #1}
\newcommand{\legende}[1]{\par\vspace{6pt}{\popsemi\fontsize{9}{11.5}\selectfont\color{popmuted}#1}}

% ============================================================
% Signature OnBuch+
% ============================================================
\newcommand{\poplogoinline}[1]{{\poplogo\fontsize{#1}{#1}\selectfont\color{popink}OnBuch\color{poporange}+}}
\newcommand{\signatureonbuch}{%
  \begin{tcolorbox}[enhanced,colback=popink,colframe=popink,boxrule=2.2pt,arc=10pt,
      drop shadow=poporange,left=14pt,right=14pt,top=10pt,bottom=10pt,before skip=0pt,after skip=0pt]
    \tikz[baseline=-0.7ex]\node[circle,fill=popgreen,draw=popcream,line width=1.3pt,minimum size=22pt,inner sep=0]{\popcheck{white}};%
    \hspace{9pt}\begin{minipage}[c]{0.62\linewidth}
      {\popxbold\fontsize{12}{13}\selectfont\color{popcream}Signé\ \textbullet\ {\poplogo OnBuch\color{poporange}+}}\par\vspace{2pt}
      {\raggedright\popsemi\fontsize{8}{10.5}\selectfont\color{popline}\MakeUppercase{Équipe pédagogique OnBuch+}\\\MakeUppercase{Conforme au programme officiel MINESEC}\par}
    \end{minipage}\hfill
    {\poplogo\fontsize{22}{22}\selectfont\color{popcream}OnBuch\color{poporange}+}
  \end{tcolorbox}%
}

% ============================================================
% Page de garde
% #1 titre · #2 matière · #3 niveau · #4 module · #5 n° leçon · #6 durée ·
% #7 description / objectifs (texte ou itemize)
% ============================================================
\newcommand{\pagegarde}[7]{%
  \thispagestyle{empty}
  \newgeometry{a4paper,margin=1.7cm,top=1.6cm,bottom=1.4cm}%
  \begin{tikzpicture}[remember picture,overlay]
    \fill[poporange!18] ([xshift=-3.2cm,yshift=-3.6cm]current page.north east) circle (4.6cm);
    \fill[poppurple!14] ([xshift=2.4cm,yshift=7.6cm]current page.south west) circle (3.8cm);
  \end{tikzpicture}%
  {\poplogo\fontsize{30}{30}\selectfont\color{popink}OnBuch\color{poporange}+}\hfill
  \raisebox{0.6ex}{\poppill{Cours signé \textbullet\ Premium}{popink}{popcream}}\par
  \vspace{1pt}\popsquiggle{poppurple}\par
  \vspace{8pt}
  \begin{tcolorbox}[enhanced,colback=poporange,colframe=popink,boxrule=2.8pt,arc=13pt,
      drop shadow=popink,left=18pt,right=18pt,top=12pt,bottom=13pt,after skip=10pt]
    \poppill{#2 \textbullet\ #3}{popink}{popcream}\hspace{5pt}\poppill{#4}{white}{popink}\par\vspace{8pt}
    {\popdisplay\fontsize{14}{14}\selectfont\color{popink}LEÇON #5}\par\vspace{4pt}
    {\raggedright\hyphenpenalty=10000\exhyphenpenalty=10000\popdisplay\fontsize{25}{27}\selectfont\color{popcream}\MakeUppercase{#1}\par}
    \vspace{8pt}
    {\popxbold\fontsize{9.5}{9.5}\selectfont\color{popink}\MakeUppercase{Durée conseillée : #6}}
  \end{tcolorbox}
  \begin{tcolorbox}[enhanced,colback=white,colframe=popink,boxrule=2.2pt,arc=10pt,
      drop shadow=popink,left=16pt,right=16pt,top=10pt,bottom=10pt,after skip=8pt]
    \poppill{Dans cette leçon}{poppurple}{white}\par\vspace{5pt}
    \popsemi\fontsize{9.3}{12.6}\selectfont\color{popink}#7
  \end{tcolorbox}
  \noindent\begin{minipage}[t]{0.487\linewidth}
    \begin{tcolorbox}[enhanced,colback=popgreen,colframe=popink,boxrule=2.2pt,arc=10pt,
        drop shadow=popink,left=11pt,right=11pt,top=10pt,bottom=10pt,height=3.1cm,valign=center]
      \tcbox[colback=white,colframe=popink,boxrule=1.3pt,arc=5pt,boxsep=0pt,nobeforeafter,
        left=3pt,right=3pt,top=3pt,bottom=3pt,valign=center]{\qrcode[height=1.9cm]{https://play.google.com/store/apps/details?id=com.luvvix.onbuch&hl=fr}}%
      \hspace{8pt}\begin{minipage}[c]{0.5\linewidth}
        {\popdisplay\fontsize{11}{12.5}\selectfont\color{white}\MakeUppercase{Télécharge OnBuch}}\par\vspace{4pt}
        {\raggedright\popsemi\fontsize{8.4}{11}\selectfont\color{white}Gratuit \textbullet\ Google Play\\Tuteur IA, annales, concours, résultats\par}
      \end{minipage}
    \end{tcolorbox}
  \end{minipage}\hfill
  \begin{minipage}[t]{0.487\linewidth}
    \begin{tcolorbox}[enhanced,colback=poppurple,colframe=popink,boxrule=2.2pt,arc=10pt,
        drop shadow=popink,left=11pt,right=11pt,top=10pt,bottom=10pt,height=3.1cm,valign=center]
      \tikz[baseline=-0.7ex]\node[circle,fill=white,draw=popink,line width=1.3pt,minimum size=1.6cm,inner sep=0]{\popdisplay\fontsize{24}{24}\selectfont\color{poppurple}+};%
      \hspace{8pt}\begin{minipage}[c]{0.6\linewidth}
        {\popdisplay\fontsize{11}{12.5}\selectfont\color{white}\MakeUppercase{Passe à OnBuch+}}\par\vspace{4pt}
        {\raggedright\popsemi\fontsize{8.4}{11}\selectfont\color{white}Tous les cours signés, Tuteur IA illimité, fiches PDF — onglet OnBuch+ de l'app\par}
      \end{minipage}
    \end{tcolorbox}
  \end{minipage}
  \vspace{8pt}
  \popcontenu
  \vfill
  \signatureonbuch
  \restoregeometry
  \newpage
}

% Rangée « Ce cours contient » (page de garde)
\newcommand{\poptile}[3]{% #1 couleur #2 gros texte #3 légende
  \begin{tcolorbox}[enhanced,colback=#1,colframe=popink,boxrule=2pt,arc=9pt,shadow={2.5pt}{-2.5pt}{0pt}{popink},
      left=6pt,right=6pt,top=9pt,bottom=9pt,halign=center,valign=center,height=1.75cm,width=\linewidth,nobeforeafter]
    {\popdisplay\fontsize{12.5}{13}\selectfont\color{white}\MakeUppercase{#2}}\par\vspace{3pt}
    {\centering\popsemi\fontsize{7.8}{9.5}\selectfont\color{white}#3\par}
  \end{tcolorbox}}
\newcommand{\popcontenu}{%
  \par\noindent{\popxboldls\fontsize{7.5}{7.5}\selectfont\color{popmuted}CE COURS SIGNÉ CONTIENT}\par\vspace{7pt}
  \noindent\begin{minipage}[t]{0.235\linewidth}\poptile{popblue}{Cours}{complet, illustré, pas à pas}\end{minipage}\hfill
  \begin{minipage}[t]{0.235\linewidth}\poptile{poporange}{Méthodes}{et exercices résolus}\end{minipage}\hfill
  \begin{minipage}[t]{0.235\linewidth}\poptile{poppink}{Exercices}{type examen + corrigés}\end{minipage}\hfill
  \begin{minipage}[t]{0.235\linewidth}\poptile{popgreen}{Fiche bilan}{l'essentiel à retenir}\end{minipage}\par}

% Sommaire
\newtcolorbox{sommaire}{enhanced,colback=white,colframe=popink,boxrule=2.2pt,arc=10pt,
  drop shadow=popink,left=18pt,right=18pt,top=13pt,bottom=13pt,before skip=6pt,after skip=20pt,
  before upper={\poppill{Sommaire}{poppurple}{white}\par\vspace{9pt}\popsemi\fontsize{10.6}{14.5}\selectfont\color{popink}}}

% Signature de fin de cours — poussée tout en bas de la dernière page
\newcommand{\findecours}{%
  \par\vfill\nopagebreak
  \begin{center}\popsquiggle{poporange}\end{center}\vspace{4pt}
  \signatureonbuch
}
\AtBeginDocument{\def\llbracket{\mathopen{[\mkern-3.5mu[}}\def\rrbracket{\mathclose{]\mkern-3.5mu]}}} % intervalles d'entiers (glyphe ⟦ absent de la police)
% Glyphes absents de Fira Math : redéfinis à partir de glyphes disponibles
\AtBeginDocument{%
  \def\setminus{\mathbin{\backslash}}\let\smallsetminus\setminus
  \def\vdots{\mathord{\vbox{\baselineskip3.2pt\lineskiplimit0pt\kern4pt\hbox{.}\hbox{.}\hbox{.}}}}%
  \def\ddots{\mathinner{\mkern1mu\raise6.5pt\hbox{.}\mkern2mu\raise3.6pt\hbox{.}\mkern2mu\raise0.7pt\hbox{.}\mkern1mu}}%
  \def\triangle{\mathord{\scalebox{0.9}{$\symup{\Delta}$}}}%
  \def\nabla{\mathord{\raisebox{\depth}{\scalebox{1}[-1]{$\symup{\Delta}$}}}}%
  \def\checkmark{\popcheck{popdark}}%
  \def\star{\mathbin{\ast}}\def\bigstar{\mathord{\ast}}%
  \def\diamond{\mathbin{\scalebox{0.6}{$\Diamond$}}}%
  \def\bigcup{\mathop{\mathchoice{\raisebox{-0.3em}{\scalebox{1.7}{$\cup$}}}{\scalebox{1.25}{$\cup$}}{\cup}{\cup}}\displaylimits}%
  \def\bigcap{\mathop{\mathchoice{\raisebox{-0.3em}{\scalebox{1.7}{$\cap$}}}{\scalebox{1.25}{$\cap$}}{\cap}{\cap}}\displaylimits}%
  \def\lesseqqgtr{\mathrel{\lessgtr}}\def\bigoplus{\mathop{\oplus}\displaylimits}%
}
\providecommand{\ssi}{si et seulement si} % abréviation fréquente
\AtBeginDocument{\providecommand{\wideparen}{\overparen}} % arc de cercle

\begin{document}

\pagegarde{Lesson 44 — Grammar: Sequence of tenses, word order}{Anglais}{Tle SH}{Chapter 4 : Citizenship and Human Rights}{EN44}{2 h}{Cette leçon de grammaire étudie la concordance des temps (sequence of tenses) et l'ordre des mots en anglais. À travers des exemples liés à la citoyenneté et aux droits de l'homme, les élèves apprennent à rapporter des paroles et à construire des phrases correctes, en évitant les pièges typiques des francophones.
\vspace{4pt}
\begin{multicols}{2}\small
\begin{itemize}
\item À la fin de cette leçon, je sais expliquer la règle de la concordance des temps après un verbe au passé
\item je sais transformer une phrase au discours direct en discours indirect avec les bons changements de temps
\item je sais identifier les verbes qui déclenchent la concordance des temps (say, tell, think, believe, know, promise...)
\item je sais appliquer l'ordre des mots anglais : sujet – verbe – complément – lieu – temps
\item je sais placer correctement les adverbes de fréquence et les adverbes de manière dans la phrase
\item je sais repérer les cas où la concordance ne s'applique pas (vérités générales, faits encore vrais)
\end{itemize}\end{multicols}}

\begin{sommaire}
\begin{multicols}{2}
\begin{enumerate}
\item Observation et découverte
\item La concordance des temps : la règle
\item Exceptions et cas particuliers
\item L'ordre des mots en anglais
\item Comparaison français/anglais et pièges des francophones
\item Entraînement guidé et production
\item Activité d'intégration
\item Méthodes et exercices résolus
\item Exercices
\item Corrigés des exercices
\item Fiche bilan
\end{enumerate}
\end{multicols}
\end{sommaire}

\begin{objectifs}
\begin{itemize}
\item À la fin de cette leçon, je sais expliquer la règle de la concordance des temps après un verbe au passé
\item je sais transformer une phrase au discours direct en discours indirect avec les bons changements de temps
\item je sais identifier les verbes qui déclenchent la concordance des temps (say, tell, think, believe, know, promise...)
\item je sais appliquer l'ordre des mots anglais : sujet – verbe – complément – lieu – temps
\item je sais placer correctement les adverbes de fréquence et les adverbes de manière dans la phrase
\item je sais repérer les cas où la concordance ne s'applique pas (vérités générales, faits encore vrais)
\item je sais comparer les structures anglaises et françaises pour éviter les calques
\item je sais rédiger un court paragraphe rapportant un discours sur les droits des citoyens
\end{itemize}
\end{objectifs}

\begin{prerequis}
\begin{itemize}
\item Les temps principaux de l'indicatif anglais (present simple, past simple, present perfect, future)
\item Le discours direct et indirect vu en classe de Première
\item Les pronoms personnels et possessifs
\item La structure de base de la phrase affirmative anglaise (S + V + C)
\item Les verbes irréguliers courants (say–said, tell–told, know–knew...)
\end{itemize}
\end{prerequis}

\newpage

\coursec{Observation et découverte}

\courssub{Situation d'accroche : le journaliste de la CRTV}

Imagine a journalist from CRTV reporting on yesterday's Human Rights Day celebration in Yaoundé. On the radio, you hear: \eng{“The Minister said that the government had protected citizens' rights.”} Why \eng{said} and \eng{had protected} when the events are recent? And why do we say \eng{“The government protects human rights”} and not \eng{“protects human rights the government”}? Every day, when we report what someone declared about elections, strikes or court cases in Cameroon, we must shift tenses backwards and respect the English word order. This lesson gives you the keys to report speech and ideas like a professional journalist.

\textbf{Problématique :} quand on rapporte les paroles de quelqu'un après un verbe au passé comme \eng{said}, que devient le temps du verbe rapporté ? Et dans quel ordre faut-il placer les mots d'une phrase anglaise ? C'est ce que nous allons découvrir ensemble, en observant d'abord, avant d'énoncer la règle.

\begin{methode}[Comment observer comme un grammairien]
Avant toute règle, voici la démarche d'observation en quatre étapes que nous utiliserons pour chaque exemple :
\begin{enumerate}
\item \textbf{Souligner} le verbe introducteur (le verbe de parole : \eng{said}, \eng{told}, \eng{declared}) et noter son temps.
\item \textbf{Encadrer} le verbe de la subordonnée (le verbe de la parole rapportée).
\item \textbf{Comparer} les deux temps : le verbe rapporté a-t-il « reculé » dans le passé ?
\item \textbf{Repérer} les autres changements : pronoms, possessifs, expressions de temps et de lieu.
\end{enumerate}
\end{methode}

\courssub{Lecture guidée : six exemples contrastés}

Lisons d'abord ce court extrait d'un reportage, puis les six paires d'exemples qui en sont tirées.

\begin{textebox}[CRTV News — Human Rights Day in Yaoundé]
Yesterday, December 10th, Yaoundé celebrated Human Rights Day. In the morning, the Minister of Justice addressed the crowd at the 20th May Boulevard. “We protect every citizen,” he declared, “and our courts are independent.” A human rights activist from Buea told our reporter: “I am proud of my country, but we will continue the fight for justice.” Later, a student from the University of Yaoundé said: “I have never seen such a big celebration here.” In the afternoon, a police officer explained: “We arrested no one today because the march was peaceful.” Finally, a traditional ruler from the West Region announced: “My people are building a community centre next year.” Our journalist reported all these declarations on the evening news, using reported speech: the Minister said that they protected every citizen, the activist said she was proud of her country, and the student said he had never seen such a big celebration there.
\end{textebox}

\textbf{Glossaire :} \emph{to address a crowd} (v., s'adresser à une foule) ; \emph{court} (n., tribunal) ; \emph{to arrest} (v., arrêter) ; \emph{peaceful} (adj., pacifique) ; \emph{traditional ruler} (n., chef traditionnel) ; \emph{reported speech} (n., discours indirect).

\textbf{Questions de compréhension :}
\begin{enumerate}
\item Where did the celebration take place?
\item What did the Minister declare?
\item Why did the police arrest no one?
\item What three declarations did the journalist report in the evening?
\end{enumerate}

\courssub{Les six paires d'exemples : direct / indirect}

Observons maintenant les six paires. Pour chacune, appliquons la méthode : verbe introducteur souligné mentalement, verbe rapporté encadré.

\begin{exemplebox}[Six exemples contrastés]
\begin{enumerate}
\item \eng{The President said: “We protect every citizen.”} → \eng{The President said (that) they protected every citizen.} « Le Président a dit qu'ils protégeaient chaque citoyen. »
\item \eng{She said, “I am proud of my country.”} → \eng{She said she was proud of her country.} « Elle a dit qu'elle était fière de son pays. »
\item \eng{He said, “We will continue the fight.”} → \eng{He said they would continue the fight.} « Il a dit qu'ils continueraient le combat. »
\item \eng{The student said, “I have never seen such a celebration.”} → \eng{The student said he had never seen such a celebration.} « L'étudiant a dit qu'il n'avait jamais vu une telle célébration. »
\item \eng{The officer said, “We arrested no one today.”} → \eng{The officer said they had arrested no one that day.} « Le policier a dit qu'ils n'avaient arrêté personne ce jour-là. »
\item \eng{The chief said, “My people are building a centre here.”} → \eng{The chief said his people were building a centre there.} « Le chef a dit que son peuple construisait un centre là-bas. »
\end{enumerate}
\end{exemplebox}

\courssub{Repérage guidé : que devient le temps du verbe après \eng{said} ?}

Regardons les verbes encadrés. Dans les six cas, le verbe introducteur est au \cle{prétérit} (\eng{said}). Et le verbe rapporté ?

\begin{center}
\begin{tabularx}{\linewidth}{|X|X|X|}\hline
\rowcolor{popblueL} Direct speech & Reported speech & Recul observé \\ \hline
\eng{protect} (present simple) & \eng{protected} (past simple) & présent → prétérit \\ \hline
\eng{am} (present simple) & \eng{was} (past simple) & présent → prétérit \\ \hline
\eng{will continue} (future) & \eng{would continue} (conditional) & will → would \\ \hline
\eng{have never seen} (present perfect) & \eng{had never seen} (past perfect) & present perfect → past perfect \\ \hline
\eng{arrested} (past simple) & \eng{had arrested} (past perfect) & prétérit → plus-que-parfait \\ \hline
\eng{are building} (present continuous) & \eng{were building} (past continuous) & présent continu → prétérit continu \\ \hline
\end{tabularx}
\end{center}

\textbf{Constat :} dans les six cas, le verbe rapporté fait un \cle{pas en arrière} dans le temps. C'est ce que les grammairiens appellent le \emph{backshift} (le « recul du temps »). On dirait que le passé du verbe introducteur « contamine » toute la phrase.

\begin{popfigure}
\begin{tikzpicture}[
  box/.style={draw=popblue, rounded corners=3pt, fill=popblueL, align=center, font=\small, inner sep=4pt},
  rbox/.style={draw=poporange, rounded corners=3pt, fill=poporangeL, align=center, font=\small, inner sep=4pt},
  arr/.style={-{Stealth[length=3mm]}, thick, popdark}
]
\node[font=\bfseries, popblue] at (2.2,4.6) {Direct speech};
\node[font=\bfseries, poporange] at (8.2,4.6) {Reported speech};
\node[box] (a1) at (2.2,3.6) {is / am};
\node[rbox] (b1) at (8.2,3.6) {was};
\draw[arr] (a1) -- (b1);
\node[box] (a2) at (2.2,2.7) {will};
\node[rbox] (b2) at (8.2,2.7) {would};
\draw[arr] (a2) -- (b2);
\node[box] (a3) at (2.2,1.8) {has / have + p.p.};
\node[rbox] (b3) at (8.2,1.8) {had + p.p.};
\draw[arr] (a3) -- (b3);
\node[box] (a4) at (2.2,0.9) {past simple};
\node[rbox] (b4) at (8.2,0.9) {had + p.p.};
\draw[arr] (a4) -- (b4);
\node[box] (a5) at (2.2,0.0) {is building};
\node[rbox] (b5) at (8.2,0.0) {was building};
\draw[arr] (a5) -- (b5);
\node[align=center, font=\itshape\small, popmuted] at (5.2,-0.9) {After a past reporting verb (said), every tense moves one step back.};
\end{tikzpicture}
\legende{Le recul du temps (backshift) après un verbe de parole au passé}
\end{popfigure}

\courssub{Observation des pronoms, possessifs et expressions de temps}

Le temps du verbe n'est pas le seul à changer. Reprenons les exemples 2, 5 et 6 :

\begin{itemize}
\item Exemple 2 : \eng{I am proud of \textbf{my} country} → \eng{she was proud of \textbf{her} country}. Le pronom \eng{I} devient \eng{she}, le possessif \eng{my} devient \eng{her} : on adopte le point de vue du rapporteur, pas celui du locuteur original.
\item Exemple 5 : \eng{today} → \eng{that day}. Le « aujourd'hui » du policier n'est plus le « aujourd'hui » du journaliste du soir.
\item Exemple 6 : \eng{here} → \eng{there}. Le « ici » du chef n'est pas le « ici » du journaliste.
\end{itemize}

\begin{center}
\begin{tabularx}{\linewidth}{|X|X|X|}\hline
\rowcolor{popblueL} Direct speech & Reported speech & Français \\ \hline
\eng{now} & \eng{then} & maintenant → alors \\ \hline
\eng{today} & \eng{that day} & aujourd'hui → ce jour-là \\ \hline
\eng{yesterday} & \eng{the day before} & hier → la veille \\ \hline
\eng{tomorrow} & \eng{the next day} & demain → le lendemain \\ \hline
\eng{here} & \eng{there} & ici → là \\ \hline
\eng{this} & \eng{that} & ceci → cela \\ \hline
\eng{I / my} & \eng{he/she / his/her} & je/mon → il/son \\ \hline
\end{tabularx}
\end{center}

\courssub{Repérage de l'ordre des mots}

Observons maintenant la construction des phrases rapportées : \eng{The President said (that) they protected every citizen.} Qui fait l'action ? \eng{they}. Où est le verbe ? Juste après le sujet : \eng{protected}. Quoi ? \eng{every citizen}, à la fin.

L'ordre anglais est rigide : \cle{Sujet – Verbe – Complément}. Contrairement au français, on ne peut presque jamais déplacer le sujet après le verbe dans une phrase affirmative, et le sujet est \textbf{toujours exprimé}.

\begin{exemplebox}[Ordre des mots : qui ? quoi ? où ? quand ?]
\eng{The Minister (who?) announced (did what?) the new law (what?) in Yaoundé (where?) yesterday (when?).} « Le ministre a annoncé la nouvelle loi à Yaoundé hier. »

Ordre anglais : \textbf{Sujet + Verbe + Complément + Lieu + Temps}. En français, on dirait naturellement « hier, à Yaoundé, le ministre a annoncé… » : l'anglais place les compléments de lieu et de temps \emph{à la fin}, dans l'ordre lieu puis temps.
\end{exemplebox}

\begin{attention}[Piège : le recul du temps n'est pas une traduction mot à mot]
Un francophone traduit souvent mécaniquement : « Il a dit qu'il \emph{viendra} » → \eng{He said he will come}. Faux ! Le français garde parfois le futur, mais l'anglais exige le recul : \eng{He said he \textbf{would} come.} De même, ne dites jamais \eng{protects human rights the government} en calquant « protège les droits de l'homme le gouvernement » : le sujet anglais précède toujours le verbe.
\end{attention}

\courssub{À vous de formuler l'hypothèse}

\begin{experience}[Devine la règle]
En groupes de quatre, observez à nouveau les six paires d'exemples et répondez par écrit, en deux ou trois phrases :
\begin{enumerate}
\item What happens to the verb tense after \eng{said}?
\item What happens to \eng{I}, \eng{my}, \eng{today} and \eng{here}?
\item In what order do the subject, verb and complements appear?
\end{enumerate}
Chaque groupe propose ensuite sa règle à la classe. Le professeur confronte les hypothèses avant la synthèse officielle de la section suivante.
\end{experience}

\begin{exoresolu}[Premier essai de transformation]
Énoncé : transformez en discours indirect après \eng{The activist said…} la phrase suivante : \eng{“We are fighting for our rights today.”}
\tcblower
Solution détaillée :
\begin{enumerate}
\item Verbe introducteur : \eng{said} (prétérit) → il faut un recul du temps.
\item \eng{are fighting} (present continuous) → \eng{were fighting} (past continuous).
\item \eng{We} → \eng{they} ; \eng{our} → \eng{their} (point de vue du rapporteur).
\item \eng{today} → \eng{that day}.
\item Ordre des mots : sujet + verbe + compléments, sans inversion.
\end{enumerate}
Réponse : \eng{The activist said (that) they were fighting for their rights that day.} « L'activiste a dit qu'ils se battaient pour leurs droits ce jour-là. »
\end{exoresolu}

\begin{aretenir}[Ce que l'observation nous a déjà appris]
Après un verbe de parole au passé (\eng{said}, \eng{told}, \eng{declared}), le verbe rapporté recule d'un temps, les pronoms et possessifs s'adaptent au point de vue du rapporteur, les marqueurs de temps et de lieu se décalent (\eng{today → that day}, \eng{here → there}), et l'ordre des mots reste Sujet – Verbe – Complément. La section suivante transformera ces observations en règle complète, avec ses exceptions.
\end{aretenir}

\coursec{La concordance des temps : la règle}

Dans la section précédente, nous avons observé un article de presse et un dialogue dans lesquels des journalistes et des citoyens \emph{rapportaient} les paroles d'autres personnes : \eng{The mayor said that the elections had been free and fair.} Nous avons remarqué que les temps des verbes « glissaient » mystérieusement vers le passé. Le moment est venu de transformer cette observation en règle précise : c'est ce que les grammairiens anglais appellent la \cle{sequence of tenses} (la concordance des temps) ou encore le \cle{backshift} (le « recul » des temps).

\courssub{Le principe général : le recul d'un cran dans le passé}

\begin{definition}[Sequence of tenses / Backshift]
La \cle{sequence of tenses} (concordance des temps) est la règle selon laquelle, lorsque le verbe introducteur d'un discours rapporté est au \textbf{passé} (\eng{said}, \eng{told}, \eng{asked}, \eng{thought}, \eng{explained}, \eng{promised}...), le temps du verbe de la subordonnée \textbf{recule d'un cran vers le passé}. Ce phénomène s'appelle le \cle{backshift}.
\end{definition}

Observons d'abord deux exemples tirés de notre thème du chapitre, la citoyenneté et les droits de l'homme :

\eng{“I will vote on Sunday.” → He said he would vote on Sunday.}

« Je voterai dimanche. » → Il a dit qu'il voterait dimanche.

\eng{“We have won the case.” → They said they had won the case.}

« Nous avons gagné le procès. » → Ils ont dit qu'ils avaient gagné le procès.

Remarquez le mécanisme : \eng{will} devient \eng{would}, \eng{have won} devient \eng{had won}. Le verbe de la parole rapportée fait un « pas en arrière » dans le temps, comme si le fait de rapporter la parole plus tard poussait tout l'énoncé vers le passé.

\begin{propriete}[La règle de la concordance des temps]
\textbf{Règle fondamentale :} quand le verbe introducteur est au passé (\eng{said}, \eng{told}, \eng{thought}, \eng{believed}, \eng{announced}...), le temps de la subordonnée \textbf{recule d'un cran dans le passé} :
\begin{itemize}
\item \eng{am / is / are} → \eng{was / were}
\item \eng{will} → \eng{would}
\item \eng{have done} → \eng{had done}
\end{itemize}
\textbf{Condition :} la règle ne s'applique que si le verbe introducteur est au passé. S'il est au présent (\eng{says}, \eng{tells}), au \emph{present perfect} (\eng{has said}) ou au futur (\eng{will say}), il n'y a \textbf{aucun changement de temps} : \eng{“I am tired.” → She says she is tired.} (« Je suis fatiguée. » → Elle dit qu'elle est fatiguée.)
\end{propriete}

\begin{attention}[Comparaison avec le français]
Bonne nouvelle : le français possède un mécanisme très proche ! « Il dit : “Je viendrai.” » devient « Il a dit qu'il \emph{viendrait} » (futur → conditionnel, qui sert ici de « futur dans le passé »). La logique est donc familière. La différence, c'est que l'anglais applique ce recul de façon \textbf{systématique et mécanique} à tous les temps, y compris là où le français hésite : « Il a dit qu'il \emph{était} malade » correspond exactement à \eng{He said he was ill}. Ne traduisez jamais mot à mot depuis le discours direct : pensez toujours « verbe introducteur au passé = recul d'un cran ».
\end{attention}

\courssub{Le tableau des correspondances}

Voici le tableau complet des correspondances à connaître par cœur. C'est l'outil central de toute cette leçon.

\begin{center}
\begin{tabularx}{\linewidth}{|X|X|X|}
\hline
\rowcolor{popblueL} Discours direct & Discours rapporté & Exemple \\
\hline
present simple (\eng{work}) & past simple (\eng{worked}) & \eng{“I work in Douala.” → She said she worked in Douala.} \\
\hline
present continuous (\eng{is working}) & past continuous (\eng{was working}) & \eng{“I am reading the Constitution.” → He said he was reading the Constitution.} \\
\hline
present perfect (\eng{has won}) & past perfect (\eng{had won}) & \eng{“We have won the case.” → They said they had won the case.} \\
\hline
past simple (\eng{voted}) & past perfect (\eng{had voted}) & \eng{“I voted early.” → She said she had voted early.} \\
\hline
\eng{will} & \eng{would} & \eng{“I will vote.” → He said he would vote.} \\
\hline
\eng{can} & \eng{could} & \eng{“I can help you.” → She said she could help me.} \\
\hline
\eng{may} & \eng{might} & \eng{“It may rain.” → He said it might rain.} \\
\hline
\eng{must} & \eng{had to} & \eng{“You must register.” → They said we had to register.} \\
\hline
\end{tabularx}
\end{center}

\begin{aretenir}[Les trois idées à retenir]
\begin{enumerate}
\item Le past perfect (\eng{had + participe passé}) est le « fond du temps » : il ne peut plus reculer. \eng{had won} reste donc \eng{had won}.
\item \eng{would}, \eng{could}, \eng{might}, \eng{should}, \eng{ought to} ne changent pas : ils sont déjà des formes « reculées ».
\item \eng{must} devient généralement \eng{had to} (obligation passée).
\end{enumerate}
\end{aretenir}

\begin{exemplebox}[Exemples traduits en contexte camerounais]
\begin{itemize}
\item \eng{“The elections are free and fair.” → The observers said the elections were free and fair.} « Les élections sont libres et transparentes. » → Les observateurs ont dit que les élections étaient libres et transparentes.
\item \eng{“I am writing a petition.” → The student said he was writing a petition.} « Je rédige une pétition. » → L'élève a dit qu'il rédigeait une pétition.
\item \eng{“The police arrested the journalist yesterday.” → The lawyer said the police had arrested the journalist the day before.} « La police a arrêté le journaliste hier. » → L'avocat a dit que la police avait arrêté le journaliste la veille.
\item \eng{“We will defend your rights.” → The lawyers promised they would defend our rights.} « Nous défendrons vos droits. » → Les avocats ont promis qu'ils défendraient nos droits.
\item \eng{“You can vote at eighteen.” → The teacher told us we could vote at eighteen.} « On peut voter à dix-huit ans. » → Le professeur nous a dit que nous pouvions voter à dix-huit ans.
\end{itemize}
\end{exemplebox}

\begin{popfigure}
\begin{tikzpicture}[scale=1]
% ligne du temps
\draw[-{Stealth[length=3mm]}, very thick, popdark] (-0.5,0) -- (11.5,0) node[right, font=\small\itshape] {time};
\foreach \x/\lab in {2/{past perfect}, 4.5/{past simple}, 7/{present}, 9.5/{future}} {
  \fill[popblue] (\x,0) circle (2.5pt);
  \node[below, font=\small, popdark] at (\x,-0.1) {\lab};
}
% flèches de backshift
\draw[-{Stealth}, very thick, poporange] (7,0.55) to[bend right=35] node[above, font=\small\itshape, poporange] {present simple} (4.5,0.55);
\draw[-{Stealth}, very thick, poporange] (4.5,1.15) to[bend right=35] node[above, font=\small\itshape, poporange] {past simple} (2,1.15);
\draw[-{Stealth}, very thick, poporange] (9.5,0.55) to[bend right=35] node[above, font=\small\itshape, poporange] {will} (7,0.55);
\node[font=\small\bfseries, poporange] at (5.7,2.3) {BACKSHIFT : chaque temps recule d'un cran};
\node[below, font=\small\itshape, popmuted] at (7,-0.65) {said / told / thought...};
\end{tikzpicture}
\legende{Frise des temps : le verbe introducteur au passé pousse chaque temps un cran vers la gauche (vers le passé).}
\end{popfigure}

\courssub{Les changements complémentaires : pronoms, possessifs, marqueurs de temps et de lieu}

Le backshift ne touche pas seulement les verbes. Quand on rapporte une parole, tout le « point de vue » change : celui qui parle n'est plus la même personne, au même endroit, au même moment. Il faut donc adapter trois catégories de mots.

\textbf{1. Les pronoms personnels et possessifs.} Ils s'adaptent logiquement au nouveau locuteur : \eng{I} → \eng{he / she} ; \eng{we} → \eng{they} ; \eng{my} → \eng{his / her} ; \eng{our} → \eng{their}.

\eng{“I have lost my identity card,” said Marie. → Marie said (that) she had lost her identity card.}

« J'ai perdu ma carte d'identité », a dit Marie. → Marie a dit qu'elle avait perdu sa carte d'identité.

\textbf{2. Les marqueurs temporels.} Le « aujourd'hui » du locuteur d'origine n'est plus le nôtre.

\begin{center}
\begin{tabularx}{\linewidth}{|X|X|X|}
\hline
\rowcolor{popblueL} Discours direct & Discours rapporté & Français \\
\hline
\eng{today} & \eng{that day} & ce jour-là \\
\hline
\eng{yesterday} & \eng{the day before / the previous day} & la veille \\
\hline
\eng{tomorrow} & \eng{the next day / the following day} & le lendemain \\
\hline
\eng{now} & \eng{then / at that moment} & à ce moment-là \\
\hline
\eng{here} & \eng{there} & là \\
\hline
\eng{this} & \eng{that} & ce / cet \\
\hline
\eng{these} & \eng{those} & ces \\
\hline
\eng{ago} & \eng{before / earlier} & auparavant \\
\hline
\eng{last week} & \eng{the week before / the previous week} & la semaine précédente \\
\hline
\eng{next week} & \eng{the following week} & la semaine suivante \\
\hline
\end{tabularx}
\end{center}

\eng{“I will register here tomorrow.” → She said she would register there the next day.}

« Je m'inscrirai ici demain. » → Elle a dit qu'elle s'inscrirait là le lendemain.

\begin{popfigure}
\begin{tikzpicture}[mindmap, grow cyclic,
  every node/.style={concept, font=\small},
  root concept/.append style={concept color=popblue, font=\large\bfseries, text=white},
  level 1/.append style={level distance=3.6cm, sibling angle=120},
  level 2/.append style={level distance=2.6cm, sibling angle=45, font=\scriptsize},
  concept color=popblue]
\node[root concept] {Reporting}
child[concept color=poporange] { node {tenses}
  child { node[concept color=poporangeL, text=popdark] {will $\to$ would} }
  child { node[concept color=poporangeL, text=popdark] {have won $\to$ had won} }
  child { node[concept color=poporangeL, text=popdark] {work $\to$ worked} }
}
child[concept color=popgreen] { node {pronouns}
  child { node[concept color=popgreenL, text=popdark] {I $\to$ he/she} }
  child { node[concept color=popgreenL, text=popdark] {my $\to$ his/her} }
}
child[concept color=poppurple] { node {time expressions}
  child { node[concept color=poppurpleL, text=popdark] {today $\to$ that day} }
  child { node[concept color=poppurpleL, text=popdark] {tomorrow $\to$ the next day} }
  child { node[concept color=poppurpleL, text=popdark] {here $\to$ there} }
};
\end{tikzpicture}
\legende{Carte mentale : rapporter une parole, c'est transformer à la fois les temps, les pronoms et les marqueurs de temps et de lieu.}
\end{popfigure}

\courssub{That : une conjonction optionnelle}

Après \eng{say} et \eng{tell}, la conjonction \cle{that} est \textbf{facultative} en anglais courant :

\eng{He said (that) he was tired.} Il a dit qu'il était fatigué.

\eng{She told me (that) she would come.} Elle m'a dit qu'elle viendrait.

On omet surtout \eng{that} à l'oral et dans un style informel. À l'écrit formel (dissertation, article, lettre officielle), il est conseillé de le conserver : il rend la phrase plus claire et plus élégante.

\begin{attention}[That est optionnel en anglais, « que » est obligatoire en français]
Contrairement au français, où « que » ne s'omet jamais (« Il a dit \emph{qu'}il viendrait »), l'anglais peut laisser tomber \eng{that}. Mais attention au piège inverse : ne traduisez jamais \eng{that} par « cela » dans ce contexte ! \eng{He said that he was ill} = « Il a dit \textbf{qu'}il était malade », jamais « Il a dit cela il était malade ».
\end{attention}

\courssub{Say ou tell ? Une distinction capitale}

Ces deux verbes signifient tous les deux « dire », mais ils ne se construisent pas de la même manière. C'est l'une des erreurs les plus fréquentes des francophones.

\begin{propriete}[Constructions de say et tell]
\begin{itemize}
\item \cle{tell} est \textbf{toujours suivi d'un complément d'objet personne} : \eng{tell somebody (that)...} — \eng{He told \textbf{me} that he was ready.} (Il m'a dit qu'il était prêt.)
\item \cle{say} n'est \textbf{jamais suivi directement d'une personne} ; on utilise \eng{say to somebody} (rare et formel) ou rien du tout : \eng{He said that he was ready.} / \eng{He said to me that he was ready.}
\end{itemize}
\end{propriete}

\begin{center}
\begin{tabularx}{\linewidth}{|X|X|X|}
\hline
\rowcolor{popblueL} Verbe & Construction & Exemple \\
\hline
\eng{say} & \eng{say (that)...} & \eng{She said she was a citizen.} \\
\hline
\eng{say} & \eng{say something to somebody} & \eng{She said hello to the judge.} \\
\hline
\eng{tell} & \eng{tell somebody (that)...} & \eng{She told the judge she was innocent.} \\
\hline
\eng{tell} & \eng{tell somebody to do something} & \eng{He told me to sit down.} \\
\hline
\end{tabularx}
\end{center}

\begin{attention}[L'erreur n° 1 des francophones : He told that...]
\eng{tell} \textbf{exige toujours un complément d'objet personne}. On dit \eng{He told \textbf{me} that he would vote.} — jamais \emph{He told that he would vote}. Si vous ne voulez pas mentionner de destinataire, utilisez \eng{say} : \eng{He said he would vote.} Retenez le réflexe : \textbf{TELL + personne, SAY + rien}. Autres erreurs à bannir : \emph{He said me that...} (dites \eng{He told me that...} ou \eng{He said to me that...}) et \emph{She told to me...} (jamais de \eng{to} après \eng{tell}).
\end{attention}

\begin{savaistu}[Le backshift dans la presse britannique]
Dans le journalisme britannique (et de plus en plus dans la presse internationale), le backshift est parfois \textbf{omis volontairement} lorsque l'information rapportée est encore vraie au moment où l'on écrit : \eng{The Minister said the strike \textbf{is} over.} (Le ministre a dit que la grève est terminée — et elle l'est toujours.) Le recul des temps reste cependant la norme attendue à l'examen : en classe de Terminale, appliquez toujours la concordance classique, sauf si la vérité exprimée est permanente ou encore d'actualité. Nous étudierons ces exceptions en détail dans la section suivante.
\end{savaistu}

\courssub{Exercice résolu et entraînement}

\begin{exoresolu}[Mettre au discours rapporté]
Énoncé — Transformez les phrases suivantes en discours rapporté, en commençant par le verbe introducteur indiqué.

1. \eng{“I will defend human rights,” the lawyer said.} (The lawyer said...)

2. \eng{“We have signed the petition,” the students said.} (The students said...)

3. \eng{“I am registering voters today,” the officer told me.} (The officer told me...)

4. \eng{“I voted here yesterday,” my neighbour said.} (My neighbour said...)

5. \eng{“You can collect your card tomorrow,” the mayor told us.} (The mayor told us...)
\tcblower
Solution détaillée.

1. \eng{The lawyer said (that) she would defend human rights.} — \eng{will} recule en \eng{would} ; \eng{I} devient \eng{she} (l'avocate est une femme). « L'avocate a dit qu'elle défendrait les droits de l'homme. »

2. \eng{The students said (that) they had signed the petition.} — present perfect \eng{have signed} → past perfect \eng{had signed} ; \eng{we} → \eng{they}. « Les élèves ont dit qu'ils avaient signé la pétition. »

3. \eng{The officer told me (that) he was registering voters that day.} — present continuous \eng{am registering} → past continuous \eng{was registering} ; \eng{today} → \eng{that day}. Notez la construction \eng{told me} : \eng{tell} exige son complément. « L'agent m'a dit qu'il inscrivait des électeurs ce jour-là. »

4. \eng{My neighbour said (that) he had voted there the day before.} — past simple \eng{voted} → past perfect \eng{had voted} ; \eng{here} → \eng{there} ; \eng{yesterday} → \eng{the day before}. « Mon voisin a dit qu'il avait voté là la veille. »

5. \eng{The mayor told us (that) we could collect our cards the next day.} — \eng{can} → \eng{could} ; \eng{you} → \eng{we} ; \eng{your} → \eng{our} ; \eng{tomorrow} → \eng{the next day}. « Le maire nous a dit que nous pourrions retirer nos cartes le lendemain. »
\end{exoresolu}

\begin{exercice}[À vous de jouer]
Transformez au discours rapporté avec le verbe introducteur donné, puis traduisez chaque phrase en français.

1. \eng{“I know my rights,” the activist said.} (The activist said...)

2. \eng{“We are organizing a debate tomorrow,” the teacher told the class.} (The teacher told the class...)

3. \eng{“The court has released the journalist,” the reporter announced.} (The reporter announced...)

4. \eng{“I may travel to Bamenda next week,” my cousin said.} (My cousin said...)

5. \eng{“You must respect the law,” the judge told the accused.} (The judge told the accused...)
\end{exercice}

\begin{corrige}[Exercice « À vous de jouer »]
1. \eng{The activist said (that) she knew her rights.} « La militante a dit qu'elle connaissait ses droits. » (present simple → past simple ; \eng{my} → \eng{her}.)

2. \eng{The teacher told the class (that) they were organizing a debate the next day.} « Le professeur a dit à la classe qu'ils organisaient un débat le lendemain. » (present continuous → past continuous ; \eng{tomorrow} → \eng{the next day}.)

3. \eng{The reporter announced (that) the court had released the journalist.} « Le journaliste a annoncé que le tribunal avait libéré le reporter. » (present perfect → past perfect.)

4. \eng{My cousin said (that) he might travel to Bamenda the following week.} « Mon cousin a dit qu'il pourrait se rendre à Bamenda la semaine suivante. » (\eng{may} → \eng{might} ; \eng{next week} → \eng{the following week}.)

5. \eng{The judge told the accused (that) he had to respect the law.} « Le juge a dit à l'accusé qu'il devait respecter la loi. » (\eng{must} → \eng{had to}.)
\end{corrige}

\begin{attention}[Récapitulatif des pièges à éviter]
\begin{itemize}
\item Ne laissez jamais le verbe au présent après \eng{said} : \emph{He said he is tired} est fautif dans la norme scolaire ; dites \eng{He said he was tired}.
\item N'oubliez pas de transformer les pronoms : \emph{She said I was tired} ne veut pas dire la même chose que \eng{She said she was tired}.
\item N'oubliez pas les marqueurs temporels : \eng{yesterday} rapporté devient \eng{the day before}.
\item Jamais \emph{He told that...} : \eng{tell} réclame une personne (\eng{He told me that...}).
\item Jamais \emph{He said me...} : dites \eng{He told me...} ou \eng{He said to me...}.
\end{itemize}
\end{attention}

\coursec{Exceptions et cas particuliers : quand le backshift ne s'applique pas}

La règle du backshift est systématique, mais l'anglais prévoit des cas où le recul des temps est \textbf{inutile}, \textbf{incorrect} ou \textbf{optionnel}. Les identifier est indispensable pour obtenir la note maximale au Baccalauréat.

\courssub{Verbe introducteur au présent, present perfect ou futur}

Si le verbe introducteur n'est \textbf{pas} au passé, il n'y a \textbf{aucun changement de temps}. Le discours rapporté garde exactement les temps du discours direct. Seuls les pronoms et les marqueurs de temps/lieu s'adaptent si le contexte change.

\begin{propriete}[Pas de backshift avec say/says/told/will say au non-passé]
\begin{itemize}
\item \eng{Present simple} : \eng{He \textbf{says} he \textbf{is} tired.} (Il dit qu'il est fatigué.)
\item \eng{Present perfect} : \eng{She \textbf{has told} me she \textbf{works} in Yaoundé.} (Elle m'a dit qu'elle travaille à Yaoundé.)
\item \eng{Future} : \eng{They \textbf{will say} they \textbf{will come}.} (Ils diront qu'ils viendront.)
\end{itemize}
\end{propriete}

\begin{exemplebox}[Exemples]
\begin{itemize}
\item \eng{“I live in Buea.” → He says he lives in Buea.} (Pas de changement : \eng{lives}.)
\item \eng{“We have finished.” → She has just told us they have finished.} (Pas de changement : \eng{have finished}.)
\item \eng{“I will help.” → He will say he will help.} (Pas de changement : \eng{will help}.)
\end{itemize}
\end{exemplebox}

\courssub{Vérités générales, faits scientifiques, proverbes}

Lorsque la proposition rapportée exprime une vérité intemporelle, une loi naturelle, un proverbe ou un fait géographique/historique immuable, le backshift est \textbf{optionnel} (et souvent évité pour souligner la permanence de la vérité).

\begin{propriete}[Vérités intemporelles : backshift optionnel]
\begin{itemize}
\item \eng{The teacher said (that) the Earth \textbf{revolves} / \textbf{revolved} around the Sun.}
\item \eng{He said (that) water \textbf{boils} / \textbf{boiled} at 100°C.}
\item \eng{My grandmother said (that) honesty \textbf{is} / \textbf{was} the best policy.}
\end{itemize}
\textbf{Conseil d'examen :} Les deux formes sont acceptées, mais le \textbf{présent} (\eng{revolves}, \eng{boils}, \eng{is}) est stylistiquement supérieur car il marque la validité permanente de l'énoncé.
\end{propriete}

\courssub{Situation encore actuelle au moment du rapport}

Si l'action ou l'état décrit est \textbf{toujours vrai} au moment où l'on rapporte la parole, le backshift est \textbf{optionnel}. On utilise souvent le présent (ou le temps original) pour insister sur l'actualité.

\begin{exemplebox}[Situation encore vraie]
\begin{itemize}
\item \eng{“I live in Douala.” → He said he \textbf{lives} / \textbf{lived} in Douala.} (Il y habite encore.)
\item \eng{“The office is open.” → She said the office \textbf{is} / \textbf{was} open.} (Il est encore ouvert maintenant.)
\item \eng{“I am a citizen.” → He said he \textbf{is} / \textbf{was} a citizen.} (Il l'est toujours.)
\end{itemize}
\textbf{Nuance :} \eng{lives} / \eng{is} = « il y habite encore / il est encore ouvert / il l'est encore ». \eng{lived} / \eng{was} = neutre, rapporte simplement la parole passée.
\end{exemplebox}

\courssub{Verbes modaux qui ne changent pas (rappel et précision)}

Certains modaux n'ont pas de forme passée distincte ou servent déjà de « passé » d'autres modaux. Ils restent \textbf{identiques}.

\begin{center}
\begin{tabularx}{\linewidth}{|X|X|X|}
\hline
\rowcolor{popblueL} Modal \& Forme au discours rapporté \& Note \\
\hline
\eng{would} \& \eng{would} \& Déjà passé de \eng{will} \\
\hline
\eng{could} \& \eng{could} \& Déjà passé de \eng{can} (capacité) \\
\hline
\eng{might} \& \eng{might} \& Déjà passé de \eng{may} \\
\hline
\eng{should} \& \eng{should} \& Pas de changement \\
\hline
\eng{ought to} \& \eng{ought to} \& Pas de changement \\
\hline
\eng{used to} \& \eng{used to} \& Pas de changement \\
\hline
\end{tabularx}
\end{center}

\begin{attention}[Must : deux cas de figure]
\begin{itemize}
\item \textbf{Obligation} : \eng{must} → \eng{had to}. \eng{“You must vote.” → He said we \textbf{had to} vote.}
\item \textbf{Déduction logique (probabilité forte)} : \eng{must} → \eng{must} (inchangé). \eng{“He must be tired.” → She said he \textbf{must} be tired.} (Elle en déduit qu'il est fatigué.)
\end{itemize}
\end{attention}

\courssub{Discours rapporté au passé avec verbe introducteur au passé : le past perfect ne bouge pas}

Comme vu dans le tableau, le \eng{past perfect} (\eng{had done}) est le « fond de tiroir » temporel. Il ne peut pas reculer davantage.

\eng{“I had already voted.” → He said he \textbf{had already voted}.} (Inchangé.)

\coursec{L'ordre des mots en anglais (Word Order)}

La maîtrise de l'ordre des mots est la clé de la phrase anglaise correcte. Contrairement au français, l'anglais est une langue à \textbf{ordre strict} (SVO : Subject + Verb + Object). Cette section récapitule les règles essentielles pour l'expression écrite et orale, en lien avec le discours rapporté.

\courssub{L'ordre affirmatif canonique : SVOMPT}

La structure de base de la phrase déclarative anglaise suit l'acronyme \cle{SVOMPT} :

\begin{center}
\begin{tabular}{|c|c|c|c|c|c|}
\hline
\rowcolor{popblueL} \textbf{S} & \textbf{V} & \textbf{O} & \textbf{M} & \textbf{P} & \textbf{T} \\
\hline
Subject & Verb & Object & Manner & Place & Time \\
(Sujet) & (Verbe) & (Objet) & (Manière) & (Lieu) & (Temps) \\
\hline
\eng{The voters} & \eng{cast} & \eng{their ballots} & \eng{secretly} & \eng{in the hall} & \eng{yesterday.} \\
\hline
\end{tabular}
\end{center}

\begin{aretenir}[Règle d'or : Le sujet ne se sépare jamais du verbe]
En français, on peut dire : « Hier, les électeurs ont voté. » (Adverbe en tête).
En anglais, l'adverbe de temps \textbf{ne se place pas} entre le sujet et le verbe :
\eng{The voters voted yesterday.} (Correct)
\eng{Yesterday, the voters voted.} (Correct, accentuation sur le temps)
\emph{The voters yesterday voted.} (\textbf{INCORRECT} — erreur typique de francophone)
\end{aretenir}

\courssub{Position des adverbes de fréquence et de degré (Mid-position)}

Les adverbes de fréquence (\eng{always, usually, often, sometimes, never}) et certains adverbes de degré (\eng{probably, certainly, already, just, still}) occupent une place précise : la \cle{mid-position}.

\begin{propriete}[Règles de la mid-position]
\begin{enumerate}
\item \textbf{Verbe simple (be, have, modaux) :} Adverbe \textbf{après} le verbe.
  \eng{She \textbf{is} \textbf{always} on time.} / \eng{They \textbf{can} \textbf{never} agree.} / \eng{He \textbf{has} \textbf{already} voted.}
\item \textbf{Verbe à plusieurs mots (auxiliaire + participe / infinitif) :} Adverbe \textbf{après le premier auxiliaire}.
  \eng{She \textbf{has} \textbf{always} \textbf{been} honest.} / \eng{They \textbf{would} \textbf{never} \textbf{have} \textbf{accepted}.}
\item \textbf{Verbe simple (autres verbes) :} Adverbe \textbf{avant} le verbe.
  \eng{She \textbf{always} \textbf{arrives} early.} / \eng{He \textbf{probably} \textbf{knows} the law.}
\item \textbf{Négation :} L'adverbe se place \textbf{après} l'auxiliaire + \eng{not} (ou après \eng{not} pour \eng{always/never}).
  \eng{He \textbf{doesn't} \textbf{usually} \textbf{complain}.} / \eng{She \textbf{hasn't} \textbf{always} \textbf{lived} here.}
\end{enumerate}
\end{propriete}

\begin{attention}[Erreur fréquente : *He knows always the law]
Jamais : \emph{He knows always the law.} (Français : « Il connaît toujours la loi. »)
Toujours : \eng{He \textbf{always} \textbf{knows} the law.}
Jamais : \emph{She is always late.} → Correct : \eng{She \textbf{is} \textbf{always} late.} (Après \eng{be}).
\end{attention}

\courssub{L'ordre dans les questions (inversion)}

En discours direct, la question exige une inversion (Auxiliaire + Sujet + Verbe). En discours rapporté (questions indirectes), \textbf{l'ordre redevient affirmatif (Sujet + Verbe)} : il n'y a \textbf{pas d'inversion}, \textbf{pas d'auxiliaire \eng{do/does/did}}, \textbf{pas de point d'interrogation}.

\begin{center}
\begin{tabularx}{\linewidth}{|X|X|}
\hline
\rowcolor{popblueL} Question directe (inversion) & Question indirecte (ordre affirmatif) \\
\hline
\eng{“Where \textbf{do} \textbf{you} \textbf{live}?”} & \eng{He asked where \textbf{I} \textbf{lived}.} \\
\hline
\eng{“\textbf{Is} \textbf{she} \textbf{coming}?”} & \eng{He asked if \textbf{she} \textbf{was} \textbf{coming}.} \\
\hline
\eng{“\textbf{Have} \textbf{they} \textbf{finished}?”} & \eng{He wondered if \textbf{they} \textbf{had} \textbf{finished}.} \\
\hline
\eng{“What \textbf{time} \textbf{does} \textbf{it} \textbf{start}?”} & \eng{She asked what time \textbf{it} \textbf{started}.} \\
\hline
\end{tabularx}
\end{center}

\begin{methode}[Transformer une question directe en question indirecte]
\begin{enumerate}
\item Identifier le verbe introducteur (\eng{ask, wonder, want to know, inquire}) et le mettre au passé.
\item Choisir la conjonction : \eng{if} / \eng{whether} (questions fermées \emph{Yes/No}) ou le mot interrogatif (\eng{who, what, where, when, why, how}) (questions ouvertes).
\item \textbf{Remettre l'ordre sujet + verbe} (supprimer \eng{do/does/did}, remettre le verbe au temps approprié avec backshift).
\item Adapter pronoms, temps, marqueurs spatio-temporels.
\item Terminer par un point (pas de point d'interrogation).
\end{enumerate}
\end{methode}

\begin{exemplebox}[Questions indirectes en contexte]
\begin{itemize}
\item \eng{“Are you a registered voter?” → The officer asked if I \textbf{was} a registered voter.}
\item \eng{“Where is the polling station?” → She asked where the polling station \textbf{was}.}
\item \eng{“Why did you vote for him?” → They wanted to know why I \textbf{had voted} for him.}
\item \eng{“How can we help?” → We wondered how we \textbf{could} help.}
\end{itemize}
\end{exemplebox}

\courssub{L'ordre des adjectifs (OSASCOMP)}

Si plusieurs adjectifs qualifient un nom, ils suivent un ordre fixe : \cle{OSASCOMP} (\textbf{Opinion, Size, Age, Shape, Colour, Origin, Material, Purpose}).

\begin{center}
\begin{tabular}{|l|l|l|}
\hline
\rowcolor{popblueL} Catégorie & Exemples & Position \\
\hline
Opinion & \eng{beautiful, difficult, interesting} & 1 \\
\hline
Size & \eng{large, small, huge} & 2 \\
\hline
Age & \eng{old, new, young} & 3 \\
\hline
Shape & \eng{round, square, rectangular} & 4 \\
\hline
Colour & \eng{red, green, blue} & 5 \\
\hline
Origin & \eng{Cameroonian, British, African} & 6 \\
\hline
Material & \eng{wooden, metal, cotton} & 7 \\
\hline
Purpose & \eng{voting (booth), walking (stick)} & 8 \\
\hline
\end{tabular}
\end{center}

\eng{A \textbf{beautiful large old round red Cameroonian wooden voting} box.}
(Une belle grande vieille boîte de vote ronde rouge camerounaise en bois.)

\begin{attention}[Virgules et \eng{and}]
Entre adjectifs de \textbf{même catégorie} : virgule ou \eng{and} (\eng{a long, difficult process} / \eng{a long and difficult process}).
Entre adjectifs de \textbf{catégories différentes} : \textbf{pas de virgule}, \textbf{pas de \eng{and}} (\eng{a beautiful red box}, pas *\eng{a beautiful, red box}).
\end{attention}

\courssub{Ordre des mots avec \eng{enough} et \eng{too}}

\begin{propriete}[\eng{Enough} vs \eng{Too}]
\begin{itemize}
\item \eng{Adjectif / Adverbe + \textbf{enough}} : \eng{old \textbf{enough}}, \eng{quickly \textbf{enough}}.
\item \eng{\textbf{Enough} + Nom} : \eng{\textbf{enough} \textbf{time}}, \eng{\textbf{enough} \textbf{ballots}}.
\item \eng{\textbf{Too} + Adjectif / Adverbe} : \eng{\textbf{too} \textbf{old}}, \eng{\textbf{too} \textbf{slowly}}.
\item \eng{\textbf{Too} + \textbf{much/many} + Nom} : \eng{\textbf{too} \textbf{much} \textbf{corruption}}, \eng{\textbf{too} \textbf{many} \textbf{delays}}.
\end{itemize}
\end{propriete}

\eng{He was \textbf{old enough} to vote.} (Il était assez âgé pour voter.)
\eng{He was \textbf{too young} to vote.} (Il était trop jeune pour voter.)
\eng{We had \textbf{enough} \textbf{observers}.} (Nous avions assez d'observateurs.)
\eng{There was \textbf{too} \textbf{much} \textbf{violence}.} (Il y avait trop de violence.)

\coursec{Entraînement guidé et production finale}

\begin{exoresolu}[Synthèse : Backshift, Word Order, Questions indirectes]
Énoncé — Transformez en discours rapporté. Attention à l'ordre des mots dans les questions.

1. \eng{“The candidate is honest,” the journalist said.}
2. \eng{“Where do you live?” the officer asked me.}
3. \eng{“We have never seen such corruption,” the observers said.}
4. \eng{“You must not speak now,” the judge told the witness.}
5. \eng{“How many ballots were counted?” the reporter asked.}
\tcblower
Solution détaillée.

1. \eng{The journalist said (that) the candidate \textbf{was} honest.} (Backshift \eng{is}→\eng{was} ; SVO respecté.)
2. \eng{The officer asked me where I \textbf{lived}.} (Question indirecte : ordre affirmatif \eng{I lived}, backshift \eng{do live}→\eng{lived}, \eng{you}→\eng{I}.)
3. \eng{The observers said (that) they \textbf{had never seen} such corruption.} (Present perfect → past perfect ; adverbe \eng{never} en mid-position après auxiliaire \eng{had}.)
4. \eng{The judge told the witness (that) he \textbf{had not to} speak \textbf{then}.} (\eng{must not} → \eng{had not to} / \eng{did not have to} ; \eng{now} → \eng{then} ; \eng{tell} + personne.)
5. \eng{The reporter asked how many ballots \textbf{had been counted}.} (Question indirecte : ordre affirmatif, passif au past perfect \eng{had been counted} car \eng{were counted} → backshift.)
\end{exoresolu}

\begin{exercice}[Entraînement complet — Baccalauréat Blanc]
\textbf{Partie A : Discours rapporté (Backshift + Pronoms + Marqueurs)}
Transformez les phrases suivantes en commençant par le verbe introducteur indiqué.

1. \eng{“I will submit my application tomorrow,” the applicant said.} (The applicant said...)
2. \eng{“The results were announced yesterday,” the radio announced.} (The radio announced...)
3. \eng{“You can't enter without a badge,” the guard told us.} (The guard told us...)
4. \eng{“We have been waiting for two hours,” the voters complained.} (The voters complained...)
5. \eng{“Did you vote in the last election?” the researcher asked me.} (The researcher asked me...)

\textbf{Partie B : Ordre des mots (Word Order)}
Remettez les mots dans l'ordre correct pour former une phrase grammaticale.

6. \eng{always / the / president / speaks / truth / the}
7. \eng{enough / money / we / have / not / campaign / for / the}
8. \eng{asked / me / where / I / lived / the / officer}
9. \eng{beautiful / new / Cameroonian / flag / a}
10. \eng{too / difficult / the / exam / was / for / students / the}

\textbf{Partie C : Production écrite (Composition)}
\eng{“Write a letter to the Mayor of your council complaining about the lack of transparency in the recent local elections. Use reported speech to quote what witnesses told you.”} (Environ 150 mots.)
\end{exercice}

\begin{corrige}[Exercice « Entraînement complet »]
\textbf{Partie A}
1. \eng{The applicant said (that) he would submit his application the next day.}
2. \eng{The radio announced (that) the results had been announced the day before.} (Passif : \eng{were announced} → \eng{had been announced}.)
3. \eng{The guard told us (that) we couldn't enter without a badge.} (\eng{can't} → \eng{couldn't} / \eng{were not allowed to} ; \eng{you} → \eng{we}.)
4. \eng{The voters complained (that) they had been waiting for two hours.} (Present perfect continuous → past perfect continuous.)
5. \eng{The researcher asked me if I had voted in the previous election.} (Question fermée → \eng{if} ; past simple \eng{did vote} → past perfect \eng{had voted} ; \eng{last} → \eng{previous}.)

\textbf{Partie B}
6. \eng{The president \textbf{always speaks the truth}.} (Mid-position \eng{always} avant verbe simple.)
7. \eng{We \textbf{do not have enough money for the campaign}.} (\eng{not} après auxiliaire \eng{do} ; \eng{enough} avant nom \eng{money}.)
8. \eng{The officer \textbf{asked me where I lived}.} (Ordre affirmatif question indirecte.)
9. \eng{\textbf{A beautiful new Cameroonian flag}.} (OSASCOMP : Opinion \eng{beautiful}, Age \eng{new}, Origin \eng{Cameroonian}.)
10. \eng{\textbf{The exam was too difficult for the students}.} (\eng{too} + adjectif ; ordre SVO standard.)
\end{corrige}

\begin{aretenir}[Check-list finale pour l'examen]
\begin{itemize}
\item \textbf{Verbe introducteur au passé ?} → Backshift obligatoire (sauf exceptions : vérité générale, situation encore actuelle, modaux invariables).
\item \textbf{Pronoms / Possessifs} adaptés ? (\eng{I}→\eng{he/she}, \eng{my}→\eng{his/her}).
\item \textbf{Marqueurs temps/lieu} adaptés ? (\eng{now}→\eng{then}, \eng{here}→\eng{there}, \eng{tomorrow}→\eng{the next day}).
\item \textbf{Question indirecte ?} → Ordre \textbf{Sujet + Verbe} (pas d'inversion, pas de \eng{do/does/did}), conjonction \eng{if/whether} ou mot interrogatif, point final.
\item \textbf{Ordre SVOMPT} respecté dans l'affirmation ? Adverbes en \textbf{mid-position} ?
\item \textbf{\eng{Say} vs \eng{Tell}} : \eng{Tell} + personne ; \eng{Say} (+ \eng{to} personne).
\item \textbf{Adjectifs} : ordre OSASCOMP.
\end{itemize}
\end{aretenir}

La maîtrise de la concordance des temps et de l'ordre des mots vous permet désormais de rapporter fidèlement et élégamment tout témoignage, discours officiel ou information médiatique — une compétence centrale pour tout citoyen éclairé et pour réussir l'épreuve d'anglais du Baccalauréat.

\coursec{Exceptions et cas particuliers}

Dans la section précédente, nous avons établi la règle générale de la \cle{concordance des temps} (\eng{sequence of tenses}) : lorsque le verbe introducteur est au passé (\eng{said, told, asked}), les temps de la subordonnée « reculent » d'un cran vers le passé (\eng{backshift}). Mais cette règle, comme toute règle vivante, connaît des \cle{exceptions} et des \cle{cas particuliers}. C'est précisément ce que les examinateurs du baccalauréat aiment tester. Observons d'abord un court texte, puis dégageons les règles.

\begin{textebox}[A civics lesson in Yaoundé]
Last Monday, our civics teacher, Mrs Ngo Bell, gave us a memorable lesson. She told us that all human beings are born free and equal in dignity and rights. One student raised his hand and said that the Universal Declaration of Human Rights was adopted in 1948. The teacher smiled and explained that this fact would never change, so the past tense was still correct. Then she asked us if we knew our rights as citizens. Peter answered that he could name at least five fundamental freedoms. She said he ought to share them with the class. Finally, she ordered us to write a short paragraph about equality before the law. On my way home, I met my friend Amina and told her that our teacher lives in my neighbourhood. It is still true today, so I kept the present tense. That evening, my father asked me what the teacher had said, and I reported the whole lesson to him. He told me that education is the most powerful weapon we can use to change the world, and I understood that some truths never grow old.
\end{textebox}

\textbf{Glossaire :} \eng{dignity} (n., dignité) ; \eng{to raise one's hand} (v., lever la main) ; \eng{to name} (v., citer, nommer) ; \eng{neighbourhood} (n., quartier) ; \eng{to grow old} (v., vieillir).

\textbf{Questions de compréhension :}
\begin{enumerate}
\item Why did the student keep the past tense in “was adopted in 1948”?
\item Which sentence shows a general truth that does not change tense?
\item How does the narrator report the teacher's order?
\end{enumerate}

\courssub{Exception 1 : les vérités générales et scientifiques}

\begin{propriete}[Une vérité permanente reste au présent]
Lorsque la subordonnée exprime une \cle{vérité générale}, une \cle{loi scientifique}, un \cle{principe moral} ou un \cle{fait permanent}, le verbe \textbf{reste au présent} même si le verbe introducteur est au passé. Le recul des temps (\eng{backshift}) est facultatif, mais le présent est préféré et toujours correct.
\par
\eng{The teacher said that all humans are born free and equal.} « Le professeur a dit que tous les humains naissent libres et égaux. »
\end{propriete}

Cette exception est logique : une vérité éternelle ne peut pas être « passée ». Dire \eng{The teacher said that all humans were born free} suggérerait que cela n'est peut-être plus vrai aujourd'hui.

\begin{exemplebox}[Vérités générales rapportées]
\eng{“All citizens are equal before the law.” → The judge said that all citizens are equal before the law.} « Le juge a dit que tous les citoyens sont égaux devant la loi. »
\par
\eng{“Water boils at 100 degrees.” → The scientist explained that water boils at 100 degrees.} « Le scientifique a expliqué que l'eau bout à 100 degrés. »
\par
\eng{“The earth revolves around the sun.” → Galileo taught that the earth revolves around the sun.} « Galilée enseignait que la Terre tourne autour du Soleil. »
\par
\eng{“Education is a fundamental right.” → The minister declared that education is a fundamental right.} « Le ministre a déclaré que l'éducation est un droit fondamental. »
\end{exemplebox}

\begin{attention}[Le français fait la même chose — profitez-en !]
Bonne nouvelle : le français applique exactement la même logique (\og le professeur a dit que les hommes \textbf{naissent} libres et égaux \fg). Ne traduisez donc jamais mécaniquement par un imparfait en anglais : \eng{The teacher said that all humans are born free} est correct, et même préférable à \eng{were born} quand la vérité reste universelle.
\end{attention}

\courssub{Exception 2 : les faits encore vrais au moment du rapport}

\begin{propriete}[Un fait toujours vrai peut garder le présent]
Si la situation décrite est \textbf{encore vraie au moment où l'on rapporte} les paroles, le présent peut être conservé (le backshift reste possible mais n'est pas obligatoire).
\par
\eng{She said she lives in Bamenda.} « Elle a dit qu'elle habite à Bamenda. » (elle y habite encore)
\par
\eng{She said she lived in Bamenda.} « Elle a dit qu'elle habitait à Bamenda. » (backshift classique ; peut suggérer que ce n'est plus le cas)
\end{propriete}

\begin{exemplebox}[Faits encore vrais]
\eng{He told me that the elections take place next year.} « Il m'a dit que les élections ont lieu l'année prochaine. »
\par
\eng{My friend said that she works at the Douala port.} « Mon amie a dit qu'elle travaille au port de Douala. »
\par
\eng{The journalist reported that Cameroon is a bilingual country.} « Le journaliste a rapporté que le Cameroun est un pays bilingue. »
\end{exemplebox}

\courssub{Exception 3 : les verbes déjà « au maximum » du passé}

\begin{propriete}[Ce qui ne recule plus]
Certains auxiliaires modaux sont \textbf{déjà des formes de passé} et ne peuvent pas reculer davantage : \cle{could}, \cle{would}, \cle{should}, \cle{might}, \cle{ought to}. De même, le \cle{past perfect} (\eng{had + participe passé}) est déjà le temps le plus reculé : il reste inchangé dans le discours rapporté.
\end{propriete}

\begin{center}
\begin{tabularx}{\linewidth}{|X|X|X|}
\hline
\rowcolor{popblueL} Discours direct & Discours rapporté & Traduction \\
\hline
“I can vote.” & He said he could vote. & Il a dit qu'il pouvait voter. \\
\hline
“I could help.” & He said he could help. & Il a dit qu'il pouvait aider. \\
\hline
“I will protest.” & She said she would protest. & Elle a dit qu'elle protesterait. \\
\hline
“I would sign.” & She said she would sign. & Elle a dit qu'elle signerait. \\
\hline
“I may come.” & He said he might come. & Il a dit qu'il viendrait peut-être. \\
\hline
“I might join.” & He said he might join. & Il a dit qu'il se joindrait peut-être. \\
\hline
“You should vote.” & She said I should vote. & Elle a dit que je devrais voter. \\
\hline
“You ought to know your rights.” & She said I ought to know my rights. & Elle a dit que je devrais connaître mes droits. \\
\hline
“I had voted before noon.” & He said he had voted before noon. & Il a dit qu'il avait voté avant midi. \\
\hline
\end{tabularx}
\end{center}

\begin{attention}[Ne créez pas de « super-passés »]
Les formes \eng{coulded}, \eng{woulded} ou \eng{had had could} n'existent pas. Une fois que le verbe est à \eng{could}, \eng{would}, \eng{might}, \eng{should}, \eng{ought to} ou au past perfect, on s'arrête là : le recul des temps est terminé.
\end{attention}

\courssub{Cas particulier : les questions rapportées}

Rapporter une question obéit à deux règles strictes.

\begin{propriete}[Questions rapportées : forme et emploi]
\textbf{1. Pas d'inversion.} Dans une question rapportée, l'ordre redevient celui d'une phrase affirmative : \textbf{sujet + verbe}. L'auxiliaire \eng{do/does/did} disparaît.
\par
\textbf{2. Questions fermées (oui/non).} On introduit la question rapportée par \cle{if} ou \cle{whether} (« si »).
\par
\eng{“Do you respect the law?” → He asked me if I respected the law.} « Il m'a demandé si je respectais la loi. »
\par
\textbf{3. Questions ouvertes (wh-).} On conserve le mot interrogatif (\eng{what, where, when, why, who, how}), suivi de sujet + verbe.
\par
\eng{“Where do you live?” → She asked me where I lived.} « Elle m'a demandé où j'habitais. »
\end{propriete}

\begin{attention}[L'erreur n°1 des francophones]
On ne dit \textbf{jamais} \eng{He asked me did I respect the law.} ni \eng{She asked me where did I live.} L'inversion et l'auxiliaire \eng{do} n'existent que dans la question directe. Dans la question rapportée : \eng{He asked me if I respected...} / \eng{She asked me where I lived...}
\end{attention}

\begin{exemplebox}[Questions rapportées variées]
\eng{“Are you a registered voter?” → The officer asked me if I was a registered voter.} « L'agent m'a demandé si j'étais un électeur inscrit. »
\par
\eng{“Did you attend the meeting?” → My mother asked me whether I had attended the meeting.} « Ma mère m'a demandé si j'avais assisté à la réunion. »
\par
\eng{“Why are human rights important?” → The examiner asked the candidates why human rights are important.} « L'examinateur a demandé aux candidats pourquoi les droits de l'homme sont importants. » (vérité générale : présent conservé)
\par
\eng{“When will the results be published?” → The students asked when the results would be published.} « Les élèves ont demandé quand les résultats seraient publiés. »
\end{exemplebox}

\courssub{Cas particulier : les ordres et demandes rapportés}

\begin{propriete}[Ordres rapportés : tell / ask / order + objet + to + base verbale]
Pour rapporter un ordre, un conseil ou une demande, on utilise la structure : \textbf{verbe introducteur + complément d'objet + \eng{to} + base verbale}. Pour un ordre négatif : \eng{not to} + base verbale.
\par
Verbes fréquents : \eng{tell} (dire à qqn de), \eng{ask} (demander à qqn de), \eng{order} (ordonner), \eng{advise} (conseiller), \eng{warn} (avertir), \eng{invite} (inviter à), \eng{forbid} (interdire).
\end{propriete}

\begin{exemplebox}[Ordres rapportés]
\eng{“Respect the national anthem.” → The teacher told us to respect the national anthem.} « Le professeur nous a dit de respecter l'hymne national. »
\par
\eng{“Don't cheat in the exam.” → The supervisor warned the candidates not to cheat in the exam.} « Le surveillant a averti les candidats de ne pas tricher à l'examen. »
\par
\eng{“Please help me carry this box.” → She asked me to help her carry the box.} « Elle m'a demandé de l'aider à porter le carton. »
\par
\eng{“Stand up!” → The judge ordered everyone to stand up.} « Le juge a ordonné à tout le monde de se lever. »
\end{exemplebox}

\begin{attention}[Deux pièges classiques]
\textbf{1.} \eng{Say} ne prend jamais de complément de personne : on dit \eng{He told me to wait} (jamais \eng{He said me to wait}). Sans complément : \eng{He said to wait} reste lourd ; préférez \eng{tell + objet}.
\par
\textbf{2.} La négation se place \textbf{avant} \eng{to} : \eng{He told me not to be late} (jamais \eng{to not be late} dans ce contexte scolaire).
\end{attention}

\courssub{Organigramme de décision et tableau récapitulatif}

\begin{popfigure}
\begin{tikzpicture}[
  node distance=0.55cm and 1.1cm,
  decision/.style={diamond, aspect=2.2, draw=poporange, fill=poporangeL, align=center, inner sep=1.5pt, font=\small},
  action/.style={rectangle, rounded corners, draw=popblue, fill=popblueL, align=center, font=\small, inner sep=4pt},
  startstop/.style={rectangle, rounded corners, draw=popgreen, fill=popgreenL, align=center, font=\small, inner sep=4pt},
  fleche/.style={-{Stealth[length=2.5mm]}, thick, popdark}
]
\node[startstop] (start) {Phrase à rapporter};
\node[decision, below=of start, align=center] (q1){Verbe introducteur\\au passé ?};
\node[action, right=of q1, align=center] (noback){Pas de backshift\\(temps inchangés)};
\node[decision, below=of q1, align=center] (q2){Vérité générale ou\\fait encore vrai ?};
\node[action, right=of q2, align=center] (keep){Présent conservé\\(backshift facultatif)};
\node[decision, below=of q2, align=center] (q3){could / would / might /\\should / ought to /\\past perfect ?};
\node[action, right=of q3, align=center] (stop){Temps inchangé\\(déjà au maximum)};
\node[action, below=of q3, align=center] (back){Backshift : le temps\\recule d'un cran};
\draw[fleche] (start) -- (q1);
\draw[fleche] (q1) -- node[above, font=\small]{non} (noback);
\draw[fleche] (q1) -- node[right, font=\small]{oui} (q2);
\draw[fleche] (q2) -- node[above, font=\small]{oui} (keep);
\draw[fleche] (q2) -- node[right, font=\small]{non} (q3);
\draw[fleche] (q3) -- node[above, font=\small]{oui} (stop);
\draw[fleche] (q3) -- node[right, font=\small]{non} (back);
\end{tikzpicture}
\legende{Organigramme de décision : faut-il appliquer le backshift ?}
\end{popfigure}

\begin{center}
\begin{tabularx}{\linewidth}{|X|X|X|}
\hline
\rowcolor{popblueL} Cas & Règle & Exemple \\
\hline
Vérité générale / scientifique & présent conservé & The judge said that all citizens are equal before the law. \\
\hline
Fait encore vrai & présent possible & She said she lives in Bamenda. \\
\hline
could, would, might, should, ought to & inchangés & He said he could vote. \\
\hline
Past perfect & inchangé & She said she had voted. \\
\hline
Question fermée & if / whether + sujet + verbe & He asked me if I respected the law. \\
\hline
Question ouverte & wh- + sujet + verbe & She asked where I lived. \\
\hline
Ordre / demande & tell/ask/order + objet + to + BV & The teacher told us to write a paragraph. \\
\hline
Ordre négatif & not to + BV & He warned me not to cheat. \\
\hline
\end{tabularx}
\end{center}

\courssub{Entraînement guidé}

\begin{exoresolu}[Mettez les phrases au discours rapporté]
1. “Human rights are universal.” (The professor explained that...) \par
2. “I live in Yaoundé.” (My cousin told me that... — elle y habite toujours) \par
3. “I can speak two official languages.” (The candidate said that...) \par
4. “Do you know your duties as a citizen?” (The policeman asked me...) \par
5. “Where did you register to vote?” (She asked me...) \par
6. “Don't forget your identity card.” (My father told me...)
\tcblower
\textbf{Solution détaillée :}
\begin{enumerate}
\item \eng{The professor explained that human rights are universal.} — Vérité générale : le présent \eng{are} est conservé (exception 1).
\item \eng{My cousin told me that she lives in Yaoundé.} — Fait encore vrai : présent conservé possible (exception 2) ; \eng{lived} serait aussi accepté.
\item \eng{The candidate said that he could speak two official languages.} — \eng{can} recule à \eng{could} ; c'est le maximum, on s'arrête là (exception 3).
\item \eng{The policeman asked me if I knew my duties as a citizen.} — Question fermée : \eng{if} + sujet + verbe, sans inversion ni \eng{did} ; \eng{know} recule à \eng{knew}.
\item \eng{She asked me where I had registered to vote.} — Question ouverte : mot interrogatif conservé, ordre sujet + verbe, backshift du prétérit au past perfect.
\item \eng{My father told me not to forget my identity card.} — Ordre négatif : \eng{tell} + objet + \eng{not to} + base verbale.
\end{enumerate}
\end{exoresolu}

\begin{exercice}[À vous de jouer]
Rapportez les paroles suivantes en commençant par le verbe introducteur indiqué :
\begin{enumerate}
\item “All citizens must obey the law.” (The mayor said that...)
\item “I am still a member of the human rights club.” (Nadia said that...)
\item “Have you ever visited the National Assembly?” (The journalist asked the students...)
\item “What should young people do for their country?” (The teacher asked us...)
\item “Speak louder, please.” (The examiner told the candidate...)
\end{enumerate}
\end{exercice}

\begin{corrige}[Exercice : À vous de jouer]
\begin{enumerate}
\item \eng{The mayor said that all citizens must obey the law.} — Vérité générale / principe permanent : \eng{must} peut être conservé (ou reculer à \eng{had to}, accepté).
\item \eng{Nadia said that she is still a member of the human rights club.} — Fait encore vrai : présent conservé possible ; \eng{was} avec backshift est aussi correct.
\item \eng{The journalist asked the students if they had ever visited the National Assembly.} — Question fermée : \eng{if}, pas d'inversion, backshift au past perfect.
\item \eng{The teacher asked us what young people should do for their country.} — Question ouverte : ordre sujet + verbe ; \eng{should} ne recule plus.
\item \eng{The examiner told the candidate to speak louder.} — Ordre rapporté : \eng{tell} + objet + \eng{to} + base verbale.
\end{enumerate}
\end{corrige}

\begin{aretenir}[L'essentiel de la section]
\begin{itemize}
\item Une vérité générale ou scientifique reste au \textbf{présent} : \eng{The teacher said that all humans are born free.}
\item Un fait encore vrai peut garder le \textbf{présent} : \eng{She said she lives in Bamenda.}
\item \eng{could, would, might, should, ought to} et le \textbf{past perfect} ne reculent plus.
\item Question rapportée : \textbf{pas d'inversion}, \eng{if/whether} pour les questions fermées, mot interrogatif conservé pour les questions ouvertes.
\item Ordre rapporté : \eng{tell/ask/order} + objet + \eng{(not) to} + base verbale.
\end{itemize}
\end{aretenir}

\coursec{L'ordre des mots en anglais}

Nous venons de voir que la concordance des temps comporte des exceptions : une vérité générale reste au présent même après un verbe au passé (\eng{The teacher said that the earth is round.}). Autrement dit, en anglais, la \emph{forme} des verbes obéit à des règles strictes. Il en va exactement de même pour la \emph{place} des mots dans la phrase. Là où le français tolère une certaine souplesse (\eng{Hier, je suis allé au marché} ou \eng{Je suis allé au marché hier}), l'anglais impose un ordre presque immuable. C'est cette architecture rigide — et rassurante une fois qu'on la maîtrise — que nous allons maintenant étudier.

\courssub{Observation : un texte pour découvrir la structure}

\begin{textebox}[A peaceful election in Douala]
Last Saturday, the citizens of Bonabéri elected their new mayor in Douala. The crowd gathered peacefully in front of the town hall early in the morning. Voters always respect the rules during elections, and they patiently waited for their turn. The mayor gave a short speech at the stadium in the afternoon. She promised that she would always defend the rights of the people. Journalists reported the event on television the same evening. Nobody ever forgets such an important day.
\end{textebox}

\textbf{Glossaire :} \emph{mayor} (nom, maire) ; \emph{town hall} (nom, mairie/hôtel de ville) ; \emph{voters} (nom, électeurs) ; \emph{speech} (nom, discours) ; \emph{to report} (verbe, rapporter, couvrir un événement).

\textbf{Questions d'observation :}
\begin{enumerate}
\item Dans la phrase \eng{The citizens of Bonabéri elected their new mayor in Douala}, quel mot exprime le lieu ? Quel mot exprime le temps ? Dans quel ordre apparaissent-ils ?
\item Où est placé l'adverbe \eng{peacefully} par rapport au verbe \eng{gathered} ?
\item Où est placé l'adverbe \eng{always} dans \eng{Voters always respect the rules} ? Et dans \eng{she would always defend} ?
\end{enumerate}

\textbf{Réponses :} le lieu (\eng{in Douala}) précède le temps (\eng{last Saturday} est placé en tête, mais à l'intérieur de la phrase le lieu vient avant le moment) ; \eng{peacefully} suit le verbe ; \eng{always} se place avant le verbe principal mais après l'auxiliaire \eng{would}. Toute la mécanique de la phrase anglaise tient dans ces trois observations.

\courssub{La structure fondamentale : Sujet + Verbe + Complément + Lieu + Temps}

\begin{propriete}[La règle d'or de la phrase affirmative anglaise]
Une phrase affirmative anglaise se construit selon l'ordre rigide :
\begin{center}
\textbf{Sujet + Verbe + Complément + Lieu + Temps} (SVCLT)
\end{center}
Le sujet précède toujours le verbe ; le complément d'objet suit immédiatement le verbe ; le complément de lieu précède le complément de temps. Contrairement au français, on ne peut pas déplacer librement ces éléments.
\end{propriete}

\begin{popfigure}
\begin{tikzpicture}[node distance=0.35cm]
\node[draw=popblue, fill=popblueL, rounded corners, minimum height=1.1cm, align=center] (s) {\textbf{Sujet}\\ The citizens};
\node[draw=poporange, fill=poporangeL, rounded corners, minimum height=1.1cm, align=center, right=of s] (v) {\textbf{Verbe}\\ elected};
\node[draw=popgreen, fill=popgreenL, rounded corners, minimum height=1.1cm, align=center, right=of v] (c) {\textbf{Complément}\\ their mayor};
\node[draw=poppurple, fill=poppurpleL, rounded corners, minimum height=1.1cm, align=center, right=of c] (l) {\textbf{Lieu}\\ in Douala};
\node[draw=poppink, fill=poppinkL, rounded corners, minimum height=1.1cm, align=center, right=of l] (t) {\textbf{Temps}\\ in 2020.};
\draw[-{Stealth}, thick, popdark] (s) -- (v);
\draw[-{Stealth}, thick, popdark] (v) -- (c);
\draw[-{Stealth}, thick, popdark] (c) -- (l);
\draw[-{Stealth}, thick, popdark] (l) -- (t);
\end{tikzpicture}
\legende{La phrase anglaise est un alignement de blocs dans un ordre fixe : \eng{The citizens elected their mayor in Douala in 2020.} (Les citoyens ont élu leur maire à Douala en 2020.)}
\end{popfigure}

\begin{exemplebox}[Le lieu avant le temps : comparaison directe]
\eng{The meeting took place in Yaoundé last Friday.} « La réunion a eu lieu à Yaoundé vendredi dernier. »

\eng{We celebrated Human Rights Day at our school on 10 December.} « Nous avons célébré la Journée des droits de l'homme dans notre école le 10 décembre. »

\eng{She arrived at the airport at six o'clock.} « Elle est arrivée à l'aéroport à six heures. »

Remarquez qu'en français, on dirait naturellement « vendredi dernier à Yaoundé » (temps avant lieu) : l'anglais fait l'inverse. Le complément de temps peut toutefois être placé \emph{en tête de phrase}, suivi d'une virgule, pour insister : \eng{Last Friday, the meeting took place in Yaoundé.}
\end{exemplebox}

\begin{propriete}[On ne sépare jamais le verbe de son complément d'objet]
En anglais, rien ne doit s'intercaler entre le verbe et son complément d'objet direct. Les adverbes et compléments circonstanciels se placent \emph{après} le complément, jamais entre le verbe et lui.
\end{propriete}

\begin{exemplebox}[Correct et incorrect]
\eng{She speaks English fluently.} « Elle parle couramment l'anglais. » — jamais : \emph{She speaks fluently English.}

\eng{He signed the document yesterday.} « Il a signé le document hier. » — jamais : \emph{He signed yesterday the document.}

\eng{The judge explained the law clearly.} « Le juge a expliqué la loi clairement. » — jamais : \emph{The judge explained clearly the law.}
\end{exemplebox}

\courssub{La place des adverbes}

\begin{propriete}[Adverbes de fréquence : avant le verbe principal, après \eng{be} ou l'auxiliaire]
Les adverbes de fréquence (\eng{always, usually, often, sometimes, rarely, never}) se placent :
\begin{itemize}
\item \textbf{avant le verbe principal} : \eng{Citizens must always respect the law.} (ici après l'auxiliaire modal \eng{must}, avant \eng{respect}) ; \eng{I often go to school.}
\item \textbf{après le verbe \eng{be}} : \eng{He is never late.} « Il n'est jamais en retard. »
\item \textbf{après le premier auxiliaire} : \eng{She has always defended human rights.} « Elle a toujours défendu les droits de l'homme. »
\end{itemize}
\end{propriete}

\begin{propriete}[Adverbes de manière : après le verbe ou après le complément]
Les adverbes de manière (\eng{slowly, peacefully, clearly, well}) se placent après le verbe s'il n'y a pas de complément, ou après le complément d'objet :
\eng{The crowd demonstrated peacefully.} « La foule a manifesté pacifiquement. »
\eng{The president spoke clearly.} « Le président a parlé clairement. »
\eng{She answered the question calmly.} « Elle a répondu à la question calmement. »
\end{propriete}

\begin{attention}[L'erreur n° 1 des francophones : l'adverbe mal placé]
En français on dit « je vais souvent à l'école » : l'adverbe suit le verbe. En anglais, c'est l'inverse pour les adverbes de fréquence.

Ne dites pas : \emph{I go often to school.} Dites : \eng{I often go to school.} « Je vais souvent à l'école. »

Ne dites pas : \emph{She speaks fluently English.} Dites : \eng{She speaks English fluently.}

Ne dites pas : \emph{They are always happy?} — attention, ici c'est juste, car \eng{always} suit \eng{be} ! Mais : \emph{He always is late} est faux ; dites \eng{He is always late.}
\end{attention}

\begin{center}
\begin{tabularx}{\linewidth}{|X|X|X|}
\hline
\rowcolor{popblueL} Français & English & Remarque \\
\hline
Je vais souvent à l'école. & I often go to school. & fréquence avant le verbe \\
\hline
Il n'est jamais en retard. & He is never late. & fréquence après \eng{be} \\
\hline
Elle a toujours défendu les droits. & She has always defended rights. & après le 1\textsuperscript{er} auxiliaire \\
\hline
La foule a manifesté pacifiquement. & The crowd demonstrated peacefully. & manière après le verbe \\
\hline
Elle parle couramment l'anglais. & She speaks English fluently. & manière après le complément \\
\hline
La réunion a eu lieu vendredi à Douala. & The meeting took place in Douala on Friday. & lieu avant temps \\
\hline
\end{tabularx}
\end{center}

\courssub{L'ordre des mots dans la question : l'inversion}

\begin{propriete}[Question : auxiliaire + sujet + verbe]
Dans une question, l'anglais opère une \cle{inversion} : l'auxiliaire (\eng{do, does, did, have, will, can, must}...) passe \emph{avant} le sujet, et le verbe principal reste après le sujet :
\begin{center}
\textbf{(Mot interrogatif) + Auxiliaire + Sujet + Verbe + Complément ?}
\end{center}
\eng{Did the citizens elect their mayor?} « Les citoyens ont-ils élu leur maire ? »
\eng{Where did the meeting take place?} « Où la réunion a-t-elle eu lieu ? »
\eng{Must citizens always respect the law?} « Les citoyens doivent-ils toujours respecter la loi ? »
Exception : quand le mot interrogatif est le \emph{sujet} de la phrase, il n'y a pas d'inversion ni d'auxiliaire : \eng{Who signed the document?} « Qui a signé le document ? »
\end{propriete}

\begin{methode}[Construire une phrase longue en quatre étapes]
\begin{enumerate}
\item Écrivez d'abord le noyau \textbf{Sujet + Verbe + Complément} : \eng{The students / wrote / a letter.}
\item Ajoutez les adverbes à leur place : fréquence avant le verbe, manière après le complément : \eng{The students carefully wrote a letter.}
\item Ajoutez le \textbf{lieu} : \eng{The students carefully wrote a letter in the classroom.}
\item Ajoutez le \textbf{temps} à la fin (ou en tête avec une virgule) : \eng{The students carefully wrote a letter in the classroom yesterday morning.}
\end{enumerate}
Vérifiez enfin que rien ne sépare le verbe de son complément d'objet.
\end{methode}

\begin{exoresolu}[Remettre la phrase dans l'ordre]
Énoncé : remettez les mots suivants dans l'ordre correct pour former une phrase affirmative : \eng{the minister / a speech / gave / at the stadium / on Saturday / always / carefully}.
\tcblower
Solution détaillée :
\begin{enumerate}
\item Noyau S + V + C : \eng{The minister gave a speech.}
\item Adverbe de manière \eng{carefully} : après le complément. Adverbe de fréquence \eng{always} : avant le verbe principal : \eng{The minister always gave a speech carefully.}
\item Lieu : \eng{at the stadium}.
\item Temps à la fin : \eng{on Saturday}.
\end{enumerate}
Phrase finale : \eng{The minister always gave a speech carefully at the stadium on Saturday.} « Le ministre a toujours prononcé un discours avec soin au stade samedi. »
Vérification : rien ne sépare \eng{gave} de \eng{a speech} ; le lieu précède le temps ; \eng{always} est avant le verbe. La phrase est correcte.
\end{exoresolu}

\begin{exercice}[À vous de jouer]
\begin{enumerate}
\item Remettez dans l'ordre : \eng{the pupils / their rights / know / at school / well}.
\item Corrigez la phrase : \emph{I go often to the market on Saturdays.}
\item Traduisez : « Les journalistes ont interviewé le maire à Bamenda hier. »
\item Formez une question à partir de : \eng{The court protected the victim.}
\end{enumerate}
\end{exercice}

\begin{corrige}[Exercice « À vous de jouer »]
\begin{enumerate}
\item \eng{The pupils know their rights well at school.} (manière \eng{well} après le complément, lieu à la fin.)
\item \eng{I often go to the market on Saturdays.} (fréquence avant le verbe.)
\item \eng{Journalists interviewed the mayor in Bamenda yesterday.} (lieu avant temps.)
\item \eng{Did the court protect the victim?} (inversion avec l'auxiliaire \eng{did} ; le verbe revient à la base verbale.)
\end{enumerate}
\end{corrige}

Retenez l'essentiel : la phrase anglaise est une construction en blocs ordonnés. Sujet, verbe, complément, lieu, temps ; les adverbes de fréquence se glissent avant le verbe, les adverbes de manière se posent après le complément, et jamais rien ne s'intercale entre le verbe et son objet. Cette discipline de l'ordre des mots vous servira immédiatement dans la section suivante, où nous comparerons systématiquement les structures françaises et anglaises pour débusquer les pièges de traduction.

\coursec{Comparaison français/anglais et pièges des francophones}

Après avoir étudié l'ordre des mots dans la phrase simple et dans la phrase interrogative, nous devons maintenant confronter nos habitudes francophones aux règles strictes de l'anglais. Le passage du discours direct au discours rapporté (indirect) est le terrain privilégié où les interférences de la langue maternelle produisent les erreurs les plus coûteuses aux examens. Observons d'abord un court échange authentique pour ancrer le contexte.

\begin{textebox}[Dialogue : Engagement citoyen à Buea]
\textbf{Context:} \textit{At a youth forum in Buea, a student leader addresses the crowd. A journalist reports his words later.}

\textbf{Direct Speech (Monday):}
\textbf{Leader:} “We \textbf{must} protect our environment today. The government \textbf{has promised} to plant trees here. I \textbf{will} organise a clean-up campaign next week. Where \textbf{are} the volunteers?”

\textbf{Reported Speech (Tuesday, newspaper article):}
\textbf{The leader urged the youth that they \textbf{had to} protect their environment \textbf{that day}. He added that the government \textbf{had promised} to plant trees \textbf{there}. He announced that he \textbf{would} organise a clean-up campaign \textbf{the following week}. He asked where the volunteers \textbf{were}.}
\end{textebox}

\begin{definition}[Glossaire]
\textbf{to urge} (v.) : exhorter, presser ; \textbf{clean-up campaign} (n.) : campagne de nettoyage ; \textbf{volunteer} (n.) : bénévole ; \textbf{following} (adj.) : suivant, d'après (placé après le nom pour le temps).
\end{definition}

\courssub{Le piège du « backshift » : le conditionnel français n'est pas toujours \eng{would}}

En français, le recul des temps (concordance) est systématique après un verbe introducteur au passé : \eng{Il dit qu'il viendra} $\rightarrow$ \eng{Il a dit qu'il viendrait} (futur $\rightarrow$ conditionnel). En anglais, la règle mécanique \eng{will $\rightarrow$ would} s'applique, mais \textbf{le conditionnel français a d'autres emplois qui ne se traduisent jamais par \eng{would}}.

\begin{propriete}[Règle d'or : \eng{If} + prétérit $\rightarrow$ \eng{would} + infinitif (2\textsuperscript{e} conditionnel)]
Dans une phrase hypothétique introduite par \eng{if} (si), le verbe de la proposition principale prend \eng{would} (ou \eng{could/might}). Le verbe après \eng{if} est au \textbf{prétérit} (past simple), \textbf{jamais} au conditionnel (\eng{would}).
\begin{center}
\begin{tabularx}{\linewidth}{|X|X|X|}
\hline
\rowcolor{popblueL} \textbf{Français (Conditionnel)} & \textbf{English (Correct)} & \textbf{Erreur fréquente (Interdite)} \\
\hline
Si j'avais le temps, \textbf{je voterais}. & \eng{If I \textbf{had} time, I \textbf{would vote}.} & \eng{If I \textbf{would have} time, I would vote.} \\
\hline
S'il savait, \textbf{il viendrait}. & \eng{If he \textbf{knew}, he \textbf{would come}.} & \eng{If he \textbf{would know}, he would come.} \\
\hline
Nous \textbf{ferions} cela si c'était possible. & \eng{We \textbf{would do} that if it \textbf{was} possible.} & \eng{We would do that if it \textbf{would be} possible.} \\
\hline
\end{tabularx}
\end{center}
\end{propriete}

\begin{attention}[Piège majeur : \eng{Would} après \eng{If} = Zéro point au Bac]
\eng{Would} (ou \eng{should/could/might}) \textbf{ne s'écrit jamais} dans la proposition introduite par \eng{if}, \eng{unless}, \eng{supposing}, \eng{in case}. C'est la signature immédiate d'un calque du français.
\begin{itemize}
    \item \eng{\textcolor{poppink}{If I would be elected, I would fight corruption.}} $\rightarrow$ \eng{If I \textbf{was/were} elected, I would fight corruption.}
    \item \eng{\textcolor{poppink}{Unless you would help me, I cannot finish.}} $\rightarrow$ \eng{Unless you \textbf{help} me, I cannot finish.}
\end{itemize}
\end{attention}

\courssub{Les 5 pièges classiques du discours rapporté (Transformation Direct $\rightarrow$ Indirect)}

Le tableau ci-dessous synthétise les cinq erreurs structurelles que commettent 90\% des apprenants francophones. La flèche verte indique la correction ; la forme en rouge est l'erreur calquée sur le français.

\begin{popfigure}
\begin{tikzpicture}[
    node distance=0.3cm and 0.5cm,
    box/.style={rectangle, draw=popline, rounded corners=3pt, minimum height=1.1cm, text width=7.5cm, align=left, font=\small, inner sep=4pt},
    good/.style={box, fill=popgreenL, draw=popgreen},
    bad/.style={box, fill=poppinkL, draw=poppink},
    title/.style={font=\bfseries\large, text width=7.5cm, align=center, fill=popblue, text=white, rounded corners=3pt, minimum height=0.8cm},
    arrow/.style={->, thick, popgreen, shorten <=2pt, shorten >=2pt}
]
% Column Titles
\node[title] (t1) at (0,0) {FRANÇAIS (Source de l'erreur)};
\node[title] (t2) at (16,0) {ENGLISH (Correct vs Faute)};

% Piège 1
\node[good, below=of t1.south west, anchor=north west] (f1) {\textbf{Piège 1 : Oubli du backshift}\par\medskip « Il a dit qu'il \textbf{viendra} » (futur conservé à l'oral familier)};
\node[bad, right=of f1, anchor=north west] (e1) {\eng{\textcolor{poppink}{He said he \textbf{will} come.}} \par\medskip \textcolor{popgreen}{\eng{He said he \textbf{would} come.}}};
\draw[arrow] (f1.east) -- (e1.west) node[midway, above, font=\tiny, text=popdark] {Backshift obligatoire};

% Piège 2
\node[good, below=of f1.south west, anchor=north west] (f2) {\textbf{Piège 2 : Place de l'adverbe}\par\medskip « Il parle \textbf{bien} l'anglais » (Adv après verbe)};
\node[bad, right=of f2, anchor=north west] (e2) {\eng{\textcolor{poppink}{He speaks \textbf{well} English.}} \par\medskip \textcolor{popgreen}{\eng{He speaks English \textbf{well}.}}};
\draw[arrow] (f2.east) -- (e2.west) node[midway, above, font=\tiny, text=popdark] {Adv après l'objet};

% Piège 3
\node[good, below=of f2.south west, anchor=north west] (f3) {\textbf{Piège 3 : Inversion en question indirecte}\par\medskip « Il a demandé où \textbf{étais}-je ? » (Inversion sujet/verbe)};
\node[bad, right=of f3, anchor=north west] (e3) {\eng{\textcolor{poppink}{He asked where \textbf{was I}.}} \par\medskip \textcolor{popgreen}{\eng{He asked where \textbf{I was}.}}};
\draw[arrow] (f3.east) -- (e3.west) node[midway, above, font=\tiny, text=popdark] {Ordre affirmatif};

% Piège 4
\node[good, below=of f3.south west, anchor=north west] (f4) {\textbf{Piège 4 : Objet manquant après \eng{tell}}\par\medskip « Il a dit que... » (Pas d'objet indirect)};
\node[bad, right=of f4, anchor=north west] (e4) {\eng{\textcolor{poppink}{He \textbf{told that}...}} \par\medskip \textcolor{popgreen}{\eng{He \textbf{told me} that...}} \par\smallskip \textit{ou} \eng{He \textbf{said} that...}};
\draw[arrow] (f4.east) -- (e4.west) node[midway, above, font=\tiny, text=popdark] {\eng{Tell} + objet};

% Piège 5
\node[good, below=of f4.south west, anchor=north west] (f5) {\textbf{Piège 5 : Marqueurs spatio-temporels inchangés}\par\medskip « Il vient \textbf{aujourd'hui/ici} » (Garder proximal)};
\node[bad, right=of f5, anchor=north west] (e5) {\eng{\textcolor{poppink}{He comes \textbf{today/here}.}} \par\medskip \textcolor{popgreen}{\eng{He came \textbf{that day/there}.}}};
\draw[arrow] (f5.east) -- (e5.west) node[midway, above, font=\tiny, text=popdark] {Distanciation};
\end{tikzpicture}
\legende{Tableau récapitulatif des 5 pièges majeurs : forme fautive (rose) vs forme correcte (vert).}
\end{popfigure}

\courssub{Analyse détaillée de chaque piège}

\textbf{Piège 1 : L'oubli du \eng{backshift} (recul des temps).}\par
En français, à l'oral, on conserve souvent le futur (\eng{Il a dit qu'il viendra}). En anglais standard, le recul est \textbf{obligatoire} si le verbe introducteur est au passé (\eng{said, told, asked, thought, knew}).
\begin{exemplebox}[Backshift systématique]
\eng{Direct:} “I \textbf{will} defend human rights.” \\
\eng{Reported:} He said he \textbf{would} defend human rights. (Il a dit qu'il défendrait les droits de l'homme.) \\
\eng{Direct:} “The election \textbf{is} next month.” \\
\eng{Reported:} She announced the election \textbf{was} the following month.
\end{exemplebox}
\emph{Exception :} Si le fait rapporté est une vérité générale ou encore valable \textbf{au moment où l'on parle}, on peut garder le présent : \eng{The teacher said the Earth \textbf{revolves} around the Sun.}

\textbf{Piège 2 : La place de l'adverbe de manière (\eng{well, badly, quickly, fluently}).}\par
Français : Verbe + Adverbe + Objet (\eng{Il parle bien anglais}). Anglais : Verbe + Objet + Adverbe (\eng{He speaks English well}). C'est une règle absolue pour les adverbes de manière modifiant un verbe transitif.
\begin{propriete}[Ordre canonique : SVO + Adverbe]
\begin{center}
\begin{tabular}{|l|l|l|l|}
\hline
\rowcolor{popblueL} \textbf{Sujet} & \textbf{Verbe} & \textbf{Objet} & \textbf{Adverbe (Manière)} \\
\hline
The students & learn & their lessons & \textbf{quickly}. \\
\hline
The candidate & addressed & the crowd & \textbf{peacefully}. \\
\hline
\end{tabular}
\end{center}
\emph{Note :} Les adverbes de fréquence (\eng{always, often}) vont \textbf{avant} le verbe principal (mais après \eng{be/auxiliaires}) : \eng{He \textbf{always} respects the law.}
\end{propriete}

\textbf{Piège 3 : L'inversion sujet/verbe dans la question indirecte.}\par
En français, la question indirecte garde souvent l'inversion (\eng{Il a demandé où j'étais / où étais-je ?}). En anglais, la question indirecte redevient une phrase \textbf{affirmative} (Sujet + Verbe). Pas d'auxiliaire \eng{do/does/did}, pas d'inversion.
\begin{exemplebox}[Transformation Question Directe $\rightarrow$ Indirecte]
\begin{itemize}
    \item \eng{Direct:} “\textbf{Where are} the polling stations?” $\rightarrow$ \eng{Indirect:} He asked where the polling stations \textbf{were}.
    \item \eng{Direct:} “\textbf{When does} the meeting \textbf{start}?” $\rightarrow$ \eng{Indirect:} She wondered when the meeting \textbf{started}.
    \item \eng{Direct:} “\textbf{Did you vote}?” $\rightarrow$ \eng{Indirect:} He asked \textbf{if/whether I had voted}.
\end{itemize}
\end{exemplebox}

\textbf{Piège 4 : \eng{Tell} sans objet indirect.}\par
\eng{Tell} est transitif indirect : on \textbf{dit quelque chose \textit{à quelqu'un}}. \eng{Say} ne prend pas d'objet personne (on \textbf{dit quelque chose}).
\begin{center}
\begin{tabularx}{\linewidth}{|X|X|X|}
\hline
\rowcolor{popblueL} \textbf{Verbe} & \textbf{Structure} & \textbf{Exemple} \\
\hline
\eng{Say} & \eng{Say (that) + clause} & \eng{He \textbf{said} (that) he was ready.} \\
\hline
\eng{Tell} & \eng{Tell + \textbf{Person} + (that) + clause} & \eng{He \textbf{told me} (that) he was ready.} \\
\hline
\end{tabularx}
\end{center}
\begin{attention}[Erreur fatale]
\eng{\textcolor{poppink}{He told that he was a citizen.}} $\rightarrow$ \eng{He \textbf{told me} that he was a citizen.} \textbf{OU} \eng{He \textbf{said} that he was a citizen.}
\end{attention}

\textbf{Piège 5 : Les marqueurs de proximité (\eng{now, today, here, this, these}) non modifiés.}\par
Le discours rapporté implique une distanciation (décentrage). Les mots « proches » deviennent « lointains ».
\begin{center}
\begin{tabular}{|l|l|l|}
\hline
\rowcolor{popblueL} \textbf{Direct Speech (Proximal)} & \textbf{Reported Speech (Distal)} & \textbf{Français} \\
\hline
\eng{now} & \eng{then / at that time} & \eng{alors / à ce moment-là} \\
\hline
\eng{today / tonight} & \eng{that day / that night} & \eng{ce jour-là / cette nuit-là} \\
\hline
\eng{yesterday} & \eng{the day before / the previous day} & \eng{la veille} \\
\hline
\eng{tomorrow} & \eng{the next day / the following day} & \eng{le lendemain} \\
\hline
\eng{here} & \eng{there} & \eng{là-bas} \\
\hline
\eng{this / these} & \eng{that / those} & \eng{ce / ceux} \\
\hline
\eng{ago} & \eng{before} & \eng{auparavant} \\
\hline
\end{tabular}
\end{center}

\courssub{Méthode de vérification avant de rendre sa copie}

\begin{methode}[Check-list « Zéro Faute » pour le Discours Rapporté]
Appliquez ces 4 étapes systématiquement à chaque phrase transformée :
\begin{enumerate}
    \item \textbf{Verbe introducteur au passé ?} (\eng{said, told, asked, replied, promised, denied}...). Si OUI $\rightarrow$ Backshift obligatoire sur le verbe principal de la proposition rapportée.
    \item \textbf{Backshift effectué ?} Vérifiez : \eng{present $\rightarrow$ past, past $\rightarrow$ past perfect, will $\rightarrow$ would, can $\rightarrow$ could, must $\rightarrow$ had to}. Les modaux \eng{might, could, would, should, ought to} ne changent pas.
    \item \textbf{Adverbe bien placé ?} Si verbe transitif + objet : \eng{Verb + Object + Adverb}. (Ex: \eng{He reads the press \textbf{carefully}.})
    \item \textbf{Lieu avant Temps ?} Pour les compléments circonstanciels en fin de phrase : \eng{Place} avant \eng{Time}. \eng{He arrived \textbf{in Yaoundé} \textbf{yesterday}.} (Pas \eng{He arrived yesterday in Yaoundé} — possible mais moins naturel pour un lieu précis).
    \item \textbf{Mots « spatio-temporels » décalés ?} \eng{now $\rightarrow$ then, here $\rightarrow$ there, today $\rightarrow$ that day, this $\rightarrow$ that}.
    \item \textbf{Question indirecte ?} Ordre \textbf{Sujet + Verbe} (pas d'inversion, pas de \eng{do/does/did}). \eng{If/Whether} pour les questions fermées.
    \item \textbf{\eng{Tell} vs \eng{Say} ?} \eng{Told me/him/us} | \eng{Said (that)}.
\end{enumerate}
\end{methode}

\courssub{Exercice résolu : Transformation complète}

\begin{exoresolu}[Transformation Direct $\rightarrow$ Indirect : Discours électoral]
\textbf{Énoncé :} Transformez le discours direct suivant en discours rapporté (style indirect). Le verbe introducteur est \eng{The candidate declared that...}

\eng{Direct Speech:} “My opponents \textbf{are corrupt}. They \textbf{stole} public funds \textbf{last year}. I \textbf{will create} jobs for the youth \textbf{now}. \textbf{Do you trust} me? \textbf{Come} to my rally \textbf{here tomorrow}!”

\textbf{Solution détaillée :}
\begin{enumerate}
    \item \textbf{Verbe introducteur :} \eng{declared} (passé) $\rightarrow$ Backshift généralisé.
    \item \textbf{Phrase 1:} “My opponents \textbf{are corrupt}.” $\rightarrow$ \eng{his opponents \textbf{were} corrupt.} (Present $\rightarrow$ Past ; \eng{my} $\rightarrow$ \eng{his}).
    \item \textbf{Phrase 2:} “They \textbf{stole} public funds \textbf{last year}.” $\rightarrow$ \eng{they \textbf{had stolen} public funds \textbf{the previous year}.} (Past Simple $\rightarrow$ Past Perfect ; \eng{last year} $\rightarrow$ \eng{the previous year}).
    \item \textbf{Phrase 3:} “I \textbf{will create} jobs for the youth \textbf{now}.” $\rightarrow$ \eng{he \textbf{would create} jobs for the youth \textbf{then}.} (\eng{will} $\rightarrow$ \eng{would} ; \eng{I} $\rightarrow$ \eng{he} ; \eng{now} $\rightarrow$ \eng{then}).
    \item \textbf{Question (Phrase 4):} “\textbf{Do you trust} me?” $\rightarrow$ \eng{\textbf{if/whether they trusted} him.} (Question fermée $\rightarrow$ \eng{if/whether} ; \eng{do/you trust} (Present) $\rightarrow$ \eng{they trusted} (Past) ; \eng{you} $\rightarrow$ \eng{they} / \eng{me} $\rightarrow$ \eng{him}).
    \item \textbf{Impératif (Phrase 5):} “\textbf{Come} to my rally \textbf{here tomorrow}!” $\rightarrow$ \eng{\textbf{to come} to his rally \textbf{there the next day}.} (Impératif $\rightarrow$ \eng{to} + infinitif ; \eng{here} $\rightarrow$ \eng{there} ; \eng{tomorrow} $\rightarrow$ \eng{the next day}).
\end{enumerate}

\textbf{Résultat final :}
\eng{The candidate declared that his opponents \textbf{were} corrupt. He added that they \textbf{had stolen} public funds \textbf{the previous year}. He promised that he \textbf{would create} jobs for the youth \textbf{then}. He asked \textbf{if they trusted} him and urged them \textbf{to come} to his rally \textbf{there the next day}.}
\end{exoresolu}

\courssub{Entraînement guidé}

\begin{exercice}[Correction d'erreurs « Spécial Bac »]
Corrigez les erreurs dans les phrases suivantes (une erreur majeure par phrase). Réécrivez la phrase correcte.
\begin{enumerate}
    \item If the government \textbf{would invest} more in education, literacy rates \textbf{will} rise.
    \item The Minister \textbf{told that} the new law \textbf{will be} applied \textbf{today}.
    \item The journalist asked where \textbf{was} the polling station located.
    \item She speaks \textbf{fluently} English because she lived in London.
    \item He promised he \textbf{can} help us \textbf{next week}.
    \item They denied \textbf{to have} participated in the riot.
    \item The teacher asked \textbf{if we have finished} the test.
    \item \textbf{Say} to your brother \textbf{to come} here immediately.
\end{enumerate}
\end{exercice}

\begin{corrige}[Exercice : Correction d'erreurs]
\begin{enumerate}
    \item \eng{If the government \textbf{invested} more in education, literacy rates \textbf{would} rise.} (Piège \eng{If + would} + concordance 2\textsuperscript{e} conditionnel).
    \item \eng{The Minister \textbf{said that} the new law \textbf{would be} applied \textbf{that day}.} (Piège \eng{told that} $\rightarrow$ \eng{said that} ; \eng{will} $\rightarrow$ \eng{would} ; \eng{today} $\rightarrow$ \eng{that day}).
    \item \eng{The journalist asked where the polling station \textbf{was} located.} (Piège inversion question indirecte).
    \item \eng{She speaks English \textbf{fluently} because she lived in London.} (Piège place adverbe : Verbe + Objet + Adverbe).
    \item \eng{He promised he \textbf{could} help us \textbf{the following week}.} (Piège \eng{can} $\rightarrow$ \eng{could} ; \eng{next week} $\rightarrow$ \eng{the following week}).
    \item \eng{They denied \textbf{having participated} in the riot.} (Piège \eng{deny + V-ing}, pas \eng{to V}).
    \item \eng{The teacher asked \textbf{if we had finished} the test.} (Piège backshift Present Perfect $\rightarrow$ Past Perfect + ordre affirmatif).
    \item \eng{\textbf{Tell} your brother \textbf{to come} \textbf{there} immediately.} (Piège \eng{Say to} $\rightarrow$ \eng{Tell} ; Impératif indirect \eng{tell + obj + to V} ; \eng{here} $\rightarrow$ \eng{there}).
\end{enumerate}
\end{corrige}

\begin{savaistu}[Le Bac et les 5 pièges]
Au Baccalauréat série A (épreuve écrite), l'exercice de transformation \textbf{Direct Speech $\rightarrow$ Indirect Speech} (ou QCM de correction d'erreurs) cible presque systématiquement la combinaison de ces 5 pièges dans une seule phrase longue. Exemple type : \eng{“He told that he will come here tomorrow if he can.”} $\rightarrow$ 4 erreurs : \eng{told that} (manque objet), \eng{will} (pas de backshift), \eng{here tomorrow} (marqueurs non décalés), \eng{if he can} (backshift manquant \eng{could}). Maîtriser la check-list de la méthode ci-dessus garantit 4/4 points sur cet item.
\end{savaistu}

\begin{aretenir}[Synthèse : Français vs Anglais — Les 5 points de vigilance]
\begin{itemize}
    \item \textbf{Conditionnel après \eng{Si}} : FR \eng{verrais} $\neq$ EN \eng{would} dans la proposition \eng{if}. $\rightarrow$ \eng{If I \textbf{knew}, I \textbf{would} go.}
    \item \textbf{Backshift} : FR (oral) garde futur ; EN \textbf{exige} \eng{will $\rightarrow$ would, present $\rightarrow$ past}.
    \item \textbf{Adverbe (Manière)} : FR [V Adv O] ; EN \textbf{[V O Adv]}.
    \item \textbf{Question indirecte} : FR (inversion possible) ; EN \textbf{ordre affirmatif obligatoire} (S + V).
    \item \textbf{\eng{Tell} vs \eng{Say}} : FR \eng{dire à qqn} (un verbe) ; EN \textbf{\eng{Tell sb that} / \eng{Say that}}.
    \item \textbf{Déictiques} : FR (contexte partagé) ; EN \textbf{marquage morphologique de la distance} (\eng{now/then, here/there}).
\end{itemize}
\end{aretenir}

\coursec{Entraînement guidé et production}

Nous avons vu dans la section précédente les différences majeures entre le français et l'anglais : le recul des temps (\eng{backshift}) après un verbe introducteur au passé, l'adaptation des pronoms et des marqueurs temporels, et l'ordre fixe \cle{SVCLT} (Sujet -- Verbe -- Complément -- Lieu -- Temps). Il est temps de passer à la pratique. Cette dernière section vous propose un entraînement progressif : transformation de phrases, remise en ordre de mots, correction d'erreurs typiques, questions rapportées, puis une production écrite finale. Chaque exercice est suivi d'un corrigé détaillé : prenez le temps de comparer vos réponses et de comprendre \emph{pourquoi} une forme est correcte.

\courssub{Méthode : transformer une phrase au discours indirect}

Avant de commencer les exercices, fixons une démarche en quatre étapes que vous appliquerez mécaniquement jusqu'à ce qu'elle devienne un réflexe.

\begin{methode}[Transformer du discours direct au discours indirect]
\begin{enumerate}
\item \textbf{Repérer le verbe introducteur.} S'il est au présent (\eng{says, reports}), les temps ne changent pas. S'il est au passé (\eng{said, declared, announced}), le \cle{backshift} est obligatoire.
\item \textbf{Reculer le temps} du verbe rapporté d'un cran vers le passé : present simple $\rightarrow$ past simple ; present continuous $\rightarrow$ past continuous ; past simple / present perfect $\rightarrow$ past perfect ; \eng{will} $\rightarrow$ \eng{would} ; \eng{can} $\rightarrow$ \eng{could} ; \eng{must} $\rightarrow$ \eng{had to}.
\item \textbf{Adapter les pronoms, possessifs et marqueurs spatio-temporels} : \eng{I} $\rightarrow$ \eng{he/she} ; \eng{my} $\rightarrow$ \eng{his/her} ; \eng{now} $\rightarrow$ \eng{then} ; \eng{today} $\rightarrow$ \eng{that day} ; \eng{tomorrow} $\rightarrow$ \eng{the next day} ; \eng{yesterday} $\rightarrow$ \eng{the day before} ; \eng{here} $\rightarrow$ \eng{there}.
\item \textbf{Vérifier l'ordre des mots} : la phrase rapportée est affirmative (Sujet + Verbe), jamais interrogative ; l'adverbe se place selon la règle anglaise, pas française.
\end{enumerate}
\end{methode}

\begin{exemplebox}[Un exemple complet, étape par étape]
\eng{“We shall defend our rights tomorrow,” the students said.}
\begin{itemize}
\item Verbe introducteur : \eng{said} (passé) $\rightarrow$ backshift nécessaire.
\item \eng{shall/will defend} $\rightarrow$ \eng{would defend}.
\item \eng{we} $\rightarrow$ \eng{they} ; \eng{our} $\rightarrow$ \eng{their} ; \eng{tomorrow} $\rightarrow$ \eng{the next day}.
\item Ordre des mots : Sujet (\eng{they}) + Verbe (\eng{would defend}) + Complément (\eng{their rights}) + Temps (\eng{the next day}).
\end{itemize}
Résultat : \eng{The students said they would defend their rights the next day.} « Les étudiants ont dit qu'ils défendraient leurs droits le lendemain. »
\end{exemplebox}

\courssub{Exercice 1 : du discours direct au discours indirect}

\begin{exoresolu}[Reported speech — elections and civic rights]
\textbf{Énoncé.} Put the following sentences into reported speech. The introductory verb is in the past.

\begin{enumerate}
\item \eng{“The elections will be free and fair,” the president announced.}
\item \eng{“I have voted in every election since 1992,” the old woman said.}
\item \eng{“We are counting the ballots now,” the official told the journalists.}
\item \eng{“You must register before Friday,” the mayor warned the young people.}
\item \eng{“The police arrested two demonstrators yesterday,” the reporter said.}
\item \eng{“I can help you fill in the form,” the volunteer told the voter.}
\end{enumerate}

\tcblower

\textbf{Solution détaillée.}
\begin{enumerate}
\item \eng{The president announced (that) the elections would be free and fair.} (\eng{will} $\rightarrow$ \eng{would})
\item \eng{The old woman said (that) she had voted in every election since 1992.} (present perfect $\rightarrow$ past perfect ; \eng{I} $\rightarrow$ \eng{she})
\item \eng{The official told the journalists (that) they were counting the ballots then.} (present continuous $\rightarrow$ past continuous ; \eng{now} $\rightarrow$ \eng{then})
\item \eng{The mayor warned the young people (that) they had to register before Friday.} (\eng{must} $\rightarrow$ \eng{had to})
\item \eng{The reporter said (that) the police had arrested two demonstrators the day before.} (past simple $\rightarrow$ past perfect ; \eng{yesterday} $\rightarrow$ \eng{the day before})
\item \eng{The volunteer told the voter (that) he/she could help him/her fill in the form.} (\eng{can} $\rightarrow$ \eng{could})
\end{enumerate}
Remarquez que \eng{that} est facultatif après \eng{said} et \eng{told}, mais que \eng{tell} exige toujours un complément de personne : \eng{told the journalists}, jamais \eng{told that} seul.
\end{exoresolu}

\begin{attention}[Le piège de « tell » et « say »]
Les francophones confondent souvent \eng{say} et \eng{tell} car le français utilise « dire » pour les deux. Retenez : \cle{say something (to somebody)} mais \cle{tell somebody something}.
\begin{itemize}
\item \eng{He said that he was tired.} / \eng{He said to me that he was tired.} « Il a dit qu'il était fatigué. »
\item \eng{He told me that he was tired.} « Il m'a dit qu'il était fatigué. »
\item Jamais : \emph{He told that...} ni \emph{He said me that...}.
\end{itemize}
\end{attention}

\courssub{Exercice 2 : remise en ordre des mots (SVCLT)}

\begin{exercice}[Word order — build correct sentences]
Put the words in the correct order to make sentences. Remember: Subject -- Verb -- Object -- Place -- Time.

\begin{enumerate}
\item \eng{the citizens / peacefully / demonstrated / in front of the town hall / yesterday}
\item \eng{speaks / the mayor / always / clearly / during press conferences}
\item \eng{signed / the minister / the new law / in Yaoundé / last week}
\item \eng{human rights / must / we / defend / every day}
\item \eng{the journalist / asked / politely / a question / at the meeting}
\end{enumerate}
\end{exercice}

\begin{corrige}[Exercice 2]
\begin{enumerate}
\item \eng{The citizens demonstrated peacefully in front of the town hall yesterday.} (S + V + manière + lieu + temps)
\item \eng{The mayor always speaks clearly during press conferences.} (adverbe de fréquence \eng{always} avant le verbe lexical)
\item \eng{The minister signed the new law in Yaoundé last week.} (lieu avant temps)
\item \eng{We must defend human rights every day.} (le complément \eng{human rights} colle au verbe ; le temps en fin de phrase)
\item \eng{The journalist politely asked a question at the meeting.} ou \eng{The journalist asked a question politely at the meeting.}
\end{enumerate}
\end{corrige}

\begin{center}
\begin{tabularx}{\linewidth}{|X|X|X|}
\hline
\rowcolor{popblueL} Position en français & Position en anglais & Exemple anglais \\
\hline
Adverbe souvent après le verbe & Adverbe de fréquence \textbf{avant} le verbe & \eng{She always votes.} \\
\hline
« Il parle bien anglais » (adverbe entre verbe et COD) & Jamais entre le verbe et son COD & \eng{He speaks English well.} \\
\hline
Temps souvent avant le lieu & Lieu \textbf{avant} temps & \eng{in Buea last night} \\
\hline
Phrase interrogative : inversion & Question rapportée : ordre affirmatif & \eng{He asked where I lived.} \\
\hline
\end{tabularx}
\end{center}

\courssub{Exercice 3 : corriger les erreurs typiques des francophones}

\begin{exercice}[Correct the mistakes]
Each sentence contains one typical mistake made by French speakers. Correct it.

\begin{enumerate}
\item \eng{The mayor said that he will come tomorrow.}
\item \eng{She told that the elections were postponed.}
\item \eng{He explained me the new voting procedure.}
\item \eng{They asked me where did I live.}
\item \eng{The deputy speaks always about human rights.}
\end{enumerate}
\end{exercice}

\begin{corrige}[Exercice 3]
\begin{enumerate}
\item \eng{The mayor said that he \textbf{would} come \textbf{the next day}.} — Après \eng{said} (passé), \eng{will} recule en \eng{would} et \eng{tomorrow} devient \eng{the next day}.
\item \eng{She \textbf{said} that the elections were postponed.} (ou \eng{She told \textbf{us} that...}) — \eng{tell} exige un complément de personne.
\item \eng{He explained the new voting procedure \textbf{to me}.} — \eng{explain} ne prend jamais de complément de personne direct : on dit \eng{explain something to somebody}, contrairement à « expliquer quelque chose à quelqu'un » où le pronom français « me » précède le verbe.
\item \eng{They asked me where I \textbf{lived}.} — Dans une question rapportée, on revient à l'ordre affirmatif (sujet + verbe), sans inversion ni auxiliaire \eng{did}.
\item \eng{The deputy \textbf{always speaks} about human rights.} — L'adverbe de fréquence se place avant le verbe lexical, contrairement au français « parle toujours ».
\end{enumerate}
\end{corrige}

\begin{attention}[Le test de la lecture à voix haute]
Relisez toujours votre phrase anglaise \textbf{à voix haute}. Si l'adverbe « sonne » comme en français (\eng{he speaks fluently English}, \eng{she watches often television}), il est très probablement mal placé. En anglais, rien ne s'intercale entre le verbe et son complément d'objet direct : \eng{He speaks English fluently.} « Il parle couramment l'anglais. »
\end{attention}

\courssub{Exercice 4 : les questions rapportées}

Rappel de la règle : pour rapporter une question fermée (réponse oui/non), on utilise \cle{if} ou \cle{whether} ; pour une question ouverte, on garde le mot interrogatif (\eng{where, when, why, who, how}). Dans les deux cas, l'ordre redevient affirmatif et le backshift s'applique.

\begin{exoresolu}[Reported questions]
\textbf{Énoncé.} Report these questions. Begin with the words given.

\begin{enumerate}
\item \eng{“Do you have your voter's card?” the officer asked me. $\rightarrow$ The officer asked me...}
\item \eng{“Where did you register?” my friend asked me. $\rightarrow$ My friend asked me...}
\item \eng{“Will the results be announced tonight?” the journalist asked. $\rightarrow$ The journalist asked...}
\item \eng{“Why are human rights important?” the teacher asked the class. $\rightarrow$ The teacher asked the class...}
\end{enumerate}

\tcblower

\textbf{Solution détaillée.}
\begin{enumerate}
\item \eng{The officer asked me \textbf{if/whether} I had my voter's card.} (question fermée $\rightarrow$ \eng{if} ; present simple $\rightarrow$ past simple)
\item \eng{My friend asked me \textbf{where} I had registered.} (mot interrogatif conservé ; past simple $\rightarrow$ past perfect ; ordre affirmatif)
\item \eng{The journalist asked \textbf{if/whether} the results would be announced that night.} (\eng{will} $\rightarrow$ \eng{would} ; \eng{tonight} $\rightarrow$ \eng{that night})
\item \eng{The teacher asked the class \textbf{why} human rights were important.} (ordre affirmatif : \eng{human rights were}, pas \emph{were human rights} ; ici le backshift est facultatif car la phrase exprime une vérité générale)
\end{enumerate}
\end{exoresolu}

\begin{propriete}[Question rapportée = phrase affirmative]
Dans une question rapportée, il n'y a \textbf{ni inversion, ni auxiliaire do/does/did, ni point d'interrogation}. Comparez :
\begin{itemize}
\item Question directe : \eng{Where do you live?} « Où habites-tu ? »
\item Question rapportée : \eng{He asked me where I lived.} « Il m'a demandé où j'habitais. »
\end{itemize}
\end{propriete}

\courssub{Carte mentale récapitulative}

\begin{popfigure}
\begin{tikzpicture}[mindmap, grow cyclic,
  every node/.style={concept, align=center, font=\small},
  root concept/.append style={concept color=popblue, text=white, font=\small\bfseries},
  level 1 concept/.append style={sibling angle=90, level distance=3.6cm}]
\node[root concept]{Sequence of tenses \& word order}
  child[concept color=poporange] { node[align=center]{Backshift\\ will $\rightarrow$ would\\ can $\rightarrow$ could\\ present $\rightarrow$ past} }
  child[concept color=popgreen] { node[align=center]{Exceptions\\ vérités générales\\ faits encore vrais\\ verbe au présent} }
  child[concept color=poppurple] { node[align=center]{SVCLT\\ Sujet + Verbe\\ + Complément\\ + Lieu + Temps} }
  child[concept color=poppink] { node[align=center]{Pièges\\ say / tell\\ explain to\\ ordre affirmatif\\ adverbe avant verbe} };
\end{tikzpicture}
\legende{Carte mentale : les quatre points clés de la concordance des temps et de l'ordre des mots.}
\end{popfigure}

\courssub{Production finale : un compte rendu au discours indirect}

\begin{savaistu}[Les grands discours des droits, un corpus d'entraînement idéal]
Les grands discours sur les droits — ceux de Martin Luther King (\eng{I Have a Dream}, 1963) ou de Nelson Mandela lors de son procès de 1964 — sont d'excellents textes pour s'entraîner au discours rapporté : ils sont construits en phrases courtes au présent et au futur (\eng{I have a dream... we shall overcome...}), parfaites pour pratiquer le backshift. Essayez : \eng{King said that he had a dream and that one day all men would be equal.} « King a dit qu'il faisait un rêve et qu'un jour tous les hommes seraient égaux. »
\end{savaistu}

\begin{experience}[Writing task — the mayor's speech]
\textbf{Situation.} Yesterday, the mayor of your town gave a speech on citizenship during the celebration of National Youth Day. You are the reporter for your school magazine. Write a short report (6--8 lines) in reported speech.

\textbf{The mayor's words (direct speech):}

\eng{“Dear citizens, I am proud of our young people. You are the future of this nation. We have built three new schools this year, and we shall build a public library next year. Every citizen must respect the law and defend human rights. I ask all of you to register on the electoral roll before December. Together, we can make our town a model of democracy.”}

\textbf{Guidance.}
\begin{enumerate}
\item Choose your introductory verbs: \eng{the mayor said / declared / announced / reminded / urged / concluded...}
\item Apply the backshift systematically.
\item Adapt pronouns (\eng{I} $\rightarrow$ \eng{he}, \eng{you} $\rightarrow$ \eng{we/the young people}) and time markers (\eng{this year} $\rightarrow$ \eng{that year}, \eng{next year} $\rightarrow$ \eng{the following year}).
\item Vary the reporting verbs to avoid repeating \eng{said}.
\end{enumerate}
\end{experience}

\begin{exemplebox}[Modèle de production]
\eng{Yesterday, the mayor gave a speech on citizenship for National Youth Day. He declared that he was proud of the young people and said that they were the future of the nation. He announced that the council had built three new schools that year and promised that they would build a public library the following year. He reminded the audience that every citizen had to respect the law and defend human rights. Then he urged all the young people to register on the electoral roll before December. He concluded that, together, they could make their town a model of democracy.}

« Hier, le maire a prononcé un discours sur la citoyenneté pour la Journée nationale de la jeunesse. Il a déclaré qu'il était fier des jeunes et a dit qu'ils étaient l'avenir de la nation. Il a annoncé que la municipalité avait construit trois nouvelles écoles cette année-là et a promis qu'elle construirait une bibliothèque publique l'année suivante. Il a rappelé à l'auditoire que chaque citoyen devait respecter la loi et défendre les droits de l'homme. Puis il a exhorté tous les jeunes à s'inscrire sur les listes électorales avant décembre. Il a conclu qu'ensemble, ils pouvaient faire de leur ville un modèle de démocratie. »
\end{exemplebox}

\begin{aretenir}[L'essentiel de la leçon]
\begin{itemize}
\item Verbe introducteur au passé $\rightarrow$ \cle{backshift} : chaque temps recule d'un cran vers le passé.
\item Pronoms, possessifs et marqueurs temporels s'adaptent au nouveau point de vue (\eng{tomorrow} $\rightarrow$ \eng{the next day}).
\item Question rapportée : \eng{if/whether} (question fermée) ou mot interrogatif (question ouverte), puis \textbf{ordre affirmatif}.
\item Ordre des mots anglais : \cle{Sujet -- Verbe -- Complément -- Lieu -- Temps} ; l'adverbe de fréquence précède le verbe lexical.
\item Pièges des francophones : \eng{tell} + personne, \eng{explain something to somebody}, pas d'inversion dans les questions rapportées, pas d'adverbe entre le verbe et son COD.
\end{itemize}
\end{aretenir}

\newpage

\coursec{Activité d'intégration}

\courssub{Objectifs de l'activité}

À l'issue de cette activité d'intégration, tu seras capable de :

\begin{itemize}
\item \textbf{Comprendre} un court discours oral au discours direct et en relever les informations essentielles (prises de notes).
\item \textbf{Rapporter fidèlement} ce discours au discours indirect en appliquant la \cle{concordance des temps} (backshift), le changement des pronoms, des possessifs et des marqueurs spatio-temporels.
\item \textbf{Respecter l'ordre des mots anglais} dans les phrases déclaratives et les subordonnées introduites par \eng{that} (sujet -- verbe -- complément ; pas d'inversion).
\item \textbf{Distinguer les cas sans backshift} : vérités générales, faits encore vrais au moment du compte rendu.
\item \textbf{Rédiger} un court article de presse de 8 à 10 lignes pour le journal du lycée.
\item \textbf{Auto-évaluer} ta production à l'aide de la grille de vérification en 4 points.
\end{itemize}

\courssub{Situation}

Ton lycée, le \emph{Government Bilingual High School} de Buea, a invité le \emph{Minister of Human Rights} dans le cadre de la \emph{Human Rights Week}. Le club des jeunes reporters (\emph{Young Reporters Club}) couvre l'événement pour le journal du lycée, \emph{The Student Voice}. Tu fais partie de l'équipe de presse. Le Ministre vient de prononcer un court discours devant les élèves ; ta mission : rédiger, avec ton groupe, l'article qui paraîtra demain.

\begin{textebox}[The Minister's speech — direct speech (script à écouter / à lire)]
“Good morning, students. I am very happy to be in your school today.\\
First, I want to tell you that our ministry is building a new human rights centre in Buea.\\
We have already trained two hundred young volunteers this year.\\
Next month, we will launch a national campaign against discrimination in schools.\\
I believe that young people are the future of democracy.\\
Finally, remember this : all citizens are equal before the law, and your rights must always be respected.”
\end{textebox}

\begin{attention}[Avant de commencer]
Le discours ci-dessus contient volontairement les trois cas étudiés dans la leçon : des \textbf{faits passés} (\eng{have already trained}), des \textbf{promesses et projets} (\eng{will launch}), et une \textbf{vérité générale} (\eng{all citizens are equal}). Ton article devra traiter chaque cas correctement : backshift pour les deux premiers, \textbf{pas de backshift obligatoire} pour la vérité générale.
\end{attention}

\courssub{Consignes}

Travail en groupes de 4 : un élève joue le rôle du Ministre, les trois autres sont journalistes.

\begin{enumerate}
\item \textbf{Répétition du discours.} L'élève-Ministre lit le discours à voix haute, deux fois, sans le montrer. Les journalistes prennent des notes en anglais (mots-clés uniquement : verbes, chiffres, dates). Interdiction de recopier des phrases entières.
\item \textbf{Reconstitution.} Sans regarder le script, les journalistes reconstituent oralement les cinq informations du discours. Le « Ministre » confirme ou corrige.
\item \textbf{Choix des verbes introducteurs.} Pour chaque information, choisissez le verbe de parole adapté : \eng{announce}, \eng{promise}, \eng{remind}, \eng{explain}, \eng{declare}, \eng{add}. Attention : \eng{promise that...} pour un engagement futur, \eng{remind us that...} pour une vérité générale.
\item \textbf{Rédaction de l'article.} Rédigez un article de 8 à 10 lignes au discours indirect. Structure attendue : un titre, une phrase d'introduction (\eng{Yesterday, the Minister of Human Rights visited our school...}), puis le compte rendu des cinq informations.
\item \textbf{Vérification collective.} Relisez votre article avec la grille de vérification en 4 points (voir Ressources). Corrigez au stylo rouge.
\item \textbf{Lecture publique.} Chaque groupe lit son article à la classe. Les autres groupes lèvent la main à chaque erreur de concordance ou d'ordre des mots détectée, puis proposent une correction.
\end{enumerate}

\courssub{Ressources à mobiliser}

\begin{aretenir}[La grille de vérification en 4 points]
\begin{enumerate}
\item \textbf{Backshift} : present simple $\rightarrow$ past simple ; present perfect $\rightarrow$ past perfect ; will $\rightarrow$ would ; can $\rightarrow$ could.
\item \textbf{Marqueurs temporels} : today $\rightarrow$ that day ; yesterday $\rightarrow$ the day before ; next month $\rightarrow$ the following month ; here $\rightarrow$ there.
\item \textbf{Vérités générales} : pas de backshift obligatoire (\eng{He reminded us that all citizens are equal before the law.}).
\item \textbf{Ordre des mots} : dans la subordonnée en \eng{that}, toujours sujet + verbe, jamais d'inversion : \eng{He said that the centre was ready.} (et jamais \eng{He said that was the centre ready.}).
\end{enumerate}
\end{aretenir}

\begin{center}
\begin{tabularx}{\linewidth}{|X|X|}
\hline
\rowcolor{popblueL} \textbf{Discours direct (Ministre)} & \textbf{Discours indirect (journaliste)} \\
\hline
“I am very happy...” & He said (that) he was very happy. \\
\hline
“...is building a new centre...” & He announced that his ministry was building a new centre. \\
\hline
“We have trained two hundred volunteers.” & He explained that they had trained two hundred volunteers. \\
\hline
“We will launch a campaign.” & He promised that they would launch a campaign. \\
\hline
“All citizens are equal.” & He reminded us that all citizens are equal. \\
\hline
\end{tabularx}
\end{center}

\textbf{Lexique utile} : \eng{a press conference} (une conférence de presse), \eng{a reporter} (un reporter), \eng{to take notes} (prendre des notes), \eng{an article} (un article), \eng{human rights} (les droits de l'homme), \eng{equal before the law} (égaux devant la loi), \eng{to launch a campaign} (lancer une campagne), \eng{a volunteer} (un bénévole), \eng{discrimination} (la discrimination).

\courssub{Proposition de production et corrigé commenté}

\begin{exoresolu}[Article modèle pour \emph{The Student Voice}]
Rédige un article de 8 à 10 lignes rapportant le discours du Ministre au discours indirect.
\tcblower
\textbf{The Minister of Human Rights Visits Our School}

\eng{Yesterday, the Minister of Human Rights visited Government Bilingual High School, Buea, and gave a speech to the students. He said that he was very happy to be in our school that day. First, he announced that his ministry was building a new human rights centre in Buea. He also explained that they had already trained two hundred young volunteers this year. Then, he promised that they would launch a national campaign against discrimination in schools the following month. He added that he believed young people were the future of democracy. Finally, he reminded us that all citizens are equal before the law and that our rights must always be respected. The students listened carefully and asked many questions.}

\medskip
\textbf{Commentaire des choix de langue :}
\begin{itemize}
\item \eng{said that he was} : backshift du present simple (\eng{am} $\rightarrow$ \eng{was}) ; \eng{that} peut être omis mais il est conseillé de le garder à l'écrit.
\item \eng{that day} : transformation obligatoire de \eng{today} au discours indirect.
\item \eng{was building} : le present continu devient past continu (backshift).
\item \eng{had already trained} : present perfect $\rightarrow$ past perfect ; l'adverbe \eng{already} reste placé entre auxiliaire et participe, conformément à l'ordre des mots anglais.
\item \eng{would launch} : \eng{will} $\rightarrow$ \eng{would}, avec le verbe introducteur \eng{promised}, qui exprime l'engagement.
\item \eng{the following month} : \eng{next month} transformé.
\item \eng{are equal} : \textbf{pas de backshift}, car il s'agit d'une vérité générale toujours vraie ; \eng{must} est inchangé (pas de forme passée).
\item Ordre des mots respecté partout : sujet + verbe dans chaque subordonnée, aucune inversion, aucun mot français calqué.
\end{itemize}
\end{exoresolu}

\courssub{Bilan de l'activité}

Cette activité t'a permis de mobiliser en situation réelle les deux compétences centrales de la leçon :

\begin{itemize}
\item \textbf{La concordance des temps} : tu as transformé un discours authentique en compte rendu journalistique, en appliquant le backshift aux faits et aux promesses, et en le suspendant pour la vérité générale (\eng{all citizens are equal}).
\item \textbf{L'ordre des mots} : tu as construit des subordonnées en \eng{that} sans inversion, avec les adverbes et les marqueurs temporels à leur place.
\end{itemize}

\begin{attention}[Auto-évaluation finale]
Avant de rendre ton article, pose-toi quatre questions : (1) Ai-je fait le backshift partout où il est nécessaire ? (2) Ai-je transformé \eng{today}, \eng{next month}, \eng{here} ? (3) Ai-je gardé le présent pour les vérités générales ? (4) L'ordre sujet--verbe est-il respecté dans chaque phrase ? Si une réponse est « non », corrige avant la lecture publique.
\end{attention}

Le savoir-faire développé ici — rapporter fidèlement les paroles d'autrui — est exactement celui attendu à l'épreuve d'anglais du baccalauréat, notamment dans les exercices de transformation direct/indirect et dans les productions écrites de type article ou reportage. Garde la grille de vérification en 4 points dans ton cahier : elle te servira pour toutes les leçons de révision.

\newpage

\coursec{Méthodes et exercices résolus}

\coursec{Lesson 44 — Grammar: Sequence of tenses, word order}

\courssub{Observation et découverte}

\begin{textebox}[Interview d'une défenseure des droits humains à Buea]
A journalist from \eng{The Cameroon Tribune} interviewed Mrs. Ayuk, a human rights defender, in Buea yesterday. Here is an extract of the conversation:

\eng{Journalist:} “Why do you defend human rights?”
\eng{Mrs. Ayuk:} “I work for a small organisation called \eng{Voice of the Voiceless}. We have helped many victims since 2016. The situation is difficult here because the authorities do not always listen. They must understand that rights are not privileges. I will never give up.”

\eng{Journalist:} “What happened after the last protest?”
\eng{Mrs. Ayuk:} “When the police arrived, the crowd had already dispersed. If the government had listened earlier, there would have been no tension.”
\end{textebox}

\begin{methode}[Comment observer la grammaire dans un texte authentique]
\begin{enumerate}
\item \textbf{Surligne} tous les verbes des paroles rapportées (ce que Mme Ayuk a dit).
\item \textbf{Identifie} le verbe introducteur (\eng{said, replied, added, declared}) et son temps (ici : \textbf{past simple}).
\item \textbf{Compare} le temps du verbe original (discours direct) et celui du discours rapporté (dans l'article). Quel décalage observes-tu ?
\item \textbf{Repère} l'ordre des mots dans les questions (\eng{Why \textbf{do} you defend...?}) et dans les phrases affirmatives (\eng{I \textbf{work} for...}).
\item \textbf{Note} la place des compléments : \eng{helped \textbf{many victims} \textbf{since 2016}} (COD + Temps), \eng{arrived \textbf{when the police arrived}} (subordonnée temporelle).
\end{enumerate}
\end{methode}

\begin{exemplebox}[Analyse guidée de l'extrait]
\begin{itemize}
\item \eng{“I \textbf{work}…”} $\rightarrow$ She \textbf{said} she \textbf{worked}… (Present simple $\rightarrow$ Past simple)
\item \eng{“We \textbf{have helped}…”} $\rightarrow$ She \textbf{explained} they \textbf{had helped}… (Present perfect $\rightarrow$ Past perfect)
\item \eng{“The situation \textbf{is}…”} $\rightarrow$ She \textbf{added} the situation \textbf{was}… (Present $\rightarrow$ Past)
\item \eng{“They \textbf{must} understand…”} $\rightarrow$ She \textbf{said} they \textbf{had to} understand… (\eng{must} $\rightarrow$ \eng{had to})
\item \eng{“I \textbf{will} never give up.”} $\rightarrow$ She \textbf{promised} she \textbf{would} never give up. (\eng{will} $\rightarrow$ \eng{would})
\item \eng{“When the police \textbf{arrived}, the crowd \textbf{had already dispersed}.”} $\rightarrow$ Past simple + Past perfect (antériorité).
\item \eng{“If the government \textbf{had listened}…, there \textbf{would have been}…”} $\rightarrow$ 3\textsuperscript{e} conditionnel (irréel du passé).
\end{itemize}
\end{exemplebox}

\courssub{La concordance des temps : la règle (Backshifting)}

\begin{propriete}[Règle générale de la concordance des temps]
Lorsque le verbe introducteur (\eng{say, tell, ask, think, know, believe, explain, promise, declare, reply, add}…) est à un \textbf{temps du passé} (past simple, past perfect, past continuous), les temps des verbes dans la subordonnée \textbf{reculent d'un cran} (backshifting). Cela s'applique au discours rapporté (\eng{reported speech}) et aux subordonnées complétives (\eng{that}-clauses).
\end{propriete}

\begin{center}
\begin{tabularx}{\linewidth}{|X|X|X|}
\hline
\rowcolor{popblueL} \textbf{Discours direct / Temps original} & \textbf{Après verbe au passé (Backshifting)} & \textbf{Exemple traduit} \\
\hline
Present simple (\eng{I \textbf{work}}) & Past simple (\eng{he \textbf{worked}}) & « Il a dit qu'il travaillait. » \\
\hline
Present continuous (\eng{I \textbf{am working}}) & Past continuous (\eng{he \textbf{was working}}) & « Il a dit qu'il était en train de travailler. » \\
\hline
Present perfect (\eng{I \textbf{have finished}}) & Past perfect (\eng{he \textbf{had finished}}) & « Il a dit qu'il avait fini. » \\
\hline
Past simple (\eng{I \textbf{saw}}) & Past perfect (\eng{he \textbf{had seen}}) & « Il a dit qu'il avait vu. » \\
\hline
\eng{will} (\eng{I \textbf{will} come}) & \eng{would} (\eng{he \textbf{would} come}) & « Il a dit qu'il viendrait. » \\
\hline
\eng{can} (\eng{I \textbf{can} help}) & \eng{could} (\eng{he \textbf{could} help}) & « Il a dit qu'il pouvait aider. » \\
\hline
\eng{may} (\eng{It \textbf{may} rain}) & \eng{might} (\eng{it \textbf{might} rain}) & « Il a dit qu'il pourrait pleuvoir. » \\
\hline
\eng{must} (\eng{I \textbf{must} go}) & \eng{had to} (\eng{he \textbf{had to} go}) & « Il a dit qu'il devait partir. » \\
\hline
\eng{shall} (\eng{Shall \textbf{I} help?}) & \eng{should} (\eng{he \textbf{should} help}) & « Il a demandé s'il devait aider. » \\
\hline
\end{tabularx}
\end{center}

\begin{methode}[Bien construire une subordonnée après un verbe principal au passé]
\begin{enumerate}
\item \textbf{Repérer} le temps du verbe principal : est-il au passé (\eng{said, told me, believed, wondered}) ?
\item \textbf{Identifier} le temps du verbe de la subordonnée dans la phrase d'origine (ou dans la pensée du locuteur).
\item \textbf{Reculer d'un cran} selon le tableau ci-dessus.
\item \textbf{Adapter} les pronoms (\eng{I} $\rightarrow$ \eng{he/she}, \eng{we} $\rightarrow$ \eng{they}, \eng{my} $\rightarrow$ \eng{his/her}, \eng{our} $\rightarrow$ \eng{their}) et les déterminants possessifs.
\item \textbf{Adapter} les marqueurs spatio-temporels :
\begin{center}
\begin{tabular}{|l|l|}
\hline
\rowcolor{poporangeL} \textbf{Direct} & \textbf{Rapporté} \\
\hline
\eng{now} & \eng{then} \\
\hline
\eng{today / tonight} & \eng{that day / that night} \\
\hline
\eng{yesterday} & \eng{the day before / the previous day} \\
\hline
\eng{tomorrow} & \eng{the next day / the following day} \\
\hline
\eng{last week / month} & \eng{the week / month before / the previous week / month} \\
\hline
\eng{next week / month} & \eng{the following week / month} \\
\hline
\eng{here} & \eng{there} \\
\hline
\eng{this / these} & \eng{that / those} \\
\hline
\eng{ago} & \eng{before} \\
\hline
\end{tabular}
\end{center}
\end{enumerate}
\textbf{Réflexe Bac :} Souligne le verbe introducteur. S'il est au passé, le recul est \textbf{obligatoire} (sauf exceptions ci-dessous). Vérifie le triptyque : \cle{Temps} — \cle{Pronoms} — \cle{Marqueurs}.
\end{methode}

\courssub{Exceptions et cas particuliers}

\begin{propriete}[Cas où le recul n'est pas obligatoire]
\begin{enumerate}
\item \textbf{Vérité générale / fait scientifique permanent} : le présent reste présent.
\eng{The teacher said that the Earth \textbf{revolves} around the Sun.}
\item \textbf{Report immédiat} (le discours est rapporté au moment même ou très peu après, le verbe introducteur est souvent au présent \eng{says} ou au passé très récent) : les temps peuvent ne pas changer.
\eng{She \textbf{says} she \textbf{is} tired.} / \eng{He just \textbf{told} me he \textbf{is} busy.}
\item \textbf{Verbes modaux n'ayant pas de forme passée distincte ou ne changeant pas} : \eng{could, might, should, would, ought to, used to} restent identiques.
\eng{He said he \textbf{could} swim.} (pas de changement)
\item \textbf{Subordonnées introduites par \eng{when} (valeur temporelle) ou \eng{if/whether} (question indirecte)} : la règle du "pas de \eng{will}" s'applique (voir section suivante), mais le recul s'opère si le principal est au passé (\eng{will} $\rightarrow$ \eng{would}, \eng{can} $\rightarrow$ \eng{could}).
\end{enumerate}
\end{propriete}

\begin{attention}[Piège : \eng{must} et \eng{need}]
\begin{itemize}
\item \eng{must} n'a pas de passé $\rightarrow$ \eng{had to} (obligation passée) ou \eng{must} (déduction logique, rare en discours rapporté).
\item \eng{need} (verbe principal) $\rightarrow$ \eng{needed to} ; \eng{need} (semi-modal, surtout négatif \eng{need not}) $\rightarrow$ \eng{did not need to} ou \eng{need not have + V3} (selon le sens).
\end{itemize}
\end{attention}

\courssub{Gérer les temps dans les subordonnées temporelles et conditionnelles}

\begin{propriete}[Règle d'or : pas de \eng{will} après \eng{when, if, after, before, as soon as, until, unless}]
Dans une subordonnée de temps ou de condition se référant au \textbf{futur}, l'anglais utilise le \textbf{present simple} (ou present perfect pour l'antériorité), jamais \eng{will} / \eng{shall}.
\begin{itemize}
\item \eng{When I \textbf{turn} eighteen, I \textbf{will vote}.} (Pas *\eng{will turn})
\item \eng{As soon as the results \textbf{are} published, we \textbf{will know}.}
\item \eng{If the government \textbf{listens}, things \textbf{will improve}.} (1\textsuperscript{er} conditionnel)
\end{itemize}
\end{propriete}

\begin{propriete}[Antériorité dans le passé : Past Perfect]
Pour exprimer qu'une action est terminée \textbf{avant} une autre action passée (marquée par \eng{after, before, when, as soon as}), on utilise le \textbf{past perfect} (\eng{had + V3}) dans la subordonnée.
\eng{After the demonstrators \textbf{had signed} the petition, they \textbf{handed} it to the mayor.}
\eng{The journalist \textbf{left} as soon as she \textbf{had finished} her interview.}
\end{propriete}

\begin{propriete}[Les trois types de conditionnelles]
\begin{center}
\begin{tabularx}{\linewidth}{|X|X|X|X|}
\hline
\rowcolor{popgreenL} \textbf{Type} & \textbf{Si-clause (Condition)} & \textbf{Main Clause (Conséquence)} & \textbf{Sens / Exemple} \\
\hline
Type 0 (Général) & Present Simple & Present Simple & \eng{If you \textbf{heat} ice, it \textbf{melts}.} (Vérité scientifique) \\
\hline
Type 1 (Réel futur) & Present Simple & \eng{will} + BV & \eng{If it \textbf{rains}, we \textbf{will stay} home.} (Possible) \\
\hline
Type 2 (Irréel présent) & Past Simple & \eng{would} + BV & \eng{If I \textbf{knew}, I \textbf{would tell} you.} (Hypothétique, peu probable) \\
\hline
Type 3 (Irréel passé) & Past Perfect & \eng{would have} + V3 & \eng{If she \textbf{had studied}, she \textbf{would have passed}.} (Regret, impossible) \\
\hline
Mixte (Passé $\rightarrow$ Présent) & Past Perfect & \eng{would} + BV & \eng{If I \textbf{had taken} the job, I \textbf{would be} rich now.} \\
\hline
\end{tabularx}
\end{center}
\end{propriete}

\begin{popfigure}
\begin{tikzpicture}[
    node distance=1.2cm and 2cm,
    box/.style={rectangle, draw=popblue, thick, fill=popblueL, rounded corners, align=center, minimum width=3cm, minimum height=1.2cm},
    arrow/.style={-Stealth, thick, popdark}
]
\node[box, align=center] (direct){Discours Direct\\Présent / Futur};
\node[box, right=of direct, align=center] (reporting){Verbe Introducteur\\AU PASSÉ};
\node[box, right=of reporting, align=center] (indirect){Discours Rapporté\\RECUL D'UN CRAN};
\draw[arrow] (direct) -- (reporting) node[midway, above, font=\small] {Backshifting};
\draw[arrow] (reporting) -- (indirect);
\node[box, below=of reporting, fill=poporangeL, draw=poporange, align=center] (exceptions){EXCEPTIONS\\Vérités générales\\Report immédiat\\Modaux invariables};
\draw[arrow, dashed] (reporting) -- (exceptions);
\end{tikzpicture}
\legende{Schéma de la concordance des temps (Backshifting) et exceptions}
\end{popfigure}

\courssub{L'ordre des mots en anglais (Word Order)}

\begin{propriete}[Structure fondamentale de la phrase affirmative : SVOMPT]
L'anglais est une langue à \textbf{ordre fixe} (rigide). Le squelette de base est :
\cle{Subject} + \cle{Verb} + \cle{Object} + \cle{Manner} + \cle{Place} + \cle{Time}
(\textbf{S V O M P T})
\begin{itemize}
\item Le \textbf{Sujet} est obligatoire (même "dummy" \eng{it/there}).
\item Le \textbf{Verbe} suit immédiatement le sujet (pas de COD intercalé).
\item Les \textbf{Compléments} suivent l'ordre : \textbf{Manière} (comment?) $\rightarrow$ \textbf{Lieu} (où?) $\rightarrow$ \textbf{Temps} (quand?).
\item L'\textbf{Adjectif épithète} est \textbf{toujours placé AVANT le nom} (\eng{a democratic country}, jamais *\eng{a country democratic}).
\end{itemize}
\end{propriete}

\begin{exemplebox}[Application de l'ordre S V O M P T]
\begin{tabular}{ll}
\eng{The minister spoke \textbf{calmly} (M) \textbf{at the press conference} (P) \textbf{yesterday} (T).} & Correct. \\
\eng{The minister spoke \textbf{at the press conference} \textbf{calmly} \textbf{yesterday}.} & Incorrect (P avant M). \\
\eng{\textbf{In Yaoundé} (P), the students \textbf{organised} (V) \textbf{a debate} (O) \textbf{peacefully} (M) \textbf{last week} (T).} & Correct (Complément de lieu en tête $\rightarrow$ virgule, ordre SVOMPT respecté après).
\end{tabular}
\end{exemplebox}

\begin{propriete}[Place des adverbes de fréquence et de degré]
\begin{itemize}
\item \textbf{Fréquence} (\eng{always, usually, often, sometimes, rarely, never}) : \textbf{Avant le verbe principal} ; \textbf{Après l'auxiliaire / \eng{be} / modaux}.
\eng{Citizens \textbf{always} \textbf{respect} the law.} \quad \eng{She \textbf{is} \textbf{always} punctual.} \quad \eng{They \textbf{have} \textbf{never} \textbf{protested}.} \quad \eng{We \textbf{must} \textbf{always} \textbf{obey}.}
\item \textbf{Degré} (\eng{very, really, quite, extremely}) : Devant l'adjectif ou l'adverbe qu'ils modifient. \eng{very important}, \eng{really carefully}.
\item \textbf{Adverbes de manière en \eng{-ly}} : Souvent en fin de phrase (position M), ou avant le verbe pour l'emphase. \eng{He spoke \textbf{fluently}.} / \eng{He \textbf{fluently} \textbf{spoke} English.}
\end{itemize}
\end{propriete}

\begin{popfigure}
\begin{tikzpicture}[
    node distance=0.8cm and 1cm,
    slot/.style={rectangle, draw=poppurple, thick, fill=poppurpleL, rounded corners, minimum width=2.2cm, minimum height=1cm, align=center, font=\small},
    arrow/.style={-Stealth, poppurple, thick}
]
\node[slot, align=center] (S){Subject\\\eng{The crowd}};
\node[slot, right=of S, align=center] (V){Verb\\\eng{gathered}};
\node[slot, right=of V, align=center] (O){Object\\--};
\node[slot, right=of O, align=center] (M){Manner\\{\color{popgreen}\eng{peacefully}}};
\node[slot, right=of M, align=center] (P){Place\\{\color{popblue}\eng{in front of the ministry}}};
\node[slot, right=of P, align=center] (T){Time\\{\color{poporange}\eng{this morning}}};
\draw[arrow] (S) -- (V) -- (O) -- (M) -- (P) -- (T);
\node[below=0.5cm of O, font=\tiny, text=popmuted] {(intransitif)};
\end{tikzpicture}
\legende{Visualisation de l'ordre S V O M P T (Manière – Lieu – Temps)}
\end{popfigure}

\courssub{Maîtriser l'inversion sujet-verbe (Questions et Emphase)}

\begin{methode}[Construire l'inversion sans faute]
\begin{enumerate}
\item \textbf{Question fermée (Yes/No) :} Auxiliaire/Modal + Sujet + Verbe (BV).
\eng{\textbf{Did} the court \textbf{protect} his rights?} \quad \eng{\textbf{Has} the government \textbf{listened}?} \quad \eng{\textbf{Can} we \textbf{change} it?}
\item \textbf{Question ouverte (Wh-) :} Mot interrogatif + Auxiliaire/Modal + Sujet + Verbe (BV).
\eng{\textbf{Why did} the judge \textbf{release} him?} \quad \eng{\textbf{Where have} they \textbf{gone}?}
\item \textbf{Cas particulier : Le mot interrogatif EST le sujet} (\eng{Who, What, Which}) $\rightarrow$ \textbf{Pas d'inversion, pas de \eng{do/does/did}}. Le verbe s'accorde (3\textsuperscript{e} pers. sg. souvent).
\eng{\textbf{Who} \textbf{wrote} this article?} \quad \eng{\textbf{What} \textbf{happened} yesterday?}
\item \textbf{Question indirecte :} Après \eng{I wonder, Do you know, Could you tell me, I'd like to know} $\rightarrow$ \textbf{Ordre AFFIRMATIF} (Sujet + Verbe). \textbf{Pas de \eng{do/does/did}}.
\eng{Do you know \textbf{where the court is}?} (Jamais *\eng{where is the court})
\eng{I wonder \textbf{why they protested}.}
\item \textbf{Inversion emphatique (négative/restrictive en tête) :} \eng{Never, Rarely, Seldom, Hardly, Scarcely, Not only, Little, Only then/if/later} + Auxiliaire + Sujet + Verbe.
\eng{\textbf{Never have I seen} such courage.} \quad \eng{\textbf{Rarely does} the government \textbf{listen}.} \quad \eng{\textbf{Only then did} they \textbf{understand}.}
\end{enumerate}
\end{methode}

\begin{attention}[Les 5 erreurs fatales des francophones (Bac)]
\begin{enumerate}
\item \textbf{Oublier le recul du temps (Backshifting)} : \eng{He said he \textbf{will} come.} $\rightarrow$ \eng{He said he \textbf{would} come.} (Sanction systématique en transformation).
\item \textbf{Mettre \eng{will} après \eng{when / if / as soon as}} : \eng{When I \textbf{will be} eighteen...} $\rightarrow$ \eng{When I \textbf{am} eighteen, I \textbf{will} vote.} (Calque du français).
\item \textbf{Placer l'adjectif APRÈS le nom} : \eng{a country \textbf{democratic}} $\rightarrow$ \eng{a \textbf{democratic} country.} (Calque structurel).
\item \textbf{Garder l'inversion dans la question indirecte} : \eng{Do you know \textbf{where is} the court?} $\rightarrow$ \eng{Do you know \textbf{where the court is}?}
\item \textbf{Conjuguer le verbe après l'auxiliaire \eng{did/does/do}} : \eng{Did the police \textbf{arrested} him?} $\rightarrow$ \eng{Did the police \textbf{arrest} him?} \quad \eng{Does she \textbf{works}?} $\rightarrow$ \eng{Does she \textbf{work}?} (Double marque du temps).
\end{enumerate}
\end{attention}

\courssub{Entraînement guidé et production}

\begin{exoresolu}[Concordance des temps — Type Bac]
\textbf{Énoncé.} Rewrite the following sentences, starting with the words given. Make all the necessary changes (tense, pronouns, time/place markers).
\begin{enumerate}
\item “I am fighting for my rights,” the activist said. $\rightarrow$ The activist said that...
\item “We have never seen such injustice,” the villagers told the journalist. $\rightarrow$ The villagers told the journalist that...
\item “I will vote in the next election,” Mary declared. $\rightarrow$ Mary declared that...
\item “You must respect the law,” the judge told the accused. $\rightarrow$ The judge told the accused that...
\item “The police arrested two suspects yesterday,” the reporter announced. $\rightarrow$ The reporter announced that...
\end{enumerate}
\tcblower
\textbf{Solution détaillée.}
\begin{enumerate}
\item \eng{said} (past) $\rightarrow$ \eng{am fighting} (pres. cont.) recule à \eng{was fighting}. \eng{I} $\rightarrow$ \eng{she}. \eng{my} $\rightarrow$ \eng{her}.\\
\textbf{The activist said that she was fighting for her rights.}
\item \eng{told} (past) $\rightarrow$ \eng{have never seen} (pres. perf.) recule à \eng{had never seen}. \eng{We} $\rightarrow$ \eng{they}.\\
\textbf{The villagers told the journalist that they had never seen such injustice.}
\item \eng{declared} (past) $\rightarrow$ \eng{will} recule à \eng{would}. \eng{I} $\rightarrow$ \eng{she}. \eng{the next election} (marqueur adapté).\\
\textbf{Mary declared that she would vote in the following election.}
\item \eng{told} (past) $\rightarrow$ \eng{must} n'a pas de passé $\rightarrow$ \eng{had to}. \eng{You} $\rightarrow$ \eng{he/she} (l'accusé).\\
\textbf{The judge told the accused that he/she had to respect the law.}
\item \eng{announced} (past) $\rightarrow$ \eng{arrested} (past simple) recule à \eng{had arrested}. \eng{yesterday} $\rightarrow$ \eng{the day before / the previous day}.\\
\textbf{The reporter announced that the police had arrested two suspects the day before.}
\end{enumerate}
\end{exoresolu}

\begin{exoresolu}[Temps dans les subordonnées (when, if, after...) — Type Bac]
\textbf{Énoncé.} Complete each sentence with the correct form of the verb in brackets.
\begin{enumerate}
\item When the elections (take place), every citizen will be able to vote.
\item After the demonstrators (sign) the petition, they handed it to the mayor.
\item If the government (listen) to the people, there would be fewer protests.
\item The journalist left as soon as she (finish) her interview.
\item If the witness had told the truth, the innocent man (not / go) to prison.
\item Unless the authorities (act) now, the crisis will worsen.
\end{enumerate}
\tcblower
\textbf{Solution détaillée.}
\begin{enumerate}
\item \eng{when} + futur $\rightarrow$ Present Simple. \textbf{When the elections \textbf{take} place...}
\item Antériorité (\eng{after} + action terminée avant \eng{handed}) $\rightarrow$ Past Perfect. \textbf{After the demonstrators \textbf{had signed}...}
\item Type 2 (irréel présent) : \eng{if} + Past Simple $\rightarrow$ \eng{would}. \textbf{If the government \textbf{listened}...}
\item Antériorité (\eng{as soon as} + action terminée avant \eng{left}) $\rightarrow$ Past Perfect. \textbf{...as soon as she \textbf{had finished}...}
\item Type 3 (irréel passé) : \eng{if} + Past Perfect $\rightarrow$ \eng{would have} + V3. \textbf{...the innocent man \textbf{would not have gone} to prison.}
\item \eng{Unless} = \eng{If not}. Futur réel $\rightarrow$ Present Simple. \textbf{Unless the authorities \textbf{act} now...}
\end{enumerate}
\end{exoresolu}

\begin{exoresolu}[Ordre des mots et Questions — Type Bac]
\textbf{Énoncé.} Put the words in the correct order to make meaningful sentences.
\begin{enumerate}
\item rights / human / defends / organisation / this / actively
\item always / citizens / their / should / duties / remember
\item peacefully / the crowd / in front of the ministry / gathered / this morning
\item a / Cameroon / democratic / is / country
\item speaks / fluently / the mayor / English / at official ceremonies
\item never / have / I / seen / such / corruption / high-level
\end{enumerate}
\tcblower
\textbf{Solution détaillée.}
\begin{enumerate}
\item Sujet (\eng{this organisation}) + Adv. fréquence/manière (\eng{actively}) + Verbe (\eng{defends}) + COD (\eng{human rights}). \textbf{This organisation actively defends human rights.}
\item Sujet (\eng{citizens}) + Modal (\eng{should}) + Adv. fréquence (\eng{always}) + Verbe (\eng{remember}) + COD (\eng{their duties}). \textbf{Citizens should always remember their duties.}
\item Sujet (\eng{the crowd}) + Verbe (\eng{gathered}) + Manière (\eng{peacefully}) + Lieu (\eng{in front of the ministry}) + Temps (\eng{this morning}). \textbf{The crowd gathered peacefully in front of the ministry this morning.}
\item Sujet (\eng{Cameroon}) + Verbe (\eng{is}) + Det (\eng{a}) + Adj (\eng{democratic}) + Nom (\eng{country}). \textbf{Cameroon is a democratic country.}
\item Sujet (\eng{the mayor}) + Verbe (\eng{speaks}) + COD (\eng{English}) + Manière (\eng{fluently}) + Lieu (\eng{at official ceremonies}). \textbf{The mayor speaks English fluently at official ceremonies.}
\item Inversion emphatique (\eng{Never} + Aux \eng{have} + Sujet \eng{I} + V3 \eng{seen} + COD). \textbf{Never have I seen such high-level corruption.}
\end{enumerate}
\end{exoresolu}

\begin{exoresolu}[Questions et inversion — Type Bac]
\textbf{Énoncé.} Transform as indicated.
\begin{enumerate}
\item The court protected the witness. (Ask a question about \underline{the witness} — use \eng{who} as object.)
\item The lawyer defended the accused. (Write the corresponding question with \eng{who} as subject.)
\item “Where does the trial take place?” (Turn into an indirect question after \eng{Do you know}...)
\item The government has rarely listened to the opposition. (Begin the sentence with \eng{Rarely}...)
\item The police arrested the suspects last night. (Write a yes/no question.)
\item The activist said: “I can help you.” (Report the sentence.)
\end{enumerate}
\tcblower
\textbf{Solution détaillée.}
\begin{enumerate}
\item \eng{Who} = COD $\rightarrow$ Aux \eng{did} + Sujet \eng{the court} + BV \eng{protect}. \textbf{Who did the court protect?}
\item \eng{Who} = Sujet $\rightarrow$ Pas d'aux, verbe au passé. \textbf{Who defended the accused?}
\item Question indirecte $\rightarrow$ Ordre affirmatif, \eng{does} disparaît, \eng{take} $\rightarrow$ \eng{takes} (3\textsuperscript{e} pers. sg). \textbf{Do you know where the trial takes place?}
\item \eng{Rarely} en tête $\rightarrow$ Inversion auxiliaire \eng{has}. \textbf{Rarely has the government listened to the opposition.}
\item Yes/No question Past Simple $\rightarrow$ \eng{Did} + Sujet + BV. \textbf{Did the police arrest the suspects last night?}
\item \eng{said} (past) + \eng{can} $\rightarrow$ \eng{could}. \eng{I} $\rightarrow$ \eng{she}, \eng{you} $\rightarrow$ \eng{him/her/me}. \textbf{The activist said that she could help him/her/me.}
\end{enumerate}
\end{exoresolu}

\begin{exoresolu}[Production guidée : Discours rapporté — Type Bac]
\textbf{Énoncé.} A journalist interviewed a human rights defender in Bamenda. Using the notes below, write a short paragraph (5–6 sentences) reporting what she said. Begin: \emph{The activist told the journalist that...}\\
Notes (her exact words): “I work for a small organisation. We have helped many victims since 2015. The situation is difficult here. The authorities must listen to the people. I will never give up.”
\tcblower
\textbf{Solution détaillée — Raisonnement pas à pas.}
\begin{enumerate}
\item \textbf{Verbe introducteur :} \eng{told} (Past Simple) $\rightarrow$ Backshifting obligatoire.
\item \textbf{Phrase 1 :} \eng{I work} (Pres. Simple) $\rightarrow$ \eng{she worked}. \eng{for a small organisation} (inchangé).
\item \textbf{Phrase 2 :} \eng{We have helped} (Pres. Perf.) $\rightarrow$ \eng{they had helped}. \eng{since 2015} (inchangé, repère temporel absolu).
\item \textbf{Phrase 3 :} \eng{The situation is} (Pres. Simple) $\rightarrow$ \eng{the situation was}. \eng{difficult here} $\rightarrow$ \eng{difficult there}.
\item \textbf{Phrase 4 :} \eng{The authorities must listen} $\rightarrow$ \eng{the authorities had to listen}. (\eng{must} $\rightarrow$ \eng{had to}).
\item \textbf{Phrase 5 :} \eng{I will never give up} $\rightarrow$ \eng{she would never give up}. (\eng{will} $\rightarrow$ \eng{would}).
\item \textbf{Style :} Varier les verbes introducteurs (\eng{explained, added, stressed, promised}).
\end{enumerate}
\textbf{Réponse rédigée :}
\emph{The activist told the journalist that she worked for a small organisation. She explained that they had helped many victims since 2015. She added that the situation was difficult there. She stressed that the authorities had to listen to the people. Finally, she promised that she would never give up.}
\end{exoresolu}

\courssub{Exercices d'application (À faire seul)}

\begin{exercice}[Grammaire : Concordance, Ordre, Questions]
\textbf{Exercice 1 : Discours rapporté.} Rewrite in reported speech. The introductory verb is in the past.
\begin{enumerate}
\item “I am writing a report on corruption,” the journalist said.
\item “We have arrested the corrupt official,” the police chief announced.
\item “The funds will be recovered soon,” the minister promised.
\item “You must submit your declaration by Friday,” the clerk told me.
\item “I can’t attend the meeting tomorrow,” the delegate said.
\end{enumerate}

\textbf{Exercice 2 : Temps des subordonnées.} Put the verbs in brackets into the correct tense.
\begin{enumerate}
\item If the President \_\_\_\_\_\_\_\_ (sign) the bill, it will become law.
\item After the voters \_\_\_\_\_\_\_\_ (cast) their ballots, the counting began.
\item The opposition leader said he \_\_\_\_\_\_\_\_ (not / participate) in the dialogue unless his conditions \_\_\_\_\_\_\_\_ (be) met.
\item Hardly \_\_\_\_\_\_\_\_ (the judge / enter) the courtroom when the accused shouted.
\item By the time the UN observers \_\_\_\_\_\_\_\_ (arrive), the election \_\_\_\_\_\_\_\_ (already / end).
\end{enumerate}

\textbf{Exercice 3 : Ordre des mots.} Reorder the words to form correct sentences.
\begin{enumerate}
\item transparently / the funds / manages / organisation / this / very
\item rarely / the press / criticises / openly / government / the
\item a / fair / trial / every / citizen / deserves
\item yesterday / at the stadium / peacefully / the match / ended
\item not only / the law / protects / citizens / but / it / also / guides / them
\end{enumerate}

\textbf{Exercice 4 : Questions et inversion.} Transform as indicated.
\begin{enumerate}
\item The court sentenced the thief to five years. (Question on \eng{the thief} — object).
\item The judge sentenced the thief. (Question on \eng{the judge} — subject).
\item “When will the results be published?” (Indirect question after \eng{Could you tell me}...).
\item The government has never ignored a court ruling. (Begin with \eng{Never}...).
\item They have been investigating the case for months. (Yes/No question).
\end{enumerate}

\textbf{Exercice 5 : Production écrite (Mini-composition).}
A local NGO organised a workshop on “Youth and Civic Engagement” in Douala. You attended it. Write a short report (80–100 words) for your school magazine. Use reported speech to summarise what the speakers said. Include at least:
\begin{itemize}
\item One \eng{that}-clause with backshifting (past $\rightarrow$ past perfect).
\item One time clause with \eng{when} or \eng{after} + correct tense.
\item One conditional sentence (Type 1 or 2).
\item Correct word order (SVOMPT) and adjective placement.
\end{itemize}
\end{exercice}

\begin{corrige}[Exercices d'application]
\textbf{Exercice 1 :}
\begin{enumerate}
\item The journalist said that he/she was writing a report on corruption.
\item The police chief announced that they had arrested the corrupt official.
\item The minister promised that the funds would be recovered soon.
\item The clerk told me that I had to submit my declaration by Friday. (\eng{must} $\rightarrow$ \eng{had to})
\item The delegate said that he/she could not attend the meeting the next day / the following day. (\eng{can't} $\rightarrow$ \eng{couldn't}, \eng{tomorrow} $\rightarrow$ \eng{the next day})
\end{enumerate}

\textbf{Exercice 2 :}
\begin{enumerate}
\item \textbf{signs} (Type 1 : \eng{if} + Present Simple).
\item \textbf{had cast} (Antériorité : \eng{after} + Past Perfect).
\item \textbf{would not participate} ... \textbf{were} (Discours rapporté + Type 2 : \eng{said he would not}... \eng{unless} = \eng{if not}, condition irréelle présente \eng{were}).
\item \textbf{had the judge entered} (Inversion après \eng{Hardly} + Past Perfect).
\item \textbf{arrived} ... \textbf{had already ended} (Antériorité : \eng{By the time} + Past Simple / Past Perfect).
\end{enumerate}

\textbf{Exercice 3 :}
\begin{enumerate}
\item This organisation manages the funds very transparently. (S V O M)
\item Rarely does the government criticise the press openly. (Inversion emphatique \eng{Rarely} + Aux \eng{does} + S + V + Adv M). \textit{Note : sens inversé par rapport à l'énoncé désordonné, mais structure grammaticale imposée par \eng{Rarely}. Si phrase affirmative : The government rarely criticises the press openly.}
\item Every citizen deserves a fair trial. (Adj \eng{fair} avant \eng{trial}).
\item The match ended peacefully at the stadium yesterday. (S V M P T).
\item Not only does the law protect citizens, but it also guides them. (Inversion \eng{Not only} + Aux \eng{does} + S + V ; seconde clause ordre normal \eng{but it also guides them}).
\end{enumerate}

\textbf{Exercice 4 :}
\begin{enumerate}
\item Who did the court sentence to five years? (\eng{Who} = COD $\rightarrow$ \eng{did} + S + BV).
\item Who sentenced the thief? (\eng{Who} = Sujet $\rightarrow$ Pas d'aux, V au passé).
\item Could you tell me when the results will be published? (Ordre affirmatif, \eng{will} conservé car verbe introducteur \eng{Could you tell me} n'est pas au passé \textit{ici}, ou futur rapporté. Si principal au passé : \eng{...when the results would be published}. On accepte les deux selon contexte, mais \eng{will} frequent si la question est encore d'actualité).
\item Never has the government ignored a court ruling. (Inversion \eng{Never} + Aux \eng{has} + S + V3).
\item Have they been investigating the case for months? (Yes/No : Aux \eng{Have} + S + V-ing).
\end{enumerate}

\textbf{Exercice 5 : Proposition de corrigé type}
\emph{Last Saturday, an NGO organised a workshop on “Youth and Civic Engagement” in Douala. The coordinator told the participants that young people had a crucial role to play. She explained that when citizens \textbf{vote} (Type 1: present), they \textbf{will shape} their future. A lawyer added that the law \textbf{protected} everyone's rights and that the authorities \textbf{had to} respect court decisions. He stressed that if the youth \textbf{organised} themselves (Type 2: past), they \textbf{would be} heard. After the workshop \textbf{had ended}, many students \textbf{promised} they \textbf{would join} the NGO. It was a very \textbf{inspiring} event.}
\end{corrige}

\begin{savaistu}[Le Cameroun et la langue des droits]
Le Cameroun est un pays \textbf{bilingue} (anglais/français) par sa Constitution. Le vocabulaire des droits de l'homme (\eng{human rights}) y circule dans les deux langues. On entend souvent : \eng{rule of law} (État de droit), \eng{due process} (procédure régulière), \eng{freedom of speech} (liberté d'expression), \eng{accountability} (reddition de comptes / responsabilité). Maîtriser la grammaire du discours rapporté (\eng{reported speech}) est essentiel pour comprendre les communiqués officiels, les rapports d'ONG (comme \eng{CHRDA}, \eng{Amnesty International}) et les débats parlementaires retransmis à la CRTV ou sur \eng{The Cameroon Tribune}.
\end{savaistu}

\begin{aretenir}[Check-list finale pour le Bac — Lesson 44]
\begin{itemize}
\item \textbf{Backshifting :} Verbe principal au passé $\rightarrow$ Subordonnée recule (Tableau p.2). Exceptions : vérités générales, report immédiat.
\item \textbf{Marqueurs :} \eng{now$\rightarrow$then, today$\rightarrow$that day, here$\rightarrow$there, ago$\rightarrow$before, this$\rightarrow$that}.
\item \textbf{Subordonnées \eng{when/if} + Futur :} \textbf{Present Simple} (pas de \eng{will}).
\item \textbf{Antériorité passée :} \eng{After/When/Before} + \textbf{Past Perfect}.
\item \textbf{Conditionnelles :} Type 1 (\eng{if} + Pres $\rightarrow$ \eng{will}), Type 2 (\eng{if} + Past $\rightarrow$ \eng{would}), Type 3 (\eng{if} + Past Perf $\rightarrow$ \eng{would have}).
\item \textbf{Ordre SVOMPT :} Sujet - Verbe - Objet - Manière - Lieu - Temps. Adjectif \textbf{avant} nom.
\item \textbf{Adverbes fréquence :} Avant V principal, Après Aux/Be/Modaux.
\item \textbf{Questions :} \eng{Who} sujet = pas d'inversion. Question indirecte = ordre affirmatif (pas de \eng{do/does/did}).
\item \textbf{Inversion emphatique :} \eng{Never/Rarely/Seldom/Hardly} + Aux + Sujet + V.
\item \textbf{5 fautes interdites :} Oubli backshifting, \eng{will} après \eng{when/if}, Adj après Nom, Inversion question indirecte, Double passé (\eng{did} + V-ed).
\end{itemize}
\end{aretenir}

\newpage

\coursec{Exercices}

\courssub{Je vérifie mes connaissances}

\begin{exercice}[QCM : la bonne concordance]
Pour chaque phrase, choisis la bonne forme du verbe entre parenthèses.
\begin{enumerate}
\item The teacher said that the earth \textbf{(is / was / were)} round.
\item He told me that he \textbf{(will / would / shall)} protect my rights.
\item She asked me if I \textbf{(am / was / had been)} registered on the electoral list.
\item The journalist reported that the minister \textbf{(has announced / had announced / announces)} new measures the day before.
\item They said they \textbf{(can / could / may)} not attend the meeting.
\item \eng{Human rights are universal,} the speaker reminded us. He added that they \textbf{(are / were)} still violated in many countries.
\end{enumerate}
\end{exercice}

\begin{exercice}[Vrai ou faux ? Justifie ta réponse]
Dis si chaque phrase est \eng{correct} ou \eng{incorrect}, puis justifie en citant la règle de grammaire concernée.
\begin{enumerate}
\item \eng{He told that he will come tomorrow.}
\item \eng{She asked me where did I live.}
\item \eng{The lawyer said that justice must be respected.}
\item \eng{My father always speaks very well English.}
\item \eng{They announced that they had voted the previous day.}
\item \eng{Citizens must respect always their duties.}
\end{enumerate}
\end{exercice}

\begin{exercice}[Vocabulaire : les verbes introducteurs]
Complète chaque phrase avec le verbe introducteur qui convient, choisi dans la liste : \eng{asked, told, said, explained, promised, warned}.
\begin{enumerate}
\item The mayor \dots\ the crowd that he would build a new market.
\item The policeman \dots\ the driver not to park in front of the school.
\item \eng{Where do you vote?} She \dots\ him where he voted.
\item The candidate \dots\ that he would reduce taxes if he was elected.
\item The teacher \dots\ us that the exam would take place in June.
\item The human rights activist \dots\ why every citizen should register to vote.
\end{enumerate}
\end{exercice}

\begin{exercice}[Conjugaison : le backshift]
Mets le verbe entre parenthèses à la forme correcte dans le discours rapporté.
\begin{enumerate}
\item \eng{I am tired,} she said. $\rightarrow$ She said that she \dots\ (be) tired.
\item \eng{We have finished the work,} they told him. $\rightarrow$ They told him that they \dots\ (finish) the work.
\item \eng{I will vote tomorrow,} he said. $\rightarrow$ He said that he \dots\ (vote) the next day.
\item \eng{I can help you,} she told me. $\rightarrow$ She told me that she \dots\ (can) help me.
\item \eng{We were watching the debate,} they said. $\rightarrow$ They said that they \dots\ (watch) the debate.
\item \eng{I bought a house in Buea last year,} my uncle said. $\rightarrow$ My uncle said that he \dots\ (buy) a house in Buea the year before.
\end{enumerate}
\end{exercice}

\courssub{Je m'entraîne}

\begin{exercice}[Traduction vers l'anglais]
Traduis les phrases suivantes en anglais, en respectant la concordance des temps et l'ordre des mots.
\begin{enumerate}
\item L'avocat a dit à la foule qu'il protégerait leurs droits.
\item Elle m'a demandé si j'étais inscrit sur la liste électorale.
\item Le professeur a expliqué que la concordance des temps était importante.
\item Ils ont annoncé qu'ils avaient voté la veille à Douala.
\item Ma mère parle très bien l'anglais.
\item Il m'a demandé où j'habitais et ce que je faisais dans la vie.
\end{enumerate}
\end{exercice}

\begin{exercice}[Transformation : du discours direct au discours indirect]
Rapporte les phrases suivantes au discours indirect, en commençant par les verbes introducteurs indiqués.
\begin{enumerate}
\item \eng{I will protect your rights,} the lawyer told the crowd. $\rightarrow$ The lawyer told the crowd that \dots
\item \eng{Are you registered on the electoral list?} $\rightarrow$ He asked me \dots
\item \eng{Where do you vote?} $\rightarrow$ She asked him \dots
\item \eng{We have never seen such a large crowd,} the journalists said. $\rightarrow$ The journalists said that \dots
\item \eng{Don't cross the road here,} the policeman told the children. $\rightarrow$ The policeman told the children \dots
\item \eng{I am preparing my speech for tomorrow,} the candidate said. $\rightarrow$ The candidate said that \dots
\end{enumerate}
\end{exercice}

\begin{exercice}[Texte à trous : la concordance en contexte]
Complète le texte suivant avec la forme correcte des verbes entre parenthèses.

\begin{textebox}[At the town hall]
Last Saturday, the mayor of Bamenda announced that the town hall \dots\ (organise) a public meeting the following week. He explained that every citizen \dots\ (can) attend and ask questions. He added that the council \dots\ (already / discuss) the problem of water supply and that a solution \dots\ (be) found soon. A journalist asked him whether the elections \dots\ (take place) in December as planned. The mayor replied that he \dots\ (not / know) yet, but that he \dots\ (inform) the population as soon as possible.
\end{textebox}
\end{exercice}

\begin{exercice}[Remets les mots dans l'ordre]
Remets les mots suivants dans le bon ordre pour former des phrases anglaises correctes. Attention à la place des adverbes et des compléments.
\begin{enumerate}
\item always / citizens / their / respect / must / duties
\item speaks / very well / my sister / English
\item me / where / asked / he / lived / I
\item yesterday / to the market / went / early / my mother
\item told / the truth / never / the politician / us
\item in Yaoundé / for five years / have lived / we
\end{enumerate}
\end{exercice}

\courssub{Je me prépare au Bac}

\begin{exercice}[Texte et questions de compréhension — type Bac]
Lis le texte suivant, puis réponds aux questions en anglais, par des phrases complètes.

\begin{textebox}[A speech for the future]
On Saturday morning, more than two thousand people gathered in front of the community hall in Douala to listen to Mrs Etame, a candidate in the municipal elections. She began by thanking the crowd and said that she had always dreamed of serving her city. She explained that when she was a child, her grandmother had told her that a good citizen never stayed silent in front of injustice. She promised that if she was elected, she would improve the roads, build a new library and create jobs for young people. She added that the town council had wasted too much money in the past and that things would change under her leadership. A young man in the crowd asked her whether she would reduce market taxes. She answered that she would study the question carefully and that she would consult the traders before taking any decision. At the end of her speech, she declared that the future of Douala belonged to its youth, and she invited everyone to register on the electoral list before the end of the month.
\end{textebox}

\textbf{Comprehension questions}
\begin{enumerate}
\item Where did the crowd gather, and why? (2 marks)
\item What did Mrs Etame's grandmother tell her when she was a child? (2 marks)
\item Mention three promises made by the candidate. (3 marks)
\item What question did the young man ask, and how did Mrs Etame answer? (2 marks)
\item What did the candidate invite the people to do at the end of her speech? (1 mark)
\end{enumerate}

\textbf{Language work}
\begin{enumerate}
\item Find in the text three verbs in the past perfect and write them down. (3 marks)
\item Report this sentence from the text: \eng{I will improve the roads,} she said. (2 marks)
\item Report this question: \eng{Will you reduce market taxes?} the young man asked her. (2 marks)
\item Correct the mistakes in this sentence: \eng{She told to the crowd that she will build a library.} (2 marks)
\end{enumerate}
\end{exercice}

\begin{exercice}[Traduction — type Bac]
Traduis en anglais le passage suivant. Soigne la concordance des temps, l'ordre des mots et le choix des verbes introducteurs. (10 marks)

\begin{textebox}[Texte à traduire]
Le candidat a déclaré qu'il avait toujours défendu les droits des citoyens. Il a promis qu'il construirait une nouvelle école et qu'il aiderait les jeunes à trouver du travail. Une femme dans la foule lui a demandé s'il réduirait les impôts. Il a répondu qu'il étudierait la question avec soin et qu'il prendrait une décision le mois suivant. À la fin de son discours, il a dit que l'avenir du pays appartenait à sa jeunesse.
\end{textebox}
\end{exercice}

\begin{exercice}[Composition — type Bac]
\textbf{Topic:} Imagine you are a journalist for your school magazine. You attended a speech given by a candidate in the municipal elections of Douala. Write a report of about \textbf{150 to 180 words} in which you report what the candidate said, using reported speech.

\textbf{Your report must:}
\begin{enumerate}
\item begin with a short introduction saying who the candidate is, where and when the speech took place;
\item report at least \textbf{four statements} made by the candidate, with at least \textbf{three different tenses} in the backshift (past simple, past perfect, conditional);
\item report at least \textbf{one question} asked by a member of the audience and the candidate's answer;
\item end with your personal opinion about the speech.
\end{enumerate}

\textbf{Useful language:} \eng{the candidate declared that... / she promised that... / he explained that... / a voter asked whether... / she replied that... / in my opinion...}
\end{exercice}

\newpage

\coursec{Corrigés des exercices}

\begin{corrige}[Exercice 1 — QCM : la bonne concordance]
\begin{enumerate}
\item \eng{The teacher said that the earth \textbf{is} round.} — Vérité générale et scientifique : le présent est conservé malgré le verbe introducteur au passé. (On accepte aussi \eng{was} dans un style très formel, mais \eng{is} est la réponse attendue.)
\item \eng{He told me that he \textbf{would} protect my rights.} — Backshift : \eng{will} $\rightarrow$ \eng{would} après un verbe au passé.
\item \eng{She asked me if I \textbf{was} registered on the electoral list.} — Question rapportée : présent $\rightarrow$ prétérit (\eng{am} $\rightarrow$ \eng{was}).
\item \eng{The journalist reported that the minister \textbf{had announced} new measures the day before.} — L'action est antérieure au moment du discours : \eng{has announced} $\rightarrow$ \eng{had announced} (past perfect).
\item \eng{They said they \textbf{could} not attend the meeting.} — Backshift du modal : \eng{can} $\rightarrow$ \eng{could}.
\item \eng{He added that they \textbf{were} still violated in many countries.} — Ici, on rapporte une affirmation du locuteur (et non une vérité universelle au sens strict) : prétérit \eng{were}. (On accepte \eng{are} si l'on insiste sur la validité présente du constat.)
\end{enumerate}
\end{corrige}

\begin{corrige}[Exercice 2 — Vrai ou faux ?]
\begin{enumerate}
\item \textbf{Incorrect.} Deux erreurs : (a) \eng{tell} exige un complément d'objet (\eng{told \textbf{me/us} that...}) ; (b) la concordance impose \eng{will} $\rightarrow$ \eng{would}. Correction : \eng{He told me that he would come the next day.}
\item \textbf{Incorrect.} Dans une question rapportée, on revient à l'ordre sujet–verbe, sans auxiliaire \eng{did}. Correction : \eng{She asked me where I lived.}
\item \textbf{Correct.} \eng{Must} n'a pas de forme passée : il reste invariable au discours rapporté (\eng{He said that justice must be respected}).
\item \textbf{Incorrect.} En anglais, le complément d'objet ne se sépare pas du verbe par un adverbe de manière ; l'ordre est verbe + objet + adverbe. Correction : \eng{My father always speaks English very well.}
\item \textbf{Correct.} Past perfect \eng{had voted} pour une action antérieure au discours, et \eng{the previous day} remplace correctement \eng{yesterday}.
\item \textbf{Incorrect.} L'adverbe de fréquence se place après le modal et avant le verbe lexical : \eng{Citizens must \textbf{always} respect their duties.}
\end{enumerate}
\end{corrige}

\begin{corrige}[Exercice 3 — Les verbes introducteurs]
\begin{enumerate}
\item \eng{The mayor \textbf{promised} the crowd that he would build a new market.} (promesse : \eng{would} + futur rapporté)
\item \eng{The policeman \textbf{warned} the driver not to park in front of the school.} (avertissement : \eng{warn somebody not to do})
\item \eng{She \textbf{asked} him where he voted.} (question rapportée $\rightarrow$ \eng{ask})
\item \eng{The candidate \textbf{said} that he would reduce taxes if he was elected.} (déclaration simple, sans complément de personne $\rightarrow$ \eng{said that})
\item \eng{The teacher \textbf{told} us that the exam would take place in June.} (\eng{tell} + complément d'objet \eng{us})
\item \eng{The human rights activist \textbf{explained} why every citizen should register to vote.} (\eng{explain} + \eng{why} ; jamais \eng{explain me})
\end{enumerate}
\end{corrige}

\begin{corrige}[Exercice 4 — Le backshift]
\begin{enumerate}
\item \eng{She said that she \textbf{was} tired.} (présent $\rightarrow$ prétérit)
\item \eng{They told him that they \textbf{had finished} the work.} (present perfect $\rightarrow$ past perfect)
\item \eng{He said that he \textbf{would vote} the next day.} (\eng{will} $\rightarrow$ \eng{would} ; \eng{tomorrow} $\rightarrow$ \eng{the next day})
\item \eng{She told me that she \textbf{could} help me.} (\eng{can} $\rightarrow$ \eng{could})
\item \eng{They said that they \textbf{had been watching} the debate.} (past continuous $\rightarrow$ past perfect continuous ; on accepte aussi \eng{were watching} si l'action est encore en cours)
\item \eng{My uncle said that he \textbf{had bought} a house in Buea the year before.} (past simple $\rightarrow$ past perfect ; \eng{last year} $\rightarrow$ \eng{the year before})
\end{enumerate}
\end{corrige}

\begin{corrige}[Exercice 5 — Traduction vers l'anglais]
\begin{enumerate}
\item \eng{The lawyer told the crowd that he would protect their rights.}
\item \eng{She asked me if I was registered on the electoral list.}
\item \eng{The teacher explained that the sequence of tenses was important.}
\item \eng{They announced that they had voted the day before in Douala.}
\item \eng{My mother speaks English very well.} (ordre : verbe + objet + adverbe de manière)
\item \eng{He asked me where I lived and what I did for a living.} (questions rapportées : ordre sujet–verbe, pas d'inversion)
\end{enumerate}
\textbf{Points de vigilance :} ne pas écrire \eng{explained me} ; ne pas traduire « la veille » par \eng{yesterday} au discours rapporté (\eng{the day before}) ; « ce que je faisais dans la vie » se rend idiomatiquement par \eng{what I did for a living}.
\end{corrige}

\begin{corrige}[Exercice 6 — Du discours direct au discours indirect]
\begin{enumerate}
\item \eng{The lawyer told the crowd that he would protect their rights.}
\item \eng{He asked me if I was registered on the electoral list.} (question fermée $\rightarrow$ \eng{if/whether}, sans inversion)
\item \eng{She asked him where he voted.} (question en \eng{wh-} : on garde le mot interrogatif, ordre sujet–verbe)
\item \eng{The journalists said that they had never seen such a large crowd.} (present perfect $\rightarrow$ past perfect)
\item \eng{The policeman told the children not to cross the road there.} (impératif négatif $\rightarrow$ \eng{not to} + base verbale ; \eng{here} $\rightarrow$ \eng{there})
\item \eng{The candidate said that he was preparing his speech for the next day.} (present continuous $\rightarrow$ past continuous ; \eng{tomorrow} $\rightarrow$ \eng{the next day})
\end{enumerate}
\end{corrige}

\begin{corrige}[Exercice 7 — Texte à trous : la concordance en contexte]
\eng{Last Saturday, the mayor of Bamenda announced that the town hall \textbf{would organise} a public meeting the following week. He explained that every citizen \textbf{could} attend and ask questions. He added that the council \textbf{had already discussed} the problem of water supply and that a solution \textbf{would be} found soon. A journalist asked him whether the elections \textbf{would take place} in December as planned. The mayor replied that he \textbf{did not know} yet, but that he \textbf{would inform} the population as soon as possible.}

\textbf{Justifications :} \eng{will organise} $\rightarrow$ \eng{would organise} ; \eng{can} $\rightarrow$ \eng{could} ; action antérieure au discours $\rightarrow$ past perfect \eng{had already discussed} ; \eng{will be found} $\rightarrow$ \eng{would be found} (conditionnel passif) ; \eng{will take place} $\rightarrow$ \eng{would take place} ; \eng{I don't know} $\rightarrow$ \eng{did not know} ; \eng{I will inform} $\rightarrow$ \eng{would inform}.
\end{corrige}

\begin{corrige}[Exercice 8 — Remets les mots dans l'ordre]
\begin{enumerate}
\item \eng{Citizens must always respect their duties.} (adverbe de fréquence : après le modal, avant le verbe lexical)
\item \eng{My sister speaks English very well.} (verbe + objet + adverbe de manière)
\item \eng{He asked me where I lived.} (question rapportée : pas d'inversion)
\item \eng{My mother went to the market early yesterday.} (ordre : sujet + verbe + lieu + manière + temps)
\item \eng{The politician never told us the truth.} (\eng{never} avant le verbe lexical)
\item \eng{We have lived in Yaoundé for five years.} (lieu avant la durée)
\end{enumerate}
\textbf{Rappel :} l'ordre anglais canonique est \textbf{Sujet – Verbe – Objet – Manière – Lieu – Temps}, contrairement au français qui place souvent le temps en tête.
\end{corrige}

\begin{corrige}[Exercice 9 — Texte et questions de compréhension — type Bac]
\textbf{Comprehension questions — réponses attendues}
\begin{enumerate}
\item \eng{The crowd gathered in front of the community hall in Douala to listen to Mrs Etame, a candidate in the municipal elections.} (2 marks : lieu 1, raison 1)
\item \eng{Her grandmother had told her that a good citizen never stayed silent in front of injustice.} (2 marks)
\item \eng{She promised that she would improve the roads, build a new library and create jobs for young people.} (3 marks : 1 par promesse)
\item \eng{The young man asked her whether she would reduce market taxes. She answered that she would study the question carefully and consult the traders before taking any decision.} (2 marks)
\item \eng{She invited everyone to register on the electoral list before the end of the month.} (1 mark)
\end{enumerate}

\textbf{Language work — réponses attendues}
\begin{enumerate}
\item Trois verbes au past perfect : \eng{had always dreamed}, \eng{had told}, \eng{had wasted}. (3 marks : 1 par verbe)
\item \eng{She said that she would improve the roads.} (2 marks : \eng{will} $\rightarrow$ \eng{would})
\item \eng{The young man asked her whether she would reduce market taxes.} (2 marks : question fermée $\rightarrow$ \eng{whether/if}, ordre sujet–verbe)
\item \eng{She told the crowd that she would build a library.} (2 marks : suppression de \eng{to} après \eng{told} ; \eng{will} $\rightarrow$ \eng{would})
\end{enumerate}

\textbf{Points de vigilance :} répondre par des phrases complètes ; reprendre la concordance des temps dans les réponses elles-mêmes ; \eng{tell} ne se construit jamais avec \eng{to} + personne (contrairement à \eng{say to somebody}).
\end{corrige}

\begin{corrige}[Exercice 10 — Traduction — type Bac]
\textbf{Proposition de traduction}

\eng{The candidate declared that he had always defended the citizens' rights. He promised that he would build a new school and that he would help young people find work. A woman in the crowd asked him whether he would reduce taxes. He replied that he would study the question carefully and that he would take a decision the following month. At the end of his speech, he said that the future of the country belonged to its youth.}

\textbf{Barème indicatif (10 marks) :} 2 points par phrase (1 pour la concordance des temps, 1 pour le vocabulaire et l'ordre des mots).

\textbf{Points de vigilance :}
\begin{itemize}
\item « avait toujours défendu » $\rightarrow$ past perfect \eng{had always defended} (antériorité).
\item « construirait / aiderait / réduirait / étudierait / prendrait » $\rightarrow$ conditionnel \eng{would} + base verbale (futur dans le passé).
\item « s'il réduirait » : question fermée rapportée $\rightarrow$ \eng{whether/if}, sans inversion.
\item « le mois suivant » $\rightarrow$ \eng{the following month} (jamais \eng{the next month} calqué sur « le mois prochain » dans ce contexte).
\item « appartenait à sa jeunesse » $\rightarrow$ \eng{belonged to its youth} ; attention au possessif \eng{its} sans apostrophe.
\end{itemize}
\end{corrige}

\begin{corrige}[Exercice 11 — Composition — type Bac]
\textbf{Proposition de rédaction modèle (environ 165 mots)}

\eng{Last Saturday, I attended a speech given by Mrs Etame, a candidate in the municipal elections of Douala, in front of the community hall. More than two thousand people were present.}

\eng{The candidate declared that she had always dreamed of serving her city. She explained that her grandmother had told her that a good citizen never stayed silent in front of injustice. She promised that she would improve the roads and that she would create jobs for young people. She also said that the town council had wasted too much money in the past.}

\eng{A young voter asked her whether she would reduce market taxes. She replied that she would study the question carefully before taking any decision.}

\eng{In my opinion, it was an inspiring speech, and many people left the hall with hope.}

\textbf{Barème indicatif (sur 10 ou 20 selon le barème de l'établissement) :}
\begin{itemize}
\item Respect de la tâche et de la longueur (150–180 mots) : 2 points.
\item Introduction (qui, où, quand) : 1 point.
\item Au moins quatre déclarations rapportées avec trois temps différents (past simple, past perfect, conditionnel) : 4 points.
\item Une question rapportée + la réponse : 2 points.
\item Opinion personnelle finale, correction linguistique générale : 1 point.
\end{itemize}

\textbf{Points de vigilance :}
\begin{itemize}
\item Varier les verbes introducteurs (\eng{declared, explained, promised, added, replied}) pour éviter la répétition de \eng{said}.
\item Vérifier chaque backshift : \eng{will} $\rightarrow$ \eng{would}, \eng{have/has} $\rightarrow$ \eng{had}, \eng{can} $\rightarrow$ \eng{could}.
\item Question rapportée : \eng{asked whether she would...}, jamais \eng{asked would she...}.
\item Ne pas mélanger discours direct et indirect : pas de guillemets dans un rapport au discours indirect.
\end{itemize}
\end{corrige}

\newpage

\coursec{Fiche bilan}

\begin{popfigure}
\begin{tikzpicture}[
  mindmap,
  concept color=popblue,
  level 1/.append style={level distance=3.5cm, sibling angle=360/6},
  level 2/.append style={level distance=2.5cm},
  every node/.append style={font=\small, align=center},
  branch1/.style={concept color=poporange, grow=30},
  branch2/.style={concept color=popgreen, grow=90},
  branch3/.style={concept color=poppurple, grow=150},
  branch4/.style={concept color=poppink, grow=210},
  branch5/.style={concept color=popgold, grow=270},
  branch6/.style={concept color=popblue, grow=330}
]
\node[concept, text=white, font=\bfseries\large, align=center]{Sequence of Tenses\\ \& Word Order}
  child[branch1] { node[concept, text=white, align=center]{Reported Speech\\ (Backshifting)} }
  child[branch2] { node[concept, text=white, align=center]{Exceptions\\ (No Backshift)} }
  child[branch3] { node[concept, text=white, align=center]{Word Order\\ SVOMPT} }
  child[branch4] { node[concept, text=white, align=center]{Questions \&\\ Negation Order} }
  child[branch5] { node[concept, text=white, align=center]{Adverb\\ Placement} }
  child[branch6] { node[concept, text=white, align=center]{French vs English\\ Traps} };
\end{tikzpicture}
\legende{Carte mentale : Concordance des temps et ordre des mots (Lesson 44)}
\end{popfigure}

\begin{bilan}

\end{bilan}

\courssub{Lexique clé — Key Vocabulary}

\begin{center}
\begin{tabularx}{\linewidth}{|X|X|X|}
\hline
\rowcolor{popblueL} \textbf{English} & \textbf{Français} & \textbf{Contexte / Exemple} \\
\hline
\cle{Reporting verb} & Verbe introducteur & \eng{He \textbf{stated} that...}, \eng{She \textbf{wondered} if...} \\
\cle{Backshifting} & Recul du temps & \eng{Present $\to$ Past, Past $\to$ Past Perfect} \\
\cle{Time marker} & Marqueur temporel & \eng{now $\to$ then, today $\to$ that day, tomorrow $\to$ the next day} \\
\cle{Place marker} & Marqueur spatial & \eng{here $\to$ there, this $\to$ that, these $\to$ those} \\
\cle{SVOMPT} & Ordre canonique & \textbf{S}ubject \textbf{V}erb \textbf{O}bject \textbf{M}anner \textbf{P}lace \textbf{T}ime \\
\cle{Wh-word} & Mot interrogatif & \eng{who, what, where, when, why, how} \\
\cle{Auxiliary} & Auxiliaire & \eng{do, be, have, modals (can, must, will)} \\
\hline
\end{tabularx}
\end{center}

\courssub{Règles de grammaire à connaître par cœur}

\textbf{1. Concordance des temps (Sequence of Tenses) — Discours rapporté}

\begin{center}
\begin{tabularx}{\linewidth}{|X|X|X|}
\hline
\rowcolor{popblueL} \textbf{Direct Speech (Temps)} & \textbf{Reported Speech (Backshift)} & \textbf{Exemple} \\
\hline
Present Simple & Past Simple & \eng{``I \textbf{live} in Buea.''} $\to$ He said he \textbf{lived} in Buea. \\
Present Continuous & Past Continuous & \eng{``I \textbf{am reading}.''} $\to$ She said she \textbf{was reading}. \\
Present Perfect & Past Perfect & \eng{``I \textbf{have finished}.''} $\to$ He said he \textbf{had finished}. \\
Past Simple & Past Perfect & \eng{``I \textbf{saw} him.''} $\to$ She said she \textbf{had seen} him. \\
Past Continuous & Past Perfect Continuous & \eng{``I \textbf{was working}.''} $\to$ He said he \textbf{had been working}. \\
\cle{will} & \cle{would} & \eng{``I \textbf{will} call.''} $\to$ She said she \textbf{would} call. \\
\cle{can} & \cle{could} & \eng{``I \textbf{can} speak English.''} $\to$ He said he \textbf{could} speak English. \\
\cle{must} / \cle{have to} & \cle{had to} & \eng{``I \textbf{must} go.''} $\to$ He said he \textbf{had to} go. \\
\hline
\end{tabularx}
\end{center}

\textbf{2. Ordre des mots (Word Order) — Structure SVOMPT}

\begin{center}
\begin{tabularx}{\linewidth}{|c|X|X|}
\hline
\rowcolor{popblueL} \textbf{Position} & \textbf{Élément} & \textbf{Exemple} \\
\hline
1 & \textbf{S}ubject & \eng{The students} \\
2 & \textbf{V}erb & \eng{study} \\
3 & \textbf{O}bject & \eng{human rights} \\
4 & \textbf{M}anner & \eng{carefully} \\
5 & \textbf{P}lace & \eng{in the library} \\
6 & \textbf{T}ime & \eng{every afternoon} \\
\hline
\end{tabularx}
\end{center}

\courssub{Méthodes express}

\begin{methode}[Transformer en discours rapporté]
\begin{enumerate}
\item Identifier le \cle{reporting verb} (temps ? $\to$ détermine le backshifting).
\item Changer les \cle{pronoms} (I $\to$ he/she, my $\to$ his/her, we $\to$ they).
\item Appliquer le \cle{backshifting} aux verbes (sauf exceptions).
\item Adapter les \cle{marqueurs spatio-temporels} (now $\to$ then, here $\to$ there, this $\to$ that).
\item Vérifier l'\cle{ordre des mots} : phrase affirmative (SVOMPT), pas ordre interrogatif.
\end{enumerate}
\end{methode}

\begin{methode}[Vérifier l'ordre des mots (SVOMPT)]
\begin{enumerate}
\item Trouver le \textbf{Sujet} (qui/quoi ?).
\item Trouver le \textbf{Verbe} (action/état).
\item Placer le \textbf{Complément d'Objet} (COI puis COD si pronoms, sinon COD puis COI).
\item Placer les \textbf{Circonstances} : \textbf{Manière} $\to$ \textbf{Lieu} $\to$ \textbf{Temps}.
\item \textbf{Attention} : Adverbes de fréquence (always, often) $\to$ \emph{avant} verbe principal, \emph{après} auxiliaire/be.
\end{enumerate}
\end{methode}



\begin{aretenir}[Checklist avant le Bac]
\begin{itemize}
\item $\square$ Je sais appliquer le \textbf{backshifting} (Present $\to$ Past, Will $\to$ Would, Can $\to$ Could, etc.).
\item $\square$ Je modifie correctement les \textbf{pronoms} (personnels, possessifs, démonstratifs).
\item $\square$ J'adapte les \textbf{marqueurs de temps} (now/then, today/that day, tomorrow/the next day...).
\item $\square$ J'adapte les \textbf{marqueurs de lieu} (here/there, this/that).
\item $\square$ Je connais les \textbf{3 exceptions au backshifting} : vérité générale, verbe introducteur au présent, situation encore actuelle.
\item $\square$ Je respecte l'ordre \textbf{SVOMPT} dans la phrase affirmative (Sujet + Verbe + Objet + Manière + Lieu + Temps).
\item $\square$ Je place les \textbf{adverbes de fréquence} correctement (avant V principal, après be/auxiliaire).
\item $\square$ Je construis les \textbf{questions indirectes} sans inversion (ordre affirmatif : \eng{He asked where I \textbf{lived}}).
\item $\square$ Je n'utilise pas \eng{do/does/did} dans les questions indirectes (sauf question tag ou emphase).
\item $\square$ Je distingue \textbf{ordre interrogatif} (Aux + S + V) et \textbf{ordre affirmatif} (S + V) dans le discours rapporté.
\item $\square$ Je place correctement les \textbf{adverbes de manière} (souvent à la fin ou avant le verbe pour l'emphase).
\item $\square$ Je traduis ``Il a dit que'' par \eng{He said that} (pas \eng{He told that} $\to$ \eng{He told \textbf{me} that}).
\end{itemize}
\end{aretenir}

\findecours

\end{document}
